\documentclass[11pt,a4paper]{article}
\usepackage[margin=26mm,headheight=14pt]{geometry}
\usepackage{fontspec}
\setmainfont{Latin Modern Roman}
\usepackage{amsmath,amssymb,amsthm,mathtools}
\usepackage{microtype,booktabs,longtable,array,enumitem}
\usepackage[round,authoryear]{natbib}
\usepackage[colorlinks=true,linkcolor=blue!45!black,citecolor=blue!45!black,urlcolor=blue!45!black]{hyperref}
\usepackage{xcolor}
\usepackage{tikz-cd}
\usepackage{fancyhdr}
\pagestyle{fancy}\fancyhf{}\fancyhead[L]{\small TOTASPH}\fancyhead[R]{\small Total asphericity}\fancyfoot[C]{\thepage}
\setlength{\emergencystretch}{3em}
\renewcommand{\thesection}{\Alph{section}}
\newtheorem{theorem}{Theorem}[section]
\newtheorem{proposition}[theorem]{Proposition}
\newtheorem{lemma}[theorem]{Lemma}
\newtheorem{corollary}[theorem]{Corollary}
\theoremstyle{definition}
\newtheorem{definition}[theorem]{Definition}
\newtheorem{example}[theorem]{Example}
\newtheorem{remark}[theorem]{Remark}
\newtheorem{question}[theorem]{Question}
\newcommand{\Cat}{\mathsf{Cat}}
\newcommand{\Set}{\mathsf{Set}}
\newcommand{\Ab}{\mathsf{Ab}}
\newcommand{\Hot}{\mathsf{Hot}}
\newcommand{\op}{\mathrm{op}}
\newcommand{\id}{\mathrm{id}}
\newcommand{\Winfty}{W_{\infty}}
\newcommand{\Winf}{W_{\infty}}
\newcommand{\ZZ}{\mathbb Z}
\newcommand{\NN}{\mathbb N}
\DeclareMathOperator{\Hom}{Hom}
\DeclareMathOperator{\Ob}{Ob}
\DeclareMathOperator{\El}{El}
\DeclareMathOperator{\colim}{colim}
\DeclareMathOperator{\limop}{lim}
\DeclareMathOperator{\Lan}{Lan}
\DeclareMathOperator{\Tor}{Tor}
\DeclareMathOperator{\Ind}{Ind}
\DeclareMathOperator{\Kar}{Kar}
\DeclareMathOperator{\Aff}{Aff}
\DeclareMathOperator{\Ch}{Ch}
\title{TOTASPH: Total Asphericity, Final Images,\newline and Dwyer-Type Immersions}
\author{AI-produced by GPT-6 Astra\\\small OpenAI, ChatGPT Pro mode}
\date{10 November 2026}
\begin{document}
\maketitle
\noindent We use a fixed Grothendieck universe $\mathcal U$. Categories, sets, and diagrams are $\mathcal U$-small unless an enlargement is explicitly stated. We put $BA=|NA|$, and $\Winfty$ is the class of functors inducing weak homotopy equivalences on nerves. We write $\Hot=\Cat[\Winfty^{-1}]$. For $u:A\to B$, $A/b$ consists of arrows $u(a)\to b$, whereas $b\backslash A$ consists of arrows $b\to u(a)$. A functor is \emph{final} when each $b\backslash A$ is nonempty and connected. Presheaves are contravariant. Multiplication in an endomorphism monoid is composition, so $xy=x\circ y$. A sieve is closed under sources, and a cosieve is closed under targets. The empty category is not filtered. A fundamental localizer satisfies weak saturation, the terminal-object axiom, and the relative Theorem A axiom, in exactly the form LF1--LF3 of Definition~\ref{A:localizer}; no extra filtered-colimit axiom is assumed. The manuscript's use of ``total asphericity'' is the one in Chapter XII, manuscript pages 88--89, and not the differently defined notion in \emph{Pursuing Stacks}.

\medskip\noindent Every mathematical statement below is marked Proved, Refuted, Known, or Open. ``Open'' records a statement not resolved in this paper; it does not assert that the problem is open in the literature. We claim no priority. The dated additions in the supplied prompt are retained as source metadata; the production date above is the session date.
\tableofcontents
\clearpage
\section{Definitions and the formal theory}
\subsection{The axioms and the variance}

\begin{definition}[Known: fundamental localizer]\label{A:localizer}
We use the following form of Grothendieck's three axioms. A class $W$ of functors in $\Cat$ is a fundamental localizer if it satisfies LF1--LF3. LF1 means that identities belong to $W$, that $W$ has the two-out-of-three property, and that whenever $i:X\to Y$ and $r:Y\to X$ satisfy $ri=\id_X$ and $ir\in W$, then $r\in W$. LF2 says that $A\to e$ belongs to $W$ whenever $A$ has a terminal object. LF3 says that for a commutative triangle $A\xrightarrow{u}B\xrightarrow{v}C$, if every functor $u/c:A/c\to B/c$ belongs to $W$, then $u\in W$. Here both slices use the indicated functors to $C$. A category is $W$-aspherical when its unique functor to $e$ belongs to $W$.
\end{definition}

This is exactly \citet[Definitions 1.1.1 and 1.1.3]{Cisinski2004Minimal}. In LF1, two-out-of-three then also gives $i\in W$. LF3 is the relative Theorem A axiom, not merely its special case $C=B$. We do not add a filtered-colimit axiom. The ordinary finality used throughout this paper is a property of connected components and is independent of $W$.

\begin{definition}[Known: manuscript Definitions 1 and 2]\label{A:definitions}
For a presheaf $H:U^{\op}\to\Set$, write $U/H$ for its category of elements: its objects are $(u,h)$, and a morphism $(u,h)\to(v,k)$ is an arrow $a:u\to v$ with $H(a)k=h$. Given $p:U\to U'$, the functor $p_!:\Set^{U^{\op}}\to\Set^{U'^{\op}}$ is ordinary left Kan extension along $p^{\op}$. The presheaf $H$ is totally $W$-aspherical if $U'/p_!H$ is $W$-aspherical for every small $U'$ and every $p$. A category $A$ is totally $W$-aspherical if its terminal presheaf $1_A$ is so. We denote this property by $T_W(A)$.

For $f:X\to Y$, put $H_y(x)=\Hom_Y(fx,y)$. The functor $f$ is \emph{strictly} $W$-aspherical if every presheaf $H_y$ is totally $W$-aspherical, and \emph{totally} $W$-aspherical if every category $X/y=X/H_y$ satisfies $T_W$. For a sieve $i:U\hookrightarrow X$, let $\overline U$ be the full cosieve on the objects $x$ for which some arrow $u\to x$, $u\in U$, exists. The sieve is $W$-dwy\'erienne, respectively totally $W$-dwy\'erienne, if $U\to\overline U$ is strictly, respectively totally, $W$-aspherical. A cosieve has either property when its opposite sieve has that property.
\end{definition}

These are the definitions at XII ms.~88--89 and 109 \citep{GrothendieckDerivateursXIIb}. In particular, the phrase ``strictement'' concerns the presheaf and ``totalement'' concerns its category of elements. The tests at $x\in U$ are representable; at an object outside $\overline U$ the crossing presheaf is empty and is not tested. For a covariant functor $K:A\to\Set$ we instead write $\El_A(K)$; its arrows satisfy $K(a)(k)=k'$. These two categories of elements have opposite variance and cannot be interchanged.

\subsection{Comprehensive factorization}

\begin{proposition}[Proved: the explicit comprehensive factorization]\label{A:comprehensive}
For $f:A\to B$, define the presheaf $P_f:B^{\op}\to\Set$ by
\[
 P_f(b)=\pi_0(b\backslash A).
\]
There is a canonical factorization in $\Cat$
\[
 A\xrightarrow{q_f}D_f:=B/P_f\xrightarrow{\rho_f}B,
 \qquad q_f(a)=(fa,[a,\id_{fa}]),
\]
where $q_f$ is final and $\rho_f$ is a discrete fibration. It is the factorization into a final functor followed by a discrete fibration, unique up to a unique isomorphism over $B$ compatible with the map from $A$. Moreover $P_f\cong f_!1_A$ naturally in $b$.
\end{proposition}

\begin{proof}
For $\beta:b'\to b$, precomposition gives a functor $b\backslash A\to b'\backslash A$ and hence the required presheaf map. An arrow $b'\to b$ lifts uniquely to an arrow with prescribed target $(b,c)$ in $D_f$, with source $(b',P_f(\beta)c)$; thus $\rho_f$ is a discrete fibration. The comma category $(b,c)\backslash A$ for $q_f$ identifies with the full connected component $c$ of $b\backslash A$: its objects are exactly arrows $\alpha:b\to fa$ belonging to $c$, and its arrows are the same commuting triangles. It is nonempty and connected. This proves finality, including all empty-component cases.

The pointwise Kan formula gives
\[
 (f_!1_A)(b)=\colim_{(a,\,b\to fa)\in(b\backslash A)^{\op}}\{*\}
 =\pi_0(b\backslash A).
\]
The opposite on the indexing category does not affect its components; it matters for a nonconstant presheaf.

For uniqueness, let $f=rq$ with $q:A\to E$ final and $r:E\to B$ a discrete fibration. Write $E\cong B/Q$ over $B$. For each $b$, the category $b\backslash E$ is a disjoint union, indexed by $Q(b)$, of categories with initial objects: the component indexed by $h$ has initial object $((b,h),\id_b)$. Thus $r_!1_E\cong Q$. Since $q_!1_A\cong1_E$, composition of left adjoints gives $f_!1_A\cong Q$, with its unit giving the compatible map from $A$. This determines the asserted isomorphism and its uniqueness.
\end{proof}

This is the contravariant form of the comprehensive factorization of \citet[Theorem 3]{StreetWalters1973}; that original source was not available as full text here. The opened modern treatment \citep[Section 1.3, Theorem 1.7, and Remark 1.9]{BergerKaufmann2017} explicitly distinguishes the final/discrete-fibration version from its dual. Our proof supplies the needed version independently.

\begin{corollary}[Proved: final-image reformulation and closure]\label{A:final-images}
For any fundamental localizer $W$ and any small category $A$, $T_W(A)$ is equivalent to the assertion that every target $B$ of an ordinary final functor $A\to B$ is $W$-aspherical. If $T_W(A)$ and $A\to B$ is final, then $T_W(B)$.
\end{corollary}

\begin{proof}
If $f$ is final, $f_!1_A=1_B$, so its category of elements is $B$. Conversely every $D_f$ in Proposition~\ref{A:comprehensive} is the target of a final functor from $A$. Finally, final functors compose: for $A\xrightarrow{f}B\xrightarrow{g}C$, the identities $f_!1_A=1_B$ and $g_!1_B=1_C$ give $(gf)_!1_A=1_C$. Apply the first equivalence to this composite.
\end{proof}

\begin{proposition}[Proved: categories of elements map finally under Kan extension]\label{A:elements-Kan}
Let $H:U^{\op}\to\Set$ and $\pi:U/H\to U$. For every $p:U\to U'$ there is a canonical final functor
\[
 q_{p,H}:U/H\longrightarrow U'/p_!H
\]
over $p$. In particular, $T_W(U/H)$ implies total $W$-asphericity of $H$. Total presheaves are preserved by every ordinary $p_!$.
\end{proposition}

\begin{proof}
The lift description in Proposition~\ref{A:comprehensive} gives $\pi_!1_{U/H}\cong H$. More explicitly, $u\backslash(U/H)$ has components indexed by $h\in H(u)$, with initial object $((u,h),\id_u)$ in each component. Therefore
\[
 (p\pi)_!1_{U/H}\cong p_!H.
\]
The asserted functor is $q_{p\pi}$ of Proposition~\ref{A:comprehensive}. The first consequence follows from Corollary~\ref{A:final-images}. If $H$ is total and $r:U'\to U''$, then $r_!p_!H\cong(rp)_!H$, so $p_!H$ satisfies every required test. All these isomorphisms are natural isomorphisms of presheaves, and hence give isomorphisms of their categories of elements over the indicated base.
\end{proof}

\subsection{What filteredness supplies}

\begin{lemma}[Proved: final images of filtered categories are filtered]\label{A:filtered-final}
If $A$ is nonempty and filtered and $f:A\to B$ is final, then $B$ is nonempty and filtered.
\end{lemma}

\begin{proof}
Given $b,c\in B$, choose $b\to fa$ and $c\to fa'$ using nonemptiness of their comma categories. A common successor of $a,a'$ gives a common successor of $b,c$. Given parallel arrows $s,t:b\rightrightarrows c$, choose $v:c\to fa$. The two objects $(a,vs)$ and $(a,vt)$ in $b\backslash A$ are connected by a finite zigzag. A filtered cocone over the finitely many objects and arrows in this zigzag gives an arrow $w:a\to a''$ with $f(w)vs=f(w)vt$. To justify the same arrow $w$ at the two appearances of $a$, first choose a cocone for the finite diagram with the occurrences distinguished, and then equalize the finitely many resulting parallel arrows out of repeated objects. Thus $f(w)v:c\to fa''$ equalizes $s,t$ in $B$.
\end{proof}

\begin{theorem}[Known: filtered categories are total for every fundamental localizer]\label{A:filtered}
Every small filtered category is totally $W$-aspherical for every fundamental localizer satisfying exactly LF1--LF3 of Definition~\ref{A:localizer}. No additional filtered-colimit axiom is required.
\end{theorem}

\begin{proof}
A filtered category has weakly contractible nerve. Indeed every finite collection of simplices in its nerve uses finitely many objects and arrows, and a compatible cocone gives a contraction of their inclusion into the nerve. Equivalently, after subdivision any spherical simplicial representative and any finite homotopy problem is killed by such a cone. Nonemptiness handles degree zero. This proves weak contractibility, not the existence of a single global contraction natural transformation for an infinite category.

\citet[Theorem 2.2.11, p.~136]{Cisinski2004Minimal} proves that $\Winfty$ is the smallest fundamental localizer: it satisfies LF1--LF3 and is contained in every class satisfying those axioms. Hence every filtered category is $W$-aspherical. Lemma~\ref{A:filtered-final} and Corollary~\ref{A:final-images} now give totality. Cisinski's theorem is the cited external theorem in this proof; the stronger historical axiom about inductive limits is unnecessary under the modern three-axiom definition.
\end{proof}

\begin{remark}[Proved: invariance that is and is not automatic]\label{A:invariance}
Totality is invariant under equivalence of categories and under isomorphism of presheaves, and it passes from $T_W$ to $T_{W'}$ when $W\subseteq W'$. For $W\subseteq W_0$, it implies that the category of elements in the identity test is nonempty and connected. It is not in general invariant under $\Winfty$-equivalence, as the outgoing-fork example in the audit below shows. For the universal localizer consisting of all functors, every category and every presheaf, including the empty ones, is total.
\end{remark}

\begin{proof}
An equivalence and its quasi-inverse are final, so Corollary~\ref{A:final-images} applies in both directions. The other assertions follow directly from the quantifiers, with the noninvariance witnessed below by a specified functor to the point.
\end{proof}

\subsection{Propositions 7--9 and the gluing axioms}

\begin{definition}[Known: the full gluing axioms of XII]\label{A:audit-gluing}
For a fundamental localizer $W$, the axiom called W(7 bis) at XII ms.~75--76 asserts the following biconditional for every small category $X$ covered by two full sieves $U,V$ \citep{GrothendieckDerivateursXIIb}:
\[
 (U\cap V\longrightarrow V)\in W
 \quad\Longleftrightarrow\quad
 (U\longrightarrow X)\in W.
\]
The symmetric biconditional is included. W(7), at ms.~76--77, concerns a morphism of two such cover squares. Of the three conditions that the map on the intersection belongs to $W$, that both maps on the pieces belong to $W$, and that the map on the union belongs to $W$, any two must imply the third. These are stronger assertions than forward homotopy gluing.
\end{definition}

\begin{lemma}[Proved: a categorical model for a one-point attachment]\label{A:audit-cylinder}
Let $U$ be a small category, $H:U^{\mathrm{op}}\to\mathbf{Set}$ a presheaf, and $E=U/H$ with projection $\alpha:E\to U$. Let $X$ be the category with full sieve $U$, one complementary object $t$ with only its identity endomorphism, and crossing arrows $\operatorname{Hom}_X(u,t)=H(u)$. There is a category $M$, covered by two full sieves $M_U,M_t$ with intersection $E$, and a functor $F:M\to X$, such that
\[
 r:M_U\longrightarrow U\ \in W_\infty,
 \qquad M_t=E^{\triangleright},\qquad F\in W_\infty.
\]
The map $Br$ is a homotopy equivalence induced by a specified adjunction. Consequently $BX$ is the homotopy pushout of $BU\leftarrow BE\to *$.
\end{lemma}

\begin{proof}
Take disjoint copies of $E$ and $U$ and a new object $t$. Retain all their internal arrows, set
\[
 \operatorname{Hom}_M(e,u)=\operatorname{Hom}_U(\alpha e,u),
 \qquad \operatorname{Hom}_M(e,t)=\{*\},
\]
and allow no other crossing arrows. Composition uses composition in $U$ and the unique arrows into $t$. Put $M_U=E\cup U$ and $M_t=E\cup\{t\}$. Both are sieves and their intersection is exactly $E$. The functor $r:M_U\to U$ sending $e$ to $\alpha e$ and fixing $U$ is left adjoint to the full inclusion $U\to M_U$: the required Hom-set bijections are the defining formula above. Its counit is the identity and its unit is given by $e\to\alpha e$, so $Br$ is a homotopy equivalence.

Define $F:M\to X$ by $F(e)=\alpha e$, $F(u)=u$, and $F(t)=t$. For $e=(u,h)$, send its unique arrow to $t$ to $h\in H(u)$. Every other crossing arrow is sent to its defining arrow of $U$. The relation $H(g)(k)=h$ defining an arrow $(u,h)\to(v,k)$ in $E$ proves functoriality at $t$.

For $u\in U$, the comma category $M/u$ has terminal object $(u,\mathrm{id}_u)$. The only remaining comma category is $M/t$. Its $E$-layer objects are triples $(u,h,k)$ with $h,k\in H(u)$: $(u,h)$ is the object of $E$ and $k:F(u,h)\to t$ is the comma arrow. Its $U$-layer objects are $(u,k)$ with $k\in H(u)$, and there is one further object $(t,\mathrm{id}_t)$. An $E$-layer object has an arrow to the last object precisely when $h=k$.

Let $N\subseteq M/t$ contain the $U$-layer, the last object, and just the diagonal triples $(u,h,h)$ in the $E$-layer. Define $S:M/t\to N$ to fix $N$ and send an off-diagonal $(u,h,k)$ to the $U$-layer object $(u,k)$. Send each arrow to the arrow with the same underlying morphism of $U$, or to the unique arrow into the last object. This prescription is defined on every arrow: an arrow of $E$-layers whose target is diagonal has diagonal source, since restriction of equal sections gives equal sections. There is a natural transformation $\mathrm{Id}_{M/t}\Rightarrow\iota S$ whose only nonidentity components are $(u,h,k)\to(u,k)$ with underlying arrow $\mathrm{id}_u$. Thus $N\hookrightarrow M/t$ is a homotopy equivalence.

The category $N$ consists of two copies of $E$, the first mapping to the second by the categorical mapping cylinder of the identity, and arrows from the first copy to $t$. The endofunctor $R:N\to N$ sending the second copy to the first and fixing the first copy and $t$ has natural transformations
\[
 \mathrm{Id}_N\ \Longleftarrow\ R\ \Longrightarrow\ \mathrm{const}_t.
\]
The first uses the cylinder arrows, and the second uses the unique arrows into $t$. Hence $N$, and therefore $M/t$, is contractible. Quillen's Theorem A, in its overcategory form \citep[\S1, Theorem~A]{Quillen1973}, proves $F\in W_\infty$.

Every nerve simplex of $M$ lies in $M_U$ or $M_t$, because these are a sieve cover. Their nerve inclusions are monomorphisms of simplicial sets and their realizations are cofibrations. The literal union $BM=BM_U\cup_{BE}BM_t$ is therefore a homotopy pushout. Applying $BF$ and $Br$ proves the last assertion, with its specified comparison maps. The argument includes $H=\varnothing$ and $U=\varnothing$.
\end{proof}

\begin{proposition}[Proved: Proposition 7 with its stated additional axiom]\label{A:audit-prop7}
Let $W$ be a fundamental localizer satisfying the full W(7 bis), and let $U\subseteq X$ be a full sieve. Write $T=X\setminus U$, and let $\overline U$ be the full cosieve generated by $U$. The following are equivalent: the inclusion $U\to\overline U$ is $W$-aspherical; for every functor $g:T_1\to T$, replacing the complement by $T_1$ and the crossing bifunctor by $(u,s)\mapsto\operatorname{Hom}_X(u,g(s))$ gives an inclusion $U\to\overline U_1$ in $W$; and this last condition holds for every point $g:e\to T$ whose value belongs to $\overline U\setminus U$.

For a one-point complement, the corollary at XII ms.~85 is likewise valid under full W(7 bis): $U\to X$ belongs to $W$ if and only if $U/t$ is $W$-aspherical.
\end{proposition}

\begin{proof}
We first prove the one-point assertion. Use Lemma~\ref{A:audit-cylinder} for $H(u)=\operatorname{Hom}_X(u,t)$, so $E=U/t$. The two functors $r:M_U\to U$ and $F:M\to X$ belong to every fundamental localizer, since they belong to $W_\infty$. By two-out-of-three and the equality $F|_{M_U}=ir$, the map $M_U\to M$ belongs to $W$ exactly when $i:U\to X$ does. W(7 bis), applied to the sieve cover $M=M_U\cup M_t$, identifies this with $E\to M_t$ belonging to $W$. Since $M_t$ has terminal object $t$, two-out-of-three identifies the latter condition with $E\to e$ belonging to $W$.

For the three conditions in the proposition, an object $s$ of the replaced complement has comma category $U/s=U/g(s)$, literally with the same objects, arrows and composition. Objects of $U$ have comma categories with terminal objects. Thus the first condition implies $W$-asphericity, hence membership in $W$, of each replaced inclusion by the relative Theorem A axiom. The second condition includes the third. The one-point assertion shows that the third condition makes $U/t$ $W$-aspherical for every $t\in\overline U\setminus U$, which is the first. If this set is empty, all point tests and the comma condition are vacuous and $U=\overline U$; thus no nonemptiness hypothesis has been suppressed.
\end{proof}

\begin{example}[Proved: an explicit finite acyclic category with nontrivial fundamental group]\label{A:audit-acyclic}
There is a category $A$ with $35$ objects and $159$ arrows, including identities, such that $BA$ is connected and integrally acyclic but has nontrivial fundamental group. We specify its complete finite presentation through face operators.

Start with the two words
\[
 w_0=aaaaaBBB,\qquad w_1=bbbBABA,
\]
where capital letters denote inverses, with lengths $n_0=8,n_1=7$. A finite semisimplicial set $K$ has vertices $v,c_0,c_1$, edges $a,b$ with both endpoints $v$, and edges $r_{k,i}:c_k\to v$ for $k\in\{0,1\}$ and $i\in\mathbb Z/n_k$. It has triangles $t_{k,i}$ with ordered vertices $(c_k,v,v)$. If the $i$th letter of $w_k$ is the positive letter $\ell$, define
\[
 (d_0t_{k,i},d_1t_{k,i},d_2t_{k,i})
 =(\ell,r_{k,i+1},r_{k,i});
\]
if that letter is $\ell^{-1}$, define this triple to be $(\ell,r_{k,i},r_{k,i+1})$. The endpoint faces of every $r_{k,i}$ are $(v,c_k)$, and those of $a,b$ are $(v,v)$. These prescriptions satisfy the semisimplicial identities. There are no simplices above dimension two.

The objects of $A$ are the simplices just listed. A morphism $\sigma\to\tau$ is an increasing injection $\theta:[\dim\sigma]\hookrightarrow[\dim\tau]$ for which $\theta^*\tau=\sigma$; composition is composition of the injections. In particular, distinct occurrences of the same face give distinct arrows. These definitions specify every object, arrow and composite of $A$.
\end{example}

\begin{proof}
The realization $|NA|$ is canonically homeomorphic to the barycentric subdivision of $|K|$. Indeed, each nondegenerate nerve simplex is a flag of proper face injections; the barycentres of the nested faces specify its usual barycentric simplex in the final target simplex. These simplexes subdivide each characteristic simplex, and the constructions agree on every face identification, including repeated faces.

Geometrically $K$ triangulates the presentation complex with one vertex, two loop edges and two two-cells attached by $a^5b^{-3}$ and $b^3(ab)^{-2}$. For each attaching polygon the triangles above are its cone triangulation, with centre $c_k$ and radial edges $r_{k,i}$; reversing the order of its last two triangle vertices for a negative letter gives exactly the indicated attaching map. This is the acyclic presentation complex of \citet[Example~2.38, p.~142]{Hatcher2002}, now with explicit finite categorical data.

The cellular differential from degree two to degree one is
\[
 d_2=\begin{pmatrix}5&-2\\-3&1\end{pmatrix},
 \qquad \det(d_2)=-1.
\]
The complex is connected, $d_1=0$, and there are no higher cells. Hence its reduced integral homology vanishes. To verify nontrivial fundamental group, send
\[
 a\longmapsto(1\ 4\ 3\ 2\ 5),\qquad
 b\longmapsto(1\ 2\ 3)
\]
in $A_5$, using right-to-left permutation composition. Both relator words have identity image: the permutations have orders five and three, and $ab=(1\ 5)(3\ 4)$. Since the image of $a$ is not the identity, the presented group, and therefore $\pi_1(BA)$, is nontrivial.

The program \texttt{code/ag\_finite\_examples.py} generates all arrows and composites from these face operators and checks all category axioms. The normalized nerve chain groups have ranks $(35,124,90)$; the two boundary matrices have ranks $34,90$ and all their nonzero Smith factors are $1$. Thus the independent integral calculation gives $H_0=\mathbb Z$ and $H_i=0$ for $i>0$. It also checks the two relator permutations. Full arrow and composition tables are in \texttt{code/ag\_category\_data.json}.
\end{proof}

\begin{theorem}[Refuted: the unqualified one-point converse and W(7 bis) for $W_\infty$]\label{A:audit-prop7counter}
The converse ``$U\to X$ is a weak equivalence, therefore $U/t$ is weakly contractible'' fails for a finite sieve inclusion with a one-object complement. In particular $W_\infty$ satisfies neither the full W(7 bis) nor the full W(7), and Proposition~7 cannot be read with the additional axiom deleted.
\end{theorem}

\begin{proof}
Take $A$ from Example~\ref{A:audit-acyclic}. Adjoin two distinct terminal cone points $u,t$ over $A$, with no arrows between $u$ and $t$, and call the resulting category $X$. Thus every object $a\in A$ has one arrow to each cone point; composing any arrow of $A$ with such an arrow gives the unique corresponding arrow. Let $U=A\cup\{u\}$ and $V=A\cup\{t\}$. Both are full sieves, $X=U\cup V$, $U\cap V=A$, and $U/t\cong A$.

The nerves of $U$ and $V$ are cones on $NA$, so $BU$ and $BV$ are contractible and $BX\cong\Sigma BA$. Since $BA$ is connected, van Kampen gives $\pi_1(BX)=1$. The reduced Mayer--Vietoris sequence identifies $\widetilde H_i(BX;\mathbb Z)$ with $\widetilde H_{i-1}(BA;\mathbb Z)$, which vanishes. The simply connected acyclic CW complex $BX$ is contractible by the homological Whitehead theorem \citep[Corollary~4.33, p.~367]{Hatcher2002}. Hence $U\to X$ belongs to $W_\infty$, whereas $U/t\cong A$ is not $W_\infty$-aspherical. The cover square has $U\to X$ in $W_\infty$ and $A\to V$ outside $W_\infty$, contradicting W(7 bis). Mapping this cover square to the constant point square gives a direct failure of W(7): the two piece maps and the union map are weak equivalences, while the intersection map is not.

The category $X$ has $37$ objects and $231$ arrows. The same program checks its composition, both sieve conditions, and its integral homology. The topological argument, including the nontrivial fundamental-group certificate for $BA$, supplies the homotopy information that homology computation alone cannot establish.
\end{proof}

\begin{proposition}[Proved: the homology localizers satisfy the full gluing axioms]\label{A:audit-homologygluing}
For every commutative coefficient ring $R$, the class $W_R$ of nerve $R$-homology equivalences satisfies W(7), and consequently W(7 bis).
\end{proposition}

\begin{proof}
For a sieve cover $X=U\cup V$, any chain of arrows in $X$ lies in a piece containing its last vertex. Thus the normalized nerve chains give a natural short exact sequence
\[
 0\longrightarrow C_*(N(U\cap V);R)
 \xrightarrow{(i_*,-j_*)}C_*(NU;R)\oplus C_*(NV;R)
 \longrightarrow C_*(NX;R)\longrightarrow0.
\]
Exactness is checked on the bases of nondegenerate chains. For a morphism of cover squares, the mapping cones of the vertical chain maps again form a short exact sequence of complexes. The middle cone is the direct sum of the cones for the two piece maps. The long exact homology sequence proves that acyclicity of any two of the intersection cone, the pair of piece cones, and the union cone implies acyclicity of the third. A chain map is a homology equivalence exactly when its cone is acyclic, proving W(7). To obtain W(7 bis), map the given cover square to the cover of $X$ by $X$ and $V$. The maps on $V$ and $X$ are identities; W(7) then identifies the equivalence condition on $U\cap V\to V$ with that on $U\to X$.
\end{proof}

\begin{proposition}[Proved: Propositions 8 and 9(a)]\label{A:audit-prop8and9a}
For every fundamental localizer $W$, if $U/H$ is totally $W$-aspherical, then $H$ is totally $W$-aspherical. Every small category with a terminal object is totally $W$-aspherical, and every representable presheaf is totally $W$-aspherical.
\end{proposition}

\begin{proof}
For $\alpha:U/H\to U$, there is a canonical isomorphism $\alpha_!1_{U/H}\cong H$. Consequently, for every $p:U\to U'$, one has $p_!H\cong(p\alpha)_!1_{U/H}$; the first assertion is exactly the defining test for $U/H$ applied to $p\alpha$.

If $v_0$ is terminal in $V$ and $f:V\to B$, then
\[
 f_!1_V\cong B(-,f(v_0)).
\]
For an explicit verification, any object $(v,\beta:b\to f(v))$ of $b\backslash V$ maps to $(v_0,f(v\to v_0)\beta)$. This gives a retraction onto the discrete set $B(b,f(v_0))$, with a natural transformation from the identity to the inclusion followed by the retraction. Thus the stated component bijection is natural in $b$. Its category of elements is $B/f(v_0)$, with terminal object $(f(v_0),\mathrm{id})$, so the terminal-object axiom for $W$ applies. Finally $U/U(-,u)\cong U/u$ has a terminal object, and the first assertion gives the representable case.
\end{proof}

\begin{theorem}[Refuted: Proposition 9(b), its retraction repair, and the fork corollary]\label{A:audit-prop9b}
For $W=W_\infty$, and also for $W=W_R$ with any nonzero commutative ring $R$, a union of two totally $W$-aspherical full sieves with totally $W$-aspherical intersection need not be totally $W$-aspherical, even if both pieces retract onto the intersection. In the homology case this refutes the statements under their full stated W(7 bis) hypothesis.
\end{theorem}

\begin{proof}
Let $\Psi=(z\to a,z\to b)$. Its full sieves $\{z\to a\}$ and $\{z\to b\}$ and their intersection $\{z\}$ all have terminal objects, so Proposition~\ref{A:audit-prop8and9a} makes them totally $W$-aspherical. The constant functors from either arrow category to $\{z\}$ are retractions. Nevertheless the explicit ordinary final inclusion $j:\Psi\to C$ constructed and verified in Theorem~\ref{G:counterexample} has $BC\cong S^2$. It follows that $\Psi$ is neither totally $W_\infty$-aspherical nor totally $W_R$-aspherical when $R\ne0$. By Proposition~\ref{A:audit-homologygluing}, the latter localizer satisfies the stated gluing axiom. Thus the axiom does not rescue the union conclusion or either retraction version.

The failed step in the manuscript can be inspected before taking nerves. The subterminal presheaves of the two branches and their intersection are the representables $\Psi(-,a)$, $\Psi(-,b)$, and $\Psi(-,z)$. Their images under $j_!$ are the corresponding representables on $C$. Evaluating the first inclusion at $x$ gives the noninjective map
\[
 C(x,z)=\{\alpha,\beta\}\longrightarrow
 C(x,a)=\{\gamma\}.
\]
A retraction between full subcategories does not define a retraction between their subterminal presheaves: at an object outside a proper sieve, the latter would require a map from a singleton to the empty set. The source's attempted passage of categorical retractions through $j_!$ is therefore invalid, and the exhibited final image proves that its conclusion is false.
\end{proof}

\section{Collages, the three leads, and final-image tests}
\subsection{The collage test}

\begin{proposition}[Proved: Lead 1, including the unrestricted homotopy type]\label{B:collage}
Let $A$ be small and $K:A\to\Set$ covariant. Form the category $A_K$ by adjoining an object $x$, with $A_K(x,a)=K(a)$, $A_K(a,x)=\varnothing$, and $A_K(x,x)=\{\id_x\}$. Composition of $h:x\to a$ with $v:a\to b$ is $K(v)h$. Put $E=\El_A(K)$ and let $\pi:E\to A$ forget the element. Then $i:A\hookrightarrow A_K$ is final if and only if $E$ is nonempty and connected. Without any contractibility assumption on $BA$, there is a natural homotopy-pushout description
\[
 B(A_K)\simeq\operatorname{hocolim}\bigl(BA\xleftarrow{B\pi}BE\longrightarrow *\bigr).
\]
More precisely, the displayed realization is the mapping cone of the specified map $B\pi:BE\to BA$, not a construction determined only by the two abstract homotopy types.
\end{proposition}

\begin{proof}
For $a\in A$, the category $a\backslash A$ for the inclusion has initial object $(a,\id_a)$. For the new object $x$, the comma category is precisely $\El_A(K)$. This proves finality with the stated edge cases.

Write $CE=\Delta^0*NE$ for the simplicial cone with its new vertex first. There is an isomorphism of simplicial sets
\[
 N(A_K)\cong NA\amalg_{NE}(\Delta^0*NE).
\]
To check it, a nondegenerate simplex outside $NA$ starts with $x$, followed by an arrow $h:x\to a_0$ and a string $a_0\to\cdots\to a_n$. The data are exactly a simplex of $NE$, because the remaining elements of the string are determined from $h$ by $K$. Deleting the first vertex forgets $h$ and is the map $N\pi$; all other faces and degeneracies agree with the cone. This also proves the assertion for repeated identity arrows. Realization preserves this pushout, and $BE\hookrightarrow CBE$ is a CW cofibration. Replacing the contractible cone by a point identifies it with the indicated homotopy pushout. If $E$ is empty, the result is $BA\amalg *$, as it should be.
\end{proof}

\begin{corollary}[Proved: the integral obstruction and its converse for one collage]\label{B:acyclic-test}
Suppose $BA$ is weakly contractible and $E=\El_A(K)$ is nonempty and connected. Then
\[
 B(A_K)\simeq\Sigma BE,
 \qquad
 \widetilde H_n(B(A_K);\mathbb Z)\cong
 \widetilde H_{n-1}(BE;\mathbb Z)
 \quad(n\geq1).
\]
Consequently $B(A_K)$ is weakly contractible if and only if $BE$ is integrally acyclic. In particular, if $T_{\Winfty}(A)$, every such $E$ is acyclic. The same necessary assertion with coefficients in a commutative ring $R$ holds under $T_{W_R}(A)$.
\end{corollary}

\begin{proof}
The map $BE\to BA$ has contractible target, so its homotopy cofiber is the unreduced suspension of $BE$. The suspension of a connected CW complex is simply connected, and its reduced homology is the shifted reduced homology of that complex. A simply connected acyclic CW complex is weakly contractible by the Hurewicz theorem, applied at the least hypothetical nonzero homotopy group. For $T_{\Winfty}(A)$, both $A$ and its final image $A_K$ are weakly contractible. For $W_R$, use the homology long exact sequence of the same cofiber; no assertion of simple connectivity is needed in that case.
\end{proof}

\begin{proposition}[Proved: finite and transfinite iteration]\label{B:iteration}
If $T_W(A)$ and $\El_A(K)$ is nonempty and connected, then $T_W(A_K)$. This remains true after any small transfinite succession of such one-object collages, taking unions at limit ordinals. Thus every connected covariant element-category test on every category obtained in this way is a necessary test, not just those on the starting category.
\end{proposition}

\begin{proof}
The successor assertion is Corollary~\ref{A:final-images}. For completeness, the inclusion of an earlier stage into a union at a limit ordinal is final. If $b$ first appears at a stage $\beta$, its comma category into a fixed earlier stage becomes nonempty connected by induction on $\beta$: successor steps are composition of final inclusions, and at a limit the comma category is the increasing union of the earlier comma categories. Once nonempty, all are connected, so their union is connected. Equivalently, for the inclusion of the original stage, an object $b$ belongs to some stage $\beta$, and no later collage adds arrows from an existing object to an old object, so $b\backslash A$ is already the comma category at that stage. The original inclusion into the full union is therefore final, and Corollary~\ref{A:final-images} applies. The smallness assumption ensures that the union is still a category in our universe.
\end{proof}

The element-category characterization in Section C gives a closed form for these iterated necessary conditions in the integral and weak-homotopy cases. It avoids any inference from isomorphisms of ind-objects or from unspecified weak equivalences between stages.

\subsection{The parallel-pair homology formula}

\begin{proposition}[Proved: the integral formula in Lead 2]\label{B:lead2}
Let $a,b:p\to q$ be a parallel pair in any small category $A$. Let $h_p=A(p,-)$ and let $H$ be the pointwise quotient of $h_p$ generated by $ca\sim cb$, for every $c:q\to r$ and every $r\in A$. Then $\operatorname{El}_A(H)$ is nonempty and connected and
\[
 H_1(B\operatorname{El}_A(H);\mathbb Z)\cong\mathbb Z/I,
\]
where $I\subseteq\mathbb Z$ is the subgroup generated by all integers $\sum_cn_c$ such that
\[
 \sum_cn_c(ca-cb)=0\quad\text{in }\mathbb Z[A(p,r)]
\]
for one object $r$ and a finitely supported family $(n_c)_{c\in A(q,r)}$. Equivalently, if $I=g\mathbb Z$ with $g\geq0$, this group is $\mathbb Z/g\mathbb Z$. This formula has no asphericity assumption on $A$.
\end{proposition}
\begin{proof}
Put $P_s=\mathbb Z[A(s,-)]\in\Ab^A$. These are projective: $\operatorname{Hom}_{\Ab^A}(P_s,F)=F(s)$ naturally, and evaluation is exact. Define $d:P_q\to P_p$ by $d(c)=ca-cb$, let $K=\ker d$ and $J=\operatorname{im}d$. The free-abelian-group functor preserves coequalizers, so we have exact sequences
\[
 0\to J\to P_p\to\mathbb Z[H]\to0,
 \qquad 0\to K\to P_q\to J\to0.
\]
Both $\colim P_q$ and $\colim P_p$ identify canonically with $\mathbb Z$, and $\colim d$ is zero. The long exact sequence of derived colimits for the first sequence, and right exactness applied to the second, give
\[
 L_1\colim_A\mathbb Z[H]
 \cong\colim_A J
 \cong\operatorname{coker}(\colim_AK\to\mathbb Z).
\]
For the first identification, the map $\colim J\to\colim P_p$ is zero because it is zero after the epimorphism $\colim P_q\to\colim J$. The image of $\colim K\to\mathbb Z$ is exactly the subgroup $I$ displayed in the statement.

The bar chain complex for $\mathbb Z[H]$ has basis a string $a_0\to\cdots\to a_j$ and an element of $H(a_0)$. These are exactly the simplices of $N\operatorname{El}_A(H)$, compatibly with every face. It therefore computes the homology on the left. Finally, every object of $\operatorname{El}_A(H)$ receives an arrow from $(p,[\id_p])$: choose a representative $f:p\to r$ of its element. Thus this category is connected and nonempty, although the arrow from that object need not be unique.
\end{proof}

\subsection{The smallest band obstruction}

\begin{proposition}[Proved: Lead 3]\label{B:right-zero}
Let $M=\{1,e,f\}$, with $1$ a unit and $uv=v$ for $u,v\in\{e,f\}$. Regard $M$ as a one-object category with multiplication $uv=u\circ v$. It is cofiltered, its classifying space is contractible, and it is not filtered. Let $K:M\to\Set$ have underlying set $\{h,i\}$, with $1$ acting identically and both $e,f$ acting constantly with value $i$. Then
\[
 B\El_M(K)\simeq S^1,
 \qquad B(M_K)\simeq S^2.
\]
The inclusion $M\to M_K$ is final. In particular $M$ is not totally $\Winfty$-aspherical and is not totally $W_R$-aspherical for any nonzero commutative ring $R$.
\end{proposition}

\begin{proof}
Associativity follows from the right-zero rule, with the unit inserted. For any $a,b\in M$, $ae=be=e$, proving cofilteredness. No $c\in M$ satisfies $ce=cf$, since $ce=e$ and $cf=f$. This disproves filteredness. The element $e$ defines a natural transformation from the constant endofunctor of the one-object category (sending every arrow to $1$) to its identity: the naturality equation is $me=e$. Hence $BM$ is contractible.

The action is a functor because compositions of the two constant functions are that same constant function. Its element category has two objects, $h$ and $i$. There are only the identity at $h$, two arrows $h\rightrightarrows i$ labelled $e,f$, and the three endomorphisms $1,e,f$ at $i$. Postcomposition fixes each of the two arrows $h\to i$, by the right-zero multiplication. Therefore this element category is itself the collage of the covariant constant two-element functor on the full one-object category at $i$. The category of elements of that constant functor is the disjoint union of two copies of $M$. Proposition~\ref{B:collage} now computes its classifying space as the homotopy cofiber of
\[
 BM\amalg BM\longrightarrow BM,
\]
which is $S^1$, since each copy of $BM$ is contractible. It is connected, so $M\to M_K$ is final. A second application of the proposition gives the suspension $S^2$. Its degree-two homology with nonzero coefficients $R$ is $R$.
\end{proof}

The complete multiplication and action tables, all arrows and composites in the two collages, and normalized nerve-chain computations through degree four are recorded in the accompanying program \path{code/root_monoid_checks.py} and its output. The all-degree assertion above follows from the two homotopy-pushout calculations; a truncated chain calculation alone would not prove it.

\begin{proposition}[Proved: the whole rectangular-band family]\label{B:bands}
Let $I,J$ be nonempty small sets, let $R_{I,J}=I\times J$ with $(i,j)(k,l)=(i,l)$, and let $M_{I,J}=R_{I,J}^{1}$ have a new identity adjoined. Then $M_{I,J}$ is filtered, and totally $\Winfty$-aspherical, if and only if $|J|=1$. It is cofiltered if and only if $|I|=1$. For every nonzero $R$, total $W_R$-asphericity has the same classification. If the band is empty, adjoining an identity gives the trivial monoid and all these properties hold.
\end{proposition}

\begin{proof}
If $J$ is a singleton, any band element is a left zero and postcomposition by it equalizes all arrows; Theorem~\ref{A:filtered} applies. If $J$ has at least two elements, choose a surjection $J\twoheadrightarrow\{e,f\}$ and forget the $I$ coordinate. This gives a surjective homomorphism $M_{I,J}\to M$ to Proposition~\ref{B:right-zero}'s monoid. Any homomorphism surjective on the underlying set is a full identity-on-objects functor. Its unique comma category is connected, since its identity object has an arrow to every object via a lift of the corresponding target monoid element. Hence it is final, and totality would pass to $M$, a contradiction. This works equally for all nonzero $R$. If $|J|>1$, two band elements with different second coordinates cannot be equalized by left multiplication; this also checks filteredness directly. The cofiltered assertion is the dual argument on first coordinates. No finiteness assumption is used.
\end{proof}

\subsection{Other final tests}

\begin{proposition}[Proved: quotient, localization, and idempotent-completion tests]\label{B:other-final}
Every quotient $q:A\to A/{\sim}$ by a congruence on the hom-sets, with unchanged objects, is final. Every localization $A\to A[S^{-1}]$ at a set of arrows is final. The canonical fully faithful functor $j:A\to\Kar(A)$ into its idempotent completion is final. Consequently a totally $W$-aspherical category has totally $W$-aspherical categories as all these specified images.
\end{proposition}

\begin{proof}
For a congruence quotient, an object of $b\backslash A$ is an arrow $b\to a$ in the quotient. It lifts to an arrow in $A$, giving a morphism from $(b,\id_b)$ to that object. The comma category is nonempty connected.

For a localization, an arrow $b\to a$ is represented by a finite word of original arrows and inverses of arrows in $S$. Start at $(b,\id_b)$ and read the word. A forward arrow gives a forward morphism in $b\backslash A$. An inverse $s^{-1}$ gives a backward morphism using the original arrow $s$ of $A$. This constructs a zigzag to $(a,\alpha)$, irrespective of different representatives of $\alpha$. Thus the comma category is connected and nonempty.

An object of $\Kar(A)$ is $(a,e)$ with $e^2=e$; an arrow $(a,e)\to(b,d)$ is $v:a\to b$ with $v=dve$. The comma category $(a,e)\backslash A$ has distinguished object $(a,e:(a,e)\to j(a))$. Any object $(b,v)$ receives a morphism from it, namely $v:a\to b$, since $ve=v$. This again proves connectedness. Applying Corollary~\ref{A:final-images} gives the consequences.
\end{proof}

\begin{proposition}[Proved: totality is invariant under idempotent completion]\label{B:karoubi}
For every fundamental localizer $W$, $T_W(A)$ holds if and only if $T_W(\Kar(A))$.
\end{proposition}

\begin{proof}
Only the converse needs proof. Suppose $\Kar(A)$ is total and $f:A\to B$ is final. We claim that $\Kar(f):\Kar(A)\to\Kar(B)$ is final. Fix $(b,d)\in\Kar(B)$. Every arrow $\alpha:(b,d)\to\Kar(f)(a,e)$ gives a comma morphism to the object $(j(a),\alpha)$ via the splitting arrow $(a,e)\to j(a)$. For objects of this last form, finality of $f$ connects the corresponding arrows in $b\backslash A$. Composing every arrow of that connecting zigzag on the right with $d$ keeps it in the subcategory of arrows $\alpha$ satisfying $\alpha d=\alpha$. Nonemptiness follows by choosing any $b\to fa$ and precomposing with $d$. Thus $\Kar(f)$ is final.

It follows that $\Kar(B)$ is $W$-aspherical. The functor $B\to\Kar(B)$ belongs to $\Winfty$: for its comma category considered above, the distinguished object maps naturally to every object, with components the underlying $v$, so there is a natural transformation from its constant functor to the identity. These comma categories are contractible, and the dual form of Quillen's Theorem A applies. Cisinski's minimality theorem places this functor in $W$. Two-out-of-three gives $B\to e\in W$, proving $T_W(A)$.
\end{proof}

\begin{remark}[Proved: what the tests actually force]\label{B:scope-tests}
A total category must therefore remain aspherical after every hom-set congruence and every localization, including its groupoid completion. Ordinary quotient tests do not require a faithful comparison map. Idempotent completion is stronger: it both preserves and reflects totality by the specified functors in Proposition~\ref{B:karoubi}. None of these assertions says that an arbitrary weak equivalence preserves totality. The common-successor obstruction, the affine parallel-arrow obstruction of Lead 2, and the complete element-category criterion appear in Section C.
\end{remark}

\begin{example}[Proved: the infinite fence and the nonuniformity of contractions]\label{B:fence}
Let $F$ have objects $\mathbb N$ and nonidentity arrows $2r\to2r+1$ for $r\geq0$ and $2r\to2r-1$ for $r\geq1$, with no other nonidentity arrows. Then $BF$ is contractible, but $F$ is not totally $\Winfty$-aspherical. Its identity endofunctor is not connected to a constant endofunctor by any finite zigzag of natural transformations.
\end{example}
\begin{proof}
There are no composable pairs of nonidentity arrows, and the realization of the nerve is exactly an infinite ray. The maximal objects $1$ and $3$ have no common successor, so Lemma~\ref{C:directed-affine} rules out totality. A component of any natural transformation between two functors into $F$ moves its object by at most one in the undirected path metric. A zigzag of length $m$ therefore moves each object by at most $m$. A constant endofunctor with value $c$ would require moving the object $n$ through distance $|n-c|$, which is unbounded. This proves the last assertion. We have reproduced the relevant argument of the supplied CATHOT C.5 \citep{CATHOTDraft2026}; no passage from a finite contraction to a uniform infinite contraction is used.
\end{proof}

\begin{corollary}[Proved: homotopical ordinary colimits of simplicial sets]\label{B:homotopical-colimits}
For a small nonempty connected $A$, total $\Winfty$-asphericity is equivalent to the following condition: for every $K:A\to\mathsf{sSet}$, the canonical augmentation from its simplicial-replacement homotopy colimit to its ordinary colimit is a weak equivalence. In particular, a total category's ordinary simplicial-set colimit functor preserves all objectwise weak equivalences.
\end{corollary}
\begin{proof}
The bisimplicial replacement has bidegree $(p,q)$ equal to
\[
 \coprod_{a_0\to\cdots\to a_p}K(a_0)_q.
\]
For a fixed $q$, this is the nerve of $\El_A(K_q)$, and its augmentation is the component map to $\colim_AK_q$. Theorem~\ref{C:classification} makes that augmentation a weak equivalence, component by component. The diagonal lemma for bisimplicial sets gives the asserted comparison for $K$. It follows by two-out-of-three that ordinary colimit preserves objectwise weak equivalences, since the simplicial replacement does: in each horizontal degree it is a coproduct of the given weak equivalences, and realization preserves these degreewise comparisons. Conversely apply the stated augmentation property to a discrete set-valued diagram $K=H$. Its source is $B\El_A(H)$ and its target is the discrete set of components, so Theorem~\ref{C:classification} gives totality. All comparison maps here are the canonical augmentations with the actual transition maps of $K$.
\end{proof}

\section{The proposed filteredness characterization}
\subsection{An intrinsic criterion and the failure of Lead 4}

\begin{definition}[Proved: the action criterion]\label{C:action-definition}
Let $R$ be a nonzero commutative ring with identity. We write $\mathcal P_R(A)$ for the following property of a small category $A$: for every covariant functor $H:A\to\Set$, each connected component $E$ of the category of elements $\operatorname{El}_A(H)$ is $R$-acyclic, meaning that $H_0(BE;R)=R$ and $H_j(BE;R)=0$ for $j>0$. No condition is imposed on the empty category of elements. We write $\mathcal P_\infty(A)$ when every such component is weakly contractible. These definitions use covariant, rather than presheaf, categories of elements: an arrow $(a,h)\to(a',h')$ is an arrow $v:a\to a'$ with $H(v)(h)=h'$.
\end{definition}

\begin{lemma}[Proved: homological Theorem A in the final variance]\label{C:homology-A}
Let $f:A\to B$ be a functor. If every $b\backslash A$ is nonempty and $R$-acyclic, then $Bf:BA\to BB$ induces an isomorphism on homology with coefficients in $R$. More generally, if every $b\backslash A$ is nonempty and its reduced $R$-homology vanishes in degrees $0,\ldots,d-1$, for an integer $d\geq1$, then $Bf$ induces isomorphisms in degrees less than $d$ and an epimorphism in degree $d$.
\end{lemma}
\begin{proof}
We use the bisimplicial set whose $(p,q)$-simplices are strings
\[
 b_0\longrightarrow\cdots\longrightarrow b_p
 \longrightarrow f(a_0),\qquad
 a_0\longrightarrow\cdots\longrightarrow a_q.
\]
The forgetful map to $NA$ is a weak equivalence after diagonal realization: for a fixed string in $A$, the remaining simplicial direction is the nerve of $B/f(a_0)$, which has terminal object $\id_{f(a_0)}$. The other projection is to $NB$. Taking normalized $R$-chains in both directions and filtering by $p$ gives the spectral sequence of this projection. Its row in degree $q=0$ is the chain complex of $NB$, because each comma category is connected; the rows in degrees $1,\ldots,d-1$ vanish after taking vertical homology. The first possible additional row has degree $d$. Thus the projection induces isomorphisms below degree $d$ and an epimorphism in degree $d$. The first-quadrant filtration is finite in each total degree, so it converges without a boundedness or finiteness assumption on the categories. If all positive rows vanish, it induces isomorphisms in every degree. The two projections commute with the canonical map $Bf$ up to the homotopy supplied by the displayed arrow $b_p\to f(a_0)$; consequently the comparison just computed is the map induced by $f$.
\end{proof}

\begin{lemma}[Proved: inheritance by categories of elements]\label{C:inherit-elements}
If $A$ satisfies $\mathcal P_R$, then every connected component $E$ of $\operatorname{El}_A(H)$, for every $H:A\to\Set$, also satisfies $\mathcal P_R$. The corresponding assertion holds for $\mathcal P_\infty$.
\end{lemma}
\begin{proof}
The elements lying in $E$ define a subfunctor $H_E\subseteq H$, because arrows preserve connected components. Let $K:E\to\Set$ be any covariant functor. Define
\[
 L(a)=\coprod_{h\in H_E(a)}K(a,h).
\]
For an arrow $v:a\to a'$ in $A$, its restriction to the summand indexed by $h$ is the map $K(v):K(a,h)\to K(a',H(v)h)$, landing in the indicated summand of $L(a')$. This defines a covariant functor $L:A\to\Set$. The map
\[
 ((a,h),k)\longmapsto (a,(h,k))
\]
is an isomorphism of categories $\operatorname{El}_E(K)\cong\operatorname{El}_A(L)$, including their specified arrows. Its connected components therefore have the required property.
\end{proof}

\begin{lemma}[Proved: integral acyclicity detects every cyclic subgroup]\label{C:cyclic-covers}
Suppose that a connected, nonempty small category $E$ satisfies $\mathcal P_{\mathbb Z}$. Then $BE$ is weakly contractible. Consequently $\mathcal P_{\mathbb Z}(A)$ and $\mathcal P_\infty(A)$ are equivalent for every small category $A$.
\end{lemma}
\begin{proof}
The constant singleton functor on $E$ shows that $BE$ is integrally acyclic. Fix an object $e_0$ and let $\Pi E$ be the groupoid obtained by inverting all arrows of $E$. The usual edge-path presentation of the fundamental group identifies $\operatorname{Aut}_{\Pi E}(e_0)$ with $G=\pi_1(BE,e_0)$. For any subgroup $K\leq G$, define
\[
 H_K(e)=\operatorname{Hom}_{\Pi E}(e_0,e)/K,
\]
where $K$ acts by precomposition, and an arrow $v:e\to e'$ acts by postcomposition with its image in $\Pi E$. All these maps are bijections. The projection $\operatorname{El}_E(H_K)\to E$ is therefore a categorical covering: each arrow has a unique lift with either a prescribed source or a prescribed target. On nerves this gives the usual covering of $BE$ whose monodromy is the transitive coset action of $G$. It is connected, and its fundamental group, at the lift represented by $\id_{e_0}$, is precisely $K$. The covering classification is also recorded in \citet[Theorem~1.38]{Hatcher2002}. These assertions follow directly here by lifting the edge-path presentation; a loop lifts to a loop exactly when its class belongs to $K$.

Take $K$ to be the cyclic subgroup generated by an arbitrary element of $G$. The hypothesis gives
\[
 0=H_1\bigl(B\operatorname{El}_E(H_K);\mathbb Z\bigr)
   =K_{\mathrm{ab}}=K.
\]
The identification of first homology with abelianization is \citet[Theorem~2A.1]{Hatcher2002}. Thus $G$ is trivial. A simply connected acyclic CW complex is weakly contractible: if a first nonzero homotopy group existed, the Hurewicz theorem would identify it with a nonzero homology group \citep[Theorem~4.32]{Hatcher2002}. A nerve realization is a CW complex, so this applies to $BE$.

For general $A$ satisfying $\mathcal P_{\mathbb Z}$, apply Lemma~\ref{C:inherit-elements} and the preceding argument to each component $E$ of each category of elements. The converse follows from homotopy invariance of integral homology.
\end{proof}

\begin{theorem}[Proved: characterization of total asphericity]\label{C:classification}
For every small category $A$, the following conditions are equivalent:
\begin{enumerate}
\item $A$ is totally $\Winfty$-aspherical in the sense of XII, Definition~1.
\item $A$ is totally $W_{\mathbb Z}$-aspherical.
\item $A$ is nonempty and connected, and satisfies $\mathcal P_{\mathbb Z}$.
\item $A$ is nonempty and connected, and satisfies $\mathcal P_\infty$.
\end{enumerate}
In particular, this is a criterion in terms of all covariant set-valued diagrams on $A$, involving no quantification over target categories or final functors. Moreover every connected component of every $\operatorname{El}_A(H)$ is itself totally $\Winfty$-aspherical whenever $A$ is totally $\Winfty$-aspherical.
\end{theorem}
\begin{proof}
The implication (1)$\Rightarrow$(2) follows because a weak homotopy equivalence induces an integral homology isomorphism. Assume (2). The identity functor is final, so $A$ is integrally acyclic and hence nonempty and connected. For any connected nonempty category of elements, the collage of Part~B is a final image of $A$. Its homotopy-pushout description gives the homology exact sequence
\[
 \cdots\to\widetilde H_{j+1}(BA_H;\mathbb Z)
 \to\widetilde H_j(B\operatorname{El}_A H;\mathbb Z)
 \to\widetilde H_j(BA;\mathbb Z)\to\cdots.
\]
Both outside terms vanish. For a component of an arbitrary $H$, first replace $H$ by its component subfunctor. This proves (3). Lemma~\ref{C:cyclic-covers} proves (3)$\Rightarrow$(4).

Assume (4), and let $f:A\to B$ be any ordinary final functor. For $b\in B$, the category $b\backslash A$ is the category of elements of the covariant functor $a\mapsto\operatorname{Hom}_B(b,f(a))$. It is connected and nonempty by finality, and therefore weakly contractible by the hypothesis. Quillen's Theorem~A \citep[Section~1, Theorem~A]{Quillen1973} implies that the specified map $Bf$ is a weak homotopy equivalence. The constant singleton functor in (4) makes $BA$ weakly contractible. Thus $BB$ is weakly contractible, proving (1) by the comprehensive-factorization criterion. The final assertion follows from Lemma~\ref{C:inherit-elements} and (4).
\end{proof}

\begin{proposition}[Proved: exact integral colimits are sufficient]\label{C:exact-sufficient}
Let $A$ be nonempty and connected. If $\colim_A:\Ab^A\to\Ab$ is exact, then $A$ is totally $\Winfty$-aspherical.
\end{proposition}
\begin{proof}
For $H:A\to\Set$, define $\Sigma_H:\Ab^{\operatorname{El}_A(H)}\to\Ab^A$ by
\[
 (\Sigma_HF)(a)=\bigoplus_{h\in H(a)}F(a,h),
\]
with the transition maps specified in the proof of Lemma~\ref{C:inherit-elements}. It is exact because direct sums of abelian groups are exact. The universal property of a cocone gives the canonical natural isomorphism
\[
 \colim_{\operatorname{El}_A(H)}F
   \cong\colim_A\Sigma_HF.
\]
Thus the colimit functor on each category of elements is exact, and the same is true on any one of its components, by extending diagrams by zero on the other components. The bar resolution computing derived colimits identifies
$L_j\colim_E\underline{\mathbb Z}$ with $H_j(BE;\mathbb Z)$: its term in degree $j$ is the free abelian group on strings of $j$ composable arrows of $E$, and its differential is the alternating sum of the nerve faces. Exactness therefore makes every component $E$ integrally acyclic. Now apply Theorem~\ref{C:classification}.
\end{proof}

\subsection{An explicit nonfiltered totally aspherical category}

\begin{definition}[Proved: the category of injective ordinals]\label{C:delta-definition}
Let $\Delta_{\mathrm{inj}}$ have objects $[n]=\{0,\ldots,n\}$ for $n\geq0$, and all strictly increasing functions as arrows, with ordinary composition of functions. This is a countable small category. In the next construction $\mathbb Z[S]$ denotes the free abelian group on a set $S$, and the augmentation of a finite linear combination of arrows is the sum of its coefficients.
\end{definition}

The exact-colimit example is attributed to \citet{Isbell1974Exact} in \citet[Section~2.3, after Theorem~2.5]{Husainov2002Dimension}; its bibliographic record was verified, but the original article's body was not available in this run. We do not use that unavailable body as a proof. The following explicit verification uses polynomial interpolation. A CC0 AI-authored exposition containing the construction was found during the literature search \citep[Section~4]{OpenMaths2026Derived}; we independently verify its algebraic identities here and by the accompanying program \texttt{cd\_affine\_polynomials.py}. Its provenance is recorded in the report, and it is not used as an authority.

\begin{lemma}[Proved: interpolation on coordinate hyperplanes]\label{C:interpolation}
Let $F_0,\ldots,F_n$ be polynomials in $\mathbb Z[x_0,\ldots,x_n,u]$ such that $F_i-F_j\in(x_i,x_j)$ for every pair $i,j$. There exists $C$ in this polynomial ring with $C-F_i\in(x_i)$ for every $i$. If all $F_i$ are homogeneous of the same degree, $C$ can be chosen homogeneous of that degree.
\end{lemma}
\begin{proof}
For a subset $S\subseteq\{0,\ldots,n\}$ let $\rho_S$ set $x_i=0$ for $i\in S$. For $S\ne\varnothing$, the value $\Phi_S=\rho_S(F_i)$ is independent of the choice of $i\in S$: the difference of two choices belongs to $(x_i,x_j)$ and is killed by $\rho_S$. Set
\[
 C=\sum_{\varnothing\ne S\subseteq\{0,\ldots,n\}}
       (-1)^{|S|+1}\Phi_S.
\]
Modulo $(x_j)$ the terms indexed by $S$ and $S\cup\{j\}$ cancel whenever $j\notin S$ and $S\ne\varnothing$. The remaining term is $\Phi_{\{j\}}$, congruent to $F_j$. All substitutions preserve homogeneity.
\end{proof}

\begin{lemma}[Proved: normalized affine cones]\label{C:affine-cones}
For every $n\geq0$ there are an integer $M\geq n$ and
\[
 t\in\mathbb Z[\operatorname{Hom}_{\Delta_{\mathrm{inj}}}([n],[M])]
\]
of augmentation $1$ such that $th=th'$ whenever $h,h':[m]\to[n]$ are parallel arrows. Here composition is extended bilinearly in the free preadditive category. For $n\geq1$ one can take $M=2^{n+1}-2$.
\end{lemma}
\begin{proof}
We prove a polynomial assertion simultaneously for all integers $k\geq1$. For variables $z_0,\ldots,z_{n+1}$, let $D_i$ identify $z_i$ and $z_{i+1}$ and relabel the remaining variables in order as $y_0,\ldots,y_n$. We claim that there is a homogeneous integral polynomial $C_{n,k}$ with value $1$ when all variables are $1$, such that $y_i^kD_i(C_{n,k})$ is independent of $i$.

For $n=0$, use $C_{0,k}=1$. Assume the assertion for $n-1$ and all $k$. Let $Q=C_{n-1,2k}$, expressed in the variables $y_0,\ldots,y_n$, and prescribe
\[
 D_i(C_{n,k})=F_i:=Q\prod_{l\ne i}y_l^k.
\]
These prescriptions agree on every pairwise intersection of the hyperplanes $z_i=z_{i+1}$. Indeed, for $i<j$, one must check $D_{j-1}(F_i)=D_i(F_j)$. If $j=i+1$, both omitted-variable products become $\prod_lw_l^k$ and the images of $Q$ agree. If $j\geq i+2$, put $p=i$, $q=j-1$. The two expressions are respectively
\[
 w_q^{2k}D_q(Q)\prod_{l\ne p,q}w_l^k,
 \qquad
 w_p^{2k}D_p(Q)\prod_{l\ne p,q}w_l^k,
\]
and are equal by induction. Choose homogeneous polynomial lifts of the $F_i$. The change of coordinates
$x_i=z_i-z_{i+1}$, $u=z_{n+1}$ is invertible over $\mathbb Z$; the kernels of the $D_i$ are the coordinate ideals $(x_i)$. Lemma~\ref{C:interpolation} therefore supplies $C_{n,k}$. It has value $1$ at the all-ones point, and the required common expression is $Q\prod_l y_l^k$. Its degree satisfies
\[
 d(n,k)=d(n-1,2k)+nk=k(2^{n+1}-n-2).
\]

For an injection $g:[n]\to[M]$, let its gaps be
\[
 e_0=g(0),\quad e_i=g(i)-g(i-1)-1\ (1\leq i\leq n),
 \quad e_{n+1}=M-g(n).
\]
This gives a bijection between injections and monomials of degree $M-n$ in $z_0,\ldots,z_{n+1}$. Extend it linearly, writing $P_t$ for the polynomial of an arrow combination $t$. If $\delta^i:[n-1]\to[n]$ omits $i$, deleting $g(i)$ merges two gaps and the deleted point; hence
\[
 P_{t\delta^i}=y_iD_i(P_t).
\]
Take $P_t=C_{n,1}$ and $M=n+d(n,1)$. Its augmentation is $1$, and all $t\delta^i$ agree.

We finally check that this equalizes every parallel pair, rather than just the cofaces. If their common value is $s$, the relation
$\delta^{b+1}\delta^a=\delta^a\delta^b$ for $a\leq b$ gives $s\delta^a=s\delta^b$. Induction on $n$ therefore shows that $s$ equalizes all parallel injections into $[n-1]$. Every proper injection into $[n]$ factors through a coface; factoring two such injections proves $th=th'$. The identity is the only injection $[n]\to[n]$. The case $n=0$ is immediate.
\end{proof}

\begin{theorem}[Refuted: total asphericity need not imply filteredness]\label{C:counterexample}
The category $\Delta_{\mathrm{inj}}$ is totally $\Winfty$-aspherical and has exact colimits of abelian-group-valued diagrams, but is not filtered. Consequently Lead~4 for $\Winfty$ and manuscript row II-C-003 have a negative answer for arbitrary small categories.
\end{theorem}
\begin{proof}
The affine category $\operatorname{Aff}_{\mathbb Z}(\Delta_{\mathrm{inj}})$, defined and analyzed in Theorem~\ref{D:affine-exactness}, has the same objects and all integral arrow combinations of augmentation $1$. Any two objects map to a common larger ordinal. Given two parallel affine arrows $\alpha,\beta:[m]\to[n]$, Lemma~\ref{C:affine-cones} provides $t$ for which $th$ is independent of the ordinary injection $h:[m]\to[n]$. Since the coefficients of both $\alpha$ and $\beta$ sum to $1$, we get $t\alpha=t\beta$. Thus the affine category is filtered. Theorem~\ref{D:affine-exactness} proves exactness of integral colimits. The original category is nonempty and connected, so Proposition~\ref{C:exact-sufficient} proves total $\Winfty$-asphericity.

There are two distinct injections $a,b:[0]\to[1]$, taking $0$ to $0$ and $1$ respectively. Every arrow $c:[1]\to[n]$ is injective, so $ca\ne cb$. Hence this parallel pair has no equalizer by postcomposition, and the category is not filtered. This verifies every hypothesis and the failure in a completely specified countable category; no finite-to-infinite inference is involved.
\end{proof}

\subsection{Monoids and weakly terminal objects}



\begin{lemma}[Proved: directedness and affine equalizers]\label{C:directed-affine}
A totally $\Winfty$-aspherical category $A$ is nonempty, its objects have common successors, and for every parallel pair $a,b:p\to q$ there exist an object $r$ and an integral arrow combination $t\in\mathbb Z[A(q,r)]$ of augmentation $1$ satisfying $ta=tb$.
\end{lemma}
\begin{proof}
Nonemptiness and connectedness follow from Theorem~\ref{C:classification}. We give a direct proof of the common-successor assertion. If two objects $a,b$ have no common successor, their upward closures $U_a,U_b$ are disjoint cosieves. Define $H:A\to\Set$ as the quotient of the constant two-element functor by identifying its two elements at precisely the objects in $U_a\cup U_b$. Thus $\operatorname{El}_A(H)$ is the double of $A$ glued along the disjoint cosieve $U_a\amalg U_b$. Its nerve is the corresponding union of two nerves: a chain meeting a cosieve stays in it, so no mixed chain is created. This is a homotopy pushout of spaces, since the nerve inclusions are monomorphisms. Each $BU_a$ and $BU_b$ is connected, because every object of $U_a$ receives an arrow from $a$, and similarly for $b$. Since $BA$ is contractible, the reduced Mayer--Vietoris sequence gives a nonzero $H_1$ of this double: the kernel of $H_0(BU_a\amalg BU_b;\mathbb Z)\to H_0(BA;\mathbb Z)^{\oplus2}$ contains $(1,-1)$. The double is connected, contradicting Theorem~\ref{C:classification}.

Proposition~\ref{B:lead2} and that theorem show $I=\mathbb Z$. Thus a finite integer combination of the displayed relations has total weight $1$. Its finitely many target objects have a common successor by the first part. Postcompose each relation with a chosen arrow to that successor, multiply it by its integer coefficient, and add. This gives the required single $t$.
\end{proof}

\begin{lemma}[Proved: a signed-measure criterion for monoids]\label{C:monoid-lemma}
Let $M$ be any small monoid, with multiplication $yx$ denoting composition of $x$ followed by $y$. Suppose that for every $x,a\in M$ there is $t\in\mathbb Z[M]$ with augmentation $1$ such that $tx=txa$. Then $M$, viewed as a one-object category, is filtered.
\end{lemma}
\begin{proof}
Fix $x,a\in M$. Right multiplication $r_a:M\to M$, $m\mapsto ma$, is a map of sets. The element $v=tx$ is a finite-support integral signed measure of mass $1$, supported on the left ideal $Mx$, and satisfies $(r_a)_*v=v$.

For any self-map $r:S\to S$, an invariant finite-support signed measure is supported on finite periodic orbits, and is constant on each orbit. To verify this without positivity, take its finite support $F$. Since $r_*v=v$, every element of $F$ has at least one predecessor in $F$. There are $|F|$ possible predecessors and $|F|$ required images, so the restriction $r|_F$ maps $F$ bijectively to itself. Thus $F$ is a union of finite cycles. The invariance equation makes the coefficient constant around each cycle. Consequently the mass is an integer linear combination of the cycle lengths.

There is therefore a point of $Mx$ with a positive period under $r_a$. Choose the least attainable positive period $r$. If $r>1$, choose a prime divisor $\ell$ of $r$ and a point $z\in Mx$ of period $r$. Apply the hypothesis to $z,a$. The resulting mass-one invariant measure is supported in $Mz$; since its mass is an integral combination of its cycle lengths, at least one cycle has length $s$ not divisible by $\ell$. A point $w$ on that cycle belongs to $Mz$, so $wa^r=w$ as well as $wa^s=w$. Its positive period divides $\gcd(r,s)<r$, contradicting minimality. Thus some $y$ satisfies $yxa=yx$.

Given $a,b\in M$, first choose $y$ with $ya=y$ by applying this conclusion with $x=1$. Apply it again with $x=y$ and the element $b$, obtaining $z$ with $zyb=zy$. Then $zya=zy=zyb$. This is the filteredness axiom for the pair $a,b$.
\end{proof}

\begin{theorem}[Proved: monoids and categories with a weakly terminal object]\label{C:weak-terminal}
Let $A$ have a weakly terminal object $q$, meaning that $A(a,q)\ne\varnothing$ for every object $a$. Then $A$ is totally $\Winfty$-aspherical if and only if it is filtered. In particular the equivalence holds for every monoid and every category with finitely many objects, and hence for every category with finitely many arrows.
\end{theorem}
\begin{proof}
The filtered-to-total implication is established in Part~A. For the converse, apply Lemma~\ref{C:directed-affine}. Given $x,a\in M=\operatorname{End}_A(q)$, it supplies an affine combination of arrows $q\to r$ annihilating $x-xa$. Choose an arrow $r\to q$ and postcompose. Lemma~\ref{C:monoid-lemma} then proves that $M$ is filtered.

Let $a,b:p\to r$ be any parallel pair in $A$. Choose $d:r\to q$. There is an affine combination of arrows out of $q$ equalizing $da$ and $db$, and by weak terminality we may postcompose it back to $q$. Write the resulting combination as $t=\sum_{i=1}^N n_ic_i\in\mathbb Z[M]$, with $\sum_i n_i=1$ and $tda=tdb$. In a filtered monoid there exists $z$ with $zc_i=z$ for all $i$: successively equalize the finite family $c_i$ with the identity, preserving previous equalities by further left multiplication. Hence
\[
 zda=\sum_i n_izc_ida=\sum_i n_izc_idb=zdb.
\]
The arrow $zd:r\to q$ equalizes $a,b$. Together with nonemptiness and common successors, this proves that $A$ is filtered.

If $A$ has finitely many objects and is total, repeated use of the common-successor property gives a weakly terminal object. The empty category is neither filtered nor total, so it causes no exception. A one-object category has a weakly terminal object without any finiteness assumption on its monoid.
\end{proof}



\begin{example}[Proved: the counterexample is not sifted]\label{C:not-sifted}
The category $\Delta_{\mathrm{inj}}$ is not sifted. In the category of cospans from $[0],[0]$ to a common ordinal, a cospan is an ordinal with two distinguished points. Their comparison, $<$, $=$, or $>$, is invariant under every morphism of cospans, since all transition functions are injective and increasing. Each comparison class is connected: it receives the canonical cospan on $[1]$, $[0]$, or $[1]$ respectively, by the injection selecting the distinguished points. Thus this cospan category has exactly three connected components, and the diagonal is not final. This prevents an attempted proof of the open implication in Remark~\ref{D:exact-converse-open} by first asserting that totality forces siftedness.
\end{example}

\section{Homology and other fundamental localizers}
\subsection{Homology localizers and their exact action criterion}

\begin{proposition}[Proved: the homology localizer]\label{D:homology-localizer}
For every nonzero commutative ring $R$ with identity, let $W_R$ be the functors whose nerves induce isomorphisms on $H_j(-;R)$ for every $j\geq0$. Then $W_R$ is a fundamental localizer with the precise relative Theorem~A axiom fixed in Part~A.
\end{proposition}
\begin{proof}
Identities, two-out-of-three, and closure under retracts follow from the corresponding properties of isomorphisms of $R$-modules. A category with a terminal object has a contractible nerve and therefore maps to the point by a member of $W_R$.

For the relative axiom, consider functors $u:A\to B$, $v:B\to C$, and $f=vu:A\to C$, and assume that $A/c\to B/c$ belongs to $W_R$ for every $c\in C$. Use the bisimplicial set
\[
 S_A(p,q)=\{a_0\to\cdots\to a_q,\quad
                  f(a_q)\to c_0\to\cdots\to c_p\}.
\]
The last $A$-face composes the last arrow of the $A$-chain with the displayed arrow to $c_0$; the first $C$-face composes that displayed arrow with $c_0\to c_1$. The other faces and degeneracies have the evident nerve formulas. These formulas commute and define the indicated bisimplicial set. Define $S_B$ similarly using $v$. The functor $u$ induces a bisimplicial map $S_A\to S_B$ by applying $u$ to the $A$-chain and leaving the $C$-chain unchanged.

For a fixed $A$-chain, the $p$-direction is the nerve of $f(a_q)\backslash C$, which has initial object $\id_{f(a_q)}$. Thus projection onto the $A$-chain induces a homology isomorphism from the total normalized chains of $S_A$ to $C_*(NA;R)$. This follows by taking homology in the $p$-direction: only its degree-zero row remains. The same argument applies to $S_B$ and $B$.

For a fixed $C$-chain, the $q$-direction is the nerve of $A/c_0$. Indeed, its last arrow to $c_0$ determines the arrows from all preceding objects by composition. The analogous direction of $S_B$ is $B/c_0$. The hypothesis therefore makes $S_A\to S_B$ an isomorphism after taking homology in the $q$-direction, hence on the homology of the total complexes. Both spectral sequences are first-quadrant and converge degreewise. The two projections commute with the specified map $Nu$, so $Bu$ induces the resulting homology isomorphism. Hence $u\in W_R$.
\end{proof}

\begin{theorem}[Proved: total $R$-asphericity]\label{D:R-action}
A small category $A$ is totally $W_R$-aspherical if and only if it is nonempty and connected and satisfies $\mathcal P_R(A)$ of Definition~\ref{C:action-definition}. In particular
\[
 \{A:A\text{ is totally }W_{\mathbb Z}\text{-aspherical}\}
 =\{A:A\text{ is totally }\Winfty\text{-aspherical}\}.
\]
\end{theorem}
\begin{proof}
If $A$ is total, its identity functor first gives $R$-acyclicity. Since $R\ne0$, the induced map $R[\pi_0(A)]\to R$ is an isomorphism only when $A$ is nonempty and connected. Apply the collage construction to each component subfunctor of an arbitrary $H:A\to\Set$. The reduced Mayer--Vietoris sequence gives acyclicity of that component, exactly as in Theorem~\ref{C:classification}, with $R$ replacing $\mathbb Z$.

Conversely, for a final functor $f:A\to B$, every $b\backslash A$ is a connected category of elements and is therefore $R$-acyclic. Lemma~\ref{C:homology-A} makes $Bf$ an $R$-homology isomorphism. The constant singleton functor gives $R$-acyclicity of $A$, so $B$ is $R$-acyclic. The comprehensive-factorization characterization proves totality. The last equality was proved in Theorem~\ref{C:classification}; the use of integral coefficients there is essential in detecting nontrivial cyclic groups.
\end{proof}

\subsection{The exact-colimit condition, including its necessary quantifiers}

\begin{definition}[Proved: affine linearization]\label{D:aff-definition}
For a small category $A$ and a nonzero commutative ring $R$, define $\operatorname{Aff}_R(A)$ to have the objects of $A$ and the morphisms
\[
 \operatorname{Aff}_R(A)(a,b)=
 \left\{\sum_{f\in A(a,b)}r_ff\in R[A(a,b)]:
       \text{finite support},\ \sum_fr_f=1_R\right\}.
\]
Composition is the bilinear extension of composition in $A$; augmentations multiply, so composition preserves augmentation $1$. Each identity is the original identity with coefficient $1_R$. This is an ordinary small category, generally not an additive category. Its connected components are those of $A$.
\end{definition}

\begin{theorem}[Proved: the affine exact-colimit criterion]\label{D:affine-exactness}
For every small category $A$, including the empty category, and every nonzero commutative ring $R$, the following conditions are equivalent:
\begin{enumerate}
\item $\colim_A:(R\text{-}\mathrm{Mod})^A\to R\text{-}\mathrm{Mod}$ is exact.
\item Every connected component of $\operatorname{Aff}_R(A)$ is filtered.
\end{enumerate}
In explicit terms, condition (2) requires common successors within each component and an augmentation-one affine equalizer for every pair of parallel \emph{affine} arrows. It is not merely the condition in Lead~2 for one pair of original arrows. For $R=\mathbb Z$ this is the Isbell--Mitchell affine criterion \citep{IsbellMitchell1973Exact,IsbellMitchell1976Fixed}; a related finite-family formulation is verified in \citet[Section~1, unnumbered theorem]{Isbell1968Exact}. We supply the proof rather than relying on an unavailable theorem number from the original AMS papers.
\end{theorem}
\begin{proof}
Every functor $F:A\to R\text{-}\mathrm{Mod}$ extends to a functor $\bar F:\operatorname{Aff}_R(A)\to R\text{-}\mathrm{Mod}$ by
$\bar F(\sum r_ff)=\sum r_fF(f)$. This is functorial because $R$ is commutative and the maps are $R$-linear. Natural transformations extend in the same way, and pointwise exact sequences remain pointwise exact. A cocone $F\to\underline V$ respects every original arrow if and only if it respects every affine arrow: the latter follows from $\sum r_f=1_R$. Hence the identity on the individual values induces a canonical natural isomorphism
\[
 \colim_A F\cong\colim_{\operatorname{Aff}_R(A)}\bar F.
\]
Filtered colimits of modules and direct sums of modules are exact. Condition (2), followed by taking the direct sum over components, therefore implies (1).

Assume (1). We may work in one connected component of $A$, since extension by zero shows that its colimit functor is exact. For an object $a$, let $U_a$ be the set of objects receiving an arrow from $a$, and let $R_a\subseteq\underline R$ be the subfunctor taking value $R$ on $U_a$ and zero elsewhere. It is a subfunctor because $U_a$ is a cosieve. Its colimit is $R$, and the inclusion induces the identity map to $\colim_A\underline R=R$: the full subcategory $U_a$ is nonempty and connected, and all its coefficient maps are identities. Exact additive functors preserve finite limits, hence intersections of subobjects. Therefore
\[
 \colim_A(R_a\cap R_b)
 \cong(\colim_A R_a)\cap(\colim_A R_b)=R.
\]
Since $R\ne0$, the intersection functor cannot be zero, so $a,b$ have a common successor. Iteration supplies a common successor for any finite family of objects in this component.

Let $\alpha,\beta:p\to q$ be parallel affine arrows and write $\delta=\alpha-\beta\in R[A(p,q)]$, whose augmentation is zero. Let $P=R[A(p,-)]$ and let $I\subseteq P$ be the subfunctor generated, as a module-valued functor, by $\delta\in P(q)$. Thus $I(r)$ consists of finite linear combinations of $c\delta$ with $c:q\to r$. Exactness makes $\colim I\to\colim P=R$ injective. Its image is zero, because $\delta$ has augmentation zero and generates $I$. Hence $\colim I=0$.

We use the explicit direct-sum presentation of a colimit. The element $\delta$ in the summand $I(q)$ is consequently a finite sum
\[
 \delta=\sum_{j=1}^N
   \bigl(I(v_j)x_j-x_j\bigr)
 \quad\text{in }\bigoplus_{r\in A}I(r),
\]
where $v_j:s_j\to t_j$, $x_j\in I(s_j)$, and signs and scalar coefficients have been absorbed into the $x_j$. Choose a common successor $s$ of $q$ and all the $s_j,t_j$, and arrows $d_r:r\to s$ for these finitely many objects. Applying the direct-sum homomorphism whose $r$-component is $I(d_r)$ gives
\[
 d_q\delta=\sum_j(d_{t_j}v_j-d_{s_j})x_j.
\]
Write $x_j=\sum_l r_{jl}c_{jl}\delta$ with $c_{jl}:q\to s_j$. Then
\[
 t=d_q-\sum_{j,l}r_{jl}(d_{t_j}v_j-d_{s_j})c_{jl}
       \in R[A(q,s)]
\]
has augmentation $1_R$ and satisfies $t\delta=0$. Thus $t\alpha=t\beta$ in $\operatorname{Aff}_R(A)$. Combined with common successors, this is filteredness. If $A$ is empty, both conditions are vacuously satisfied: its colimit functor has the zero module as its only value.
\end{proof}

\begin{corollary}[Proved: exact module colimits imply total homology asphericity]\label{D:exact-R-sufficient}
If $A$ is nonempty and connected and $\colim_A:(R\text{-}\mathrm{Mod})^A\to R\text{-}\mathrm{Mod}$ is exact, then $A$ is totally $W_R$-aspherical. For $R=\mathbb Z$ the conclusion strengthens to total $\Winfty$-asphericity.
\end{corollary}
\begin{proof}
The direct-sum functor $\Sigma_H$ in Proposition~\ref{C:exact-sufficient} is exact for modules over $R$ as well. It makes the colimit functor on each component of each $\operatorname{El}_A(H)$ exact. Applying it to the bar resolution of the constant $R$-module proves $\mathcal P_R(A)$. Theorem~\ref{D:R-action} applies. The integral strengthening is Theorem~\ref{C:classification}.
\end{proof}

\begin{example}[Refuted: the literal homology half of Lead 4]\label{D:disconnected-counterexample}
Let $D$ have two objects $0,1$ and only their identity arrows. Its colimit functor on abelian groups is $(M_0,M_1)\mapsto M_0\oplus M_1$, which is exact. Nevertheless $D$ is not totally $W_{\mathbb Z}$-aspherical: the identity is final and $H_0(BD;\mathbb Z)=\mathbb Z^2\ne\mathbb Z$. The empty category provides another counterexample, with exact zero colimit. Thus the literal equivalence stated in Lead~4 is false; nonemptiness and connectedness cannot be suppressed. The finite category and its composition law are checked in \texttt{cd\_finite\_checks.py}.
\end{example}

\begin{remark}[Open: the corrected connected converse]\label{D:exact-converse-open}
The remaining implication is the precise statement
\[
 \left.
 \begin{array}{c}
 A\text{ small, nonempty, connected},\\
 A\text{ totally }\Winfty\text{-aspherical}
 \end{array}\right\}
 \quad\Longrightarrow\quad
 \colim_A:\Ab^A\to\Ab\text{ exact}.
\]
Equivalently, by Theorems~\ref{C:classification} and \ref{D:affine-exactness}, does $\mathcal P_\infty(A)$ force $\operatorname{Aff}_{\mathbb Z}(A)$ to be filtered? We have not proved this implication or produced a counterexample. This is an unresolved implication of the present work, not an assertion of a currently open problem in the literature. The attempted reduction through the acyclicity of $\mathbb Z[H]$ fails because these diagrams are not generally projective and their acyclicity does not supply acyclicity of arbitrary module diagrams. Lead~2 supplies affine equalizers of original parallel arrows, but exactness requires equalizers of all affine pairs. The proof in Lemma~\ref{C:monoid-lemma} bridges this gap for monoids and, by Theorem~\ref{C:weak-terminal}, for categories with a weakly terminal object. It also holds for $\Delta_{\mathrm{inj}}$ by its explicit affine cones. The first size not covered by the finite-object theorem is the case of countably infinitely many objects with no weakly terminal object; no particular smallest counterexample is claimed.

For general nonzero $R$, the corresponding unproved necessity is $\mathcal P_R(A)$, with $A$ connected and nonempty, implying filteredness of $\operatorname{Aff}_R(A)$. Neither assertion is replaced here by its sufficient direction.
\end{remark}

\subsection{Groups separate integral and other homology localizers}

\begin{proposition}[Proved: a subgroup criterion for groups]\label{D:group-criterion}
Let $G$ be a small group, regarded as a one-object category. It is totally $W_R$-aspherical if and only if every subgroup $K\leq G$ is $R$-acyclic. If $G$ is finite, this is equivalent to $|G|\in R^\times$. In particular a nontrivial finite group is totally $W_{\mathbb Q}$-aspherical but is not totally $\Winfty$-aspherical.
\end{proposition}
\begin{proof}
A covariant set-valued functor on $G$ is a left $G$-set. Its connected categories of elements are its transitive action groupoids. Choosing one object in a transitive orbit identifies its automorphism group with the corresponding stabilizer $K$, and its inclusion is an equivalence of categories. Every subgroup is realized by the coset action $G/K$. The first equivalence is therefore exactly Theorem~\ref{D:R-action}.

If $G$ is finite and $|G|$ is invertible, the element $e=|G|^{-1}\sum_{g\in G}g\in R[G]$ splits the augmentation as a homomorphism of right $R[G]$-modules: $eg=e$ and the augmentation of $e$ is $1$. Thus the trivial module $R$ is a direct summand of the free module $R[G]$. Tensoring with it, which computes coinvariants and hence $\colim_G$, is exact. Corollary~\ref{D:exact-R-sufficient} proves totality.

Conversely, if a prime $p$ divides $|G|$, Cauchy's theorem supplies a subgroup $C_p$. The standard periodic resolution of $C_p$ gives
$H_1(BC_p;R)=R/pR$. Totality forces this group to vanish, so $p$ is a unit in $R$. Every prime factor of $|G|$ is therefore a unit, and so is $|G|$. A nontrivial group has nontrivial fundamental group in its classifying space, which excludes total $\Winfty$-asphericity already for the identity final image. The order-two multiplication table and its periodic integral and rational chain complexes are checked in \texttt{cd\_finite\_checks.py}.
\end{proof}

\subsection{Truncated homotopy localizers}

\begin{definition}[Proved: truncation convention]\label{D:n-convention}
For an integer $n\geq0$, $W_n$ denotes functors inducing bijections on connected components and isomorphisms on all based homotopy groups in degrees $1,\ldots,n$, at every basepoint. Thus $W_0$ tests only components. A category is $W_n$-aspherical precisely when its classifying space is nonempty and $n$-connected. These are fundamental localizers: weak saturation and the terminal-object axiom are immediate, and the relative axiom follows by applying $n$-truncation to the two bisimplicial comparisons used in Proposition~\ref{D:homology-localizer}. The functor to $n$-types is left adjoint to inclusion \citep[Proposition~5.5.6.18]{Lurie2009HTT} and consequently preserves homotopy colimits in the category of $n$-types, so an objectwise $W_n$-equivalence induces a $W_n$-equivalence on those comparisons.
\end{definition}

\begin{lemma}[Proved: final functors and fundamental groups]\label{D:final-pi1}
An ordinary final functor $f:A\to B$ induces a bijection on components and, on each component with a chosen object $a_0$, an epimorphism
\[
 \pi_1(BA,a_0)\longrightarrow\pi_1(BB,f(a_0)).
\]
\end{lemma}
\begin{proof}
The assertion about components follows by applying finality to the colimit of the constant singleton functor. For each object $b$ in the chosen component, choose an object $(a_b,\beta_b:b\to f(a_b))$ of $b\backslash A$. At $b=f(a_0)$ choose $(a_0,\id)$. For an arrow $h:b\to c$, the two objects $(a_b,\beta_b)$ and $(a_c,\beta_ch)$ of $b\backslash A$ are joined by a zigzag because that comma category is connected. Its projection is a path from $a_b$ to $a_c$ in the groupoid completion of $A$, and its image under $f$ is the path obtained from $h$ by conjugating its endpoints with $\beta_b$ and $\beta_c$. Concatenating these paths along any based edge loop of $NB$ gives a based edge loop of $NA$ with the same image in $\pi_1(BB)$. The nerve's triangular relations are respected because each comma-category equality is an equality of composites in $B$. This proves surjectivity.
\end{proof}

\begin{theorem}[Proved: the cases $W_0$ and $W_1$]\label{D:zero-one}
A small category is totally $W_0$-aspherical if and only if it is nonempty and connected. It is totally $W_1$-aspherical if and only if its classifying space is nonempty and simply connected.
\end{theorem}
\begin{proof}
The identity final functor proves necessity in both cases. Lemma~\ref{D:final-pi1} proves sufficiency: a bijection on components preserves the unique component, and an epimorphic image of the trivial fundamental group is trivial.
\end{proof}

\begin{theorem}[Proved: the cases $W_n$, $n\geq2$]\label{D:truncated-classification}
For $n\geq2$, a small category $A$ is totally $W_n$-aspherical if and only if $BA$ is $n$-connected and every connected component of every $\operatorname{El}_A(H)$, for covariant $H:A\to\Set$, is $(n-1)$-connected.
\end{theorem}
\begin{proof}
Assume totality. The identity final image makes $BA$ $n$-connected. For a connected component of an element category, the collage gives, by reduced Mayer--Vietoris and $n$-connectedness of both $BA$ and its final image, vanishing of integral reduced homology in degrees $0,\ldots,n-1$. This truncated homological condition is inherited by categories of elements, by the explicit construction in Lemma~\ref{C:inherit-elements}. In particular it holds for every connected cyclic-subgroup covering of any such component. Since $n-1\geq1$, the argument in Lemma~\ref{C:cyclic-covers} proves that the component is simply connected. The Hurewicz theorem now makes it $(n-1)$-connected.

Conversely, if $f:A\to B$ is final, its comma categories are $(n-1)$-connected by the stated action condition. Lemma~\ref{C:homology-A}, with $d=n$, makes $H_j(Bf;\mathbb Z)$ an isomorphism for $j<n$ and an epimorphism for $j=n$. Thus $BB$ has zero reduced homology through degree $n$. It is connected and simply connected by Lemma~\ref{D:final-pi1}. Applying Hurewicz in the first possible nonzero degree shows that $BB$ is $n$-connected. This is the required totality.
\end{proof}

\begin{corollary}[Proved: truncated totality with a weakly terminal object]\label{D:truncated-finite}
For every $n\geq2$, a small category with a weakly terminal object is totally $W_n$-aspherical if and only if it is filtered. The same equivalence holds for every category with finitely many objects, including the empty category by separate inspection. It fails for $n=1$: the contractible nonfiltered right-zero monoid of Part~B is totally $W_1$-aspherical.
\end{corollary}
\begin{proof}
For $n\geq2$, Theorem~\ref{D:truncated-classification} makes every connected category of elements simply connected, so its first integral homology is zero. The proof of common successors in Lemma~\ref{C:directed-affine} needs only connectedness and vanishing first homology of $BA$ and of the doubled category of elements, not full contractibility. Proposition~\ref{B:lead2} therefore gives the same affine equalizers as before. The entire proof of Theorem~\ref{C:weak-terminal} now applies unchanged to these specified inputs. Filteredness implies total $\Winfty$-asphericity and therefore total $W_n$-asphericity. The final assertion follows from Theorem~\ref{D:zero-one} and the explicit example in Part~B.
\end{proof}

\begin{proposition}[Proved: dependence on the localizer and on the category]\label{D:general-localizer}
If $W\subseteq V$ are fundamental localizers, every totally $W$-aspherical category is totally $V$-aspherical. In particular every totally $\Winfty$-aspherical category is totally $W$-aspherical for every fundamental localizer satisfying exactly the axioms fixed in Part~A. Total $\Winfty$-asphericity is not invariant under a weak homotopy equivalence of categories.
\end{proposition}
\begin{proof}
Every tested category of elements whose map to the point lies in $W$ also has that map in $V$. The special case follows from the minimality theorem $\Winfty\subseteq W$ cited with its exact hypotheses in Part~A. For the last assertion, the nonfiltered three-element right-zero monoid of Part~B has contractible classifying space but is not total, whereas the point is filtered and total. The specified functor from that monoid to the point is a weak homotopy equivalence. This also places the same obstruction recorded by the companion at XII ms.~120: totality depends on ordinary final images and on set-valued actions, not only on the classifying-space homotopy type.
\end{proof}

\begin{remark}[Proved: the zero coefficient ring]\label{D:zero-ring}
If the zero ring is allowed, every homology group with its coefficients is zero, so $W_R$ is the class of all functors and every small category, including the empty category, is totally $W_R$-aspherical. Its module category is the zero category, and every colimit functor is exact. The nonzero hypothesis in the preceding affine and connectedness criteria is therefore necessary. The ring-theoretic statements above have been stated for commutative unital rings; no unannounced extension of their affine-linearization proof to noncommutative coefficients is being used.
\end{remark}

\subsection{Fixed points: an additional necessary condition and a literature correction}

\begin{proposition}[Proved: preservation of monomorphisms of set diagrams]\label{D:set-monos}
If $A$ is totally $\Winfty$-aspherical, then
$\colim_A:\Set^A\to\Set$ preserves monomorphisms. Moreover every span
$a\xrightarrow{f}b$, $a\xrightarrow{g}c$ in $A$ has a commuting completion:
there are $u:b\to d$, $v:c\to d$ with $uf=vg$.
\end{proposition}
\begin{proof}
For a monomorphism $K\hookrightarrow H$, identify $K$ with a subfunctor.
If two objects of $\operatorname{El}_A(K)$ have the same image in
$\colim_AH$, they lie in one connected component $E$ of
$\operatorname{El}_A(H)$. Theorem~\ref{C:classification} makes $E$ totally
aspherical, and Lemma~\ref{C:directed-affine} gives a common successor of
these two objects in $E$. Since $\operatorname{El}_A(K)$ is a cosieve, this
successor and the two arrows belong to $\operatorname{El}_A(K)$.
Consequently the two objects also have the same image in $\colim_AK$.
This proves injectivity.

Apply the same common-successor assertion to
$\operatorname{El}_A(A(a,-))=a\backslash A$, which is connected because it
has an initial object. Its two objects $f$ and $g$ have a common successor
precisely when the displayed commuting completion exists.
\end{proof}

\begin{proposition}[Proved: fixed points in connected representations]\label{D:fixed-points}
Let $A$ be totally $\Winfty$-aspherical, and let $H:A\to\Set$ have a
nonempty connected category of elements. Every natural endomorphism
$\tau:H\to H$ fixes an element: there are $a\in A$ and $h\in H(a)$ with
$\tau_a(h)=h$.
\end{proposition}
\begin{proof}
We use Corollary~\ref{B:homotopical-colimits}: ordinary $A$-colimits of
simplicial sets preserve objectwise weak equivalences.

First suppose that no element of any $H(a)$ is periodic under $\tau_a$.
Let $X(a)$ be the nerve of the category of the $\mathbb N$-action on
$H(a)$ generated by $\tau_a$. An arrow $x\to y$ is an integer $m\geq0$
with $\tau_a^m(x)=y$; composition adds integers. In the absence of periodic
elements there is at most one arrow with given endpoints. Each connected
component is directed: elements joined by a finite zigzag of forward
iterates have a common forward iterate. Thus each component is filtered
and weakly contractible. The quotient map
\[
 X(a)\longrightarrow Q(a):=H(a)/\langle x\sim\tau_a(x)\rangle
\]
is an objectwise weak equivalence to a discrete simplicial set.

In simplicial degree $j$ one has $X(a)_j=H(a)\times\mathbb N^j$.
Product with the fixed set $\mathbb N^j$ preserves set colimits.
As $\colim_AH$ is a singleton, the face and degeneracy maps consequently
identify $\colim_AX$ with the nerve $B\mathbb N$ of the one-object monoid
$\mathbb N$. On the other hand, $\colim_AQ$ is a singleton. Homotopicality
of colimit would make $B\mathbb N$ weakly contractible, whereas its
edge-path presentation has fundamental group the group completion
$\mathbb Z$. This contradiction proves existence of a periodic element.

Let $r\geq1$ be the least positive period occurring anywhere, and put
\[
 P(a)=\{h\in H(a):\tau_a^r(h)=h\}.
\]
Every element of $P(a)$ has exact period $r$. Naturality makes $P$ a
subfunctor, and its category of elements is connected. Indeed, any two
of its objects have a common successor in the totally aspherical
category $\operatorname{El}_A(H)$; the successor belongs to $P$.
The restriction of $\tau$ therefore gives a free action of the cyclic
group $C_r$ on each $P(a)$.

The nerve $Y(a)$ of the action groupoid $P(a)//C_r$ maps by an objectwise
weak equivalence to the orbit set $P(a)/C_r$: each orbit groupoid has
trivial stabilizers and is equivalent to a point. Its degree-$j$ set is
$P(a)\times C_r^j$. Since $\colim_AP$ is a singleton, the colimits of
these two diagrams are $BC_r$ and a point. Homotopicality now makes
$BC_r$ weakly contractible, so its fundamental group $C_r$ is trivial.
Thus $r=1$, providing a fixed element.
\end{proof}

\begin{remark}[Known: the fixed-point condition does not imply exactness]\label{D:fpp-not-exact}
The fixed-point condition of Proposition~\ref{D:fixed-points}, for all
connected representations, is strictly weaker than exactness of
abelian colimits. A small connected counterexample is constructed in
\citet[Part~C, Construction~C.1 and Theorem~C.2, pp.~384--389]{AdamekKoubekReiterman1981Embeddings}.
Their Section~A.4, p.~370, explicitly corrects the earlier identification
of these conditions. In particular, the statement in
\citet[Section~II.3, p.~100]{AdamekReiterman1976SetColim} cannot be used as
an exactness criterion. The 1981 counterexample has the fixed-point
property but fails simultaneous affine equalization of three specified
parallel arrows.

This source correction does not resolve Remark~\ref{D:exact-converse-open}:
we have not proved that the published counterexample satisfies
$\mathcal P_\infty$. The distinction between pairwise and simultaneous
affine equalization remains essential.
\end{remark}

\subsection{The branch in which the two-sphere is aspherical}

\begin{proposition}[Proved: totality above the one-type localizer]
\label{D:sphere-totality}
Let $W$ be a fundamental localizer with $W_1\subseteq W$.
For every small category $A$, including the empty category,
\[
 T_W(A)\quad\Longleftrightarrow\quad A\to e\in W.
\]
Equivalently, every ordinary final functor whose source is
$W$-aspherical has a $W$-aspherical target. By
Proposition~\ref{F:sphere-localizer}, the assumption on $W$
is precisely that $S^2$ be $W$-aspherical; it imposes no
accessibility condition.
\end{proposition}

\begin{proof}
Necessity follows from the identity final functor. For
sufficiency, suppose that $A\to e\in W$, and let
$f:A\to B$ be ordinary final. If $A$ is empty, finality
forces $B$ to be empty, and there is nothing to prove.
Assume therefore that $A$ is nonempty.

Write $\Pi A$ and $\Pi B$ for groupoid completions.
The canonical functors $A\to\Pi A$ and $B\to\Pi B$
belong to $W_1$: the edge-path presentation identifies
their components and fundamental groups. Finality gives
a bijection on components and surjections on the
corresponding fundamental groups by
Lemma~\ref{D:final-pi1}. Choose one object $a_i$ in each
component of $A$, indexed by the small set $I$.
Let $\mathcal S\subset\Pi A$ be the full subgroupoid on
the $a_i$, and let $\mathcal T\subset\Pi B$ be the full
subgroupoid on the objects $f(a_i)$. Their inclusions are
equivalences of categories, and distinct indices give
distinct objects of $\mathcal T$ because the component
map is bijective. The restriction of $\Pi f$ is thus a
functor
\[
 \psi:\mathcal S\longrightarrow\mathcal T
\]
which is bijective on objects and surjective on every
endomorphism group. There are no arrows between distinct
objects in either of these skeletal groupoids. All the
specified comparisons to $A$ and $B$ belong to $W$, and
two-out-of-three gives $\mathcal S\to e\in W$.

Theorem~\ref{F:sphere-cofibrations} makes the
object-injective functor $\psi$ a $W$-cofibration.
Its ordinary pushout along $\mathcal S\to e$ is
\[
 \begin{tikzcd}
 \mathcal S\arrow[r,"\psi"]\arrow[d]&
 \mathcal T\arrow[d]\\
 e\arrow[r,equal]&e.
 \end{tikzcd}
\]
To verify this pushout in $\Cat$, a compatible pair of
functors from $\mathcal T$ and $e$ to any category $Z$
must send every object of $\mathcal T$ to the selected
object of $Z$. Every arrow of $\mathcal T$ lifts along
$\psi$, so it must be sent to the identity. The pair is
therefore uniquely determined by that object of $Z$,
which is exactly the stated universal property.
Apply the pushout functor of $\psi$ to the morphism
$\mathcal S\to e$ from the identity object of
$\mathcal S\backslash\Cat$. Its image is
$\mathcal T\to e$, which consequently belongs to $W$.
The equivalence $\mathcal T\to\Pi B$ and the canonical
$W_1$-equivalence $B\to\Pi B$ now give $B\to e\in W$.
The final-image characterization proves $T_W(A)$.
\end{proof}

\begin{corollary}[Proved: presheaf and element-category totality
above the one-type localizer]\label{D:sphere-presheaves}
Under the same hypothesis on $W$, for every small $U$ and
every presheaf $H:U^{\op}\to\Set$, the following are
equivalent:
\[
 H\text{ is totally }W\text{-aspherical},\qquad
 U/H\text{ is }W\text{-aspherical},\qquad
 T_W(U/H).
\]
These equivalences include the empty presheaf. Strictly
and totally $W$-aspherical functors coincide, and the
$W$-dwy\'erienne and totally $W$-dwy\'erienne sieve
classes coincide.
\end{corollary}

\begin{proof}
The identity Kan-extension test gives the first
implication. Proposition~\ref{D:sphere-totality} gives
the second, and Proposition~\ref{E:forward-subterminal}
gives the third back to the first. Equivalently, for
every $p:U\to C$, the canonical functor
$U/H\to C/p_!H$ is final, and the preceding proposition
makes its target $W$-aspherical whenever its source is.
If $H$ is empty, all these element categories are empty;
the three conditions each say exactly that
$\varnothing\to e$ belongs to $W$. Applying the
equivalences to every restricted representable gives
the functor assertion and, on the nonempty boundary
weights, the sieve assertion.
\end{proof}

\begin{theorem}[Proved: sieve composition above the one-type
localizer]\label{D:sphere-sieve-composition}
Fix a fundamental localizer $W$ with $W_1\subseteq W$.
The following conditions are equivalent:
\begin{enumerate}
\item The $W$-dwy\'erienne sieves are closed under
composition.
\item The totally $W$-dwy\'erienne sieves are closed
under composition.
\item Every nonempty small category is $W$-aspherical.
\end{enumerate}
When these conditions hold, every full sieve belongs
to both classes. If condition~3 fails, the failure of
composition can be witnessed by a classical Dwyer
first factor and a second factor obtained by adjoining
one terminal object, as in
Theorem~\ref{F:total-sieve-cosieves}.
\end{theorem}

\begin{proof}
The first two conditions are equivalent by
Corollary~\ref{D:sphere-presheaves}. If condition~3
holds, every nonempty boundary category is
$W$-aspherical and hence total by
Proposition~\ref{D:sphere-totality}. Thus every full
sieve belongs to both classes, proving their
composition closure.

Conversely assume condition~2, and take an arbitrary
nonempty small category $V$. Its initial cone
$A=V^{\triangleleft}$ is weakly contractible, hence
$W$-aspherical, and consequently total by
Proposition~\ref{D:sphere-totality}. The inclusion
$V\subset A$ is a nonempty full cosieve.
The exact criterion in
Theorem~\ref{F:total-sieve-cosieves} therefore makes
$V$ totally $W$-aspherical and, by the identity test,
$W$-aspherical. This proves condition~3. If it fails,
choose such a $V$ which is not $W$-aspherical and
apply the lower-copy construction in that theorem
to $V\subset V^{\triangleleft}$. Its first factor
is classical Dwyer, its intermediate category is
total as a final image of $V^{\triangleleft}$, and
its second factor adjoins a terminal object. The
composite's boundary category there is $V$, so the
composite belongs to neither of the two coincident
classes. This verifies the stated form of failure.
\end{proof}

\section{Presheaves and their categories of elements}
\subsection{The implication that always holds}

\begin{proposition}[Proved: forward implication and the subterminal converse]\label{E:forward-subterminal}
For any fundamental localizer $W$ and any presheaf $H:U^{\op}\to\Set$, total $W$-asphericity of $U/H$ implies total $W$-asphericity of $H$. If $H$ is subterminal, the converse also holds.
\end{proposition}
\begin{proof}
The forward implication is Proposition~\ref{A:elements-Kan}. For the converse under the stated extra hypothesis, identify $E=U/H$ with the full sieve of $U$ on which $H$ is nonempty. Given any $f:E\to V$, form the ordinary pushout
\[
 \begin{tikzcd}
 E\arrow[r,hook,"\pi"]\arrow[d,"f"']&U\arrow[d,"p"]\\
 V\arrow[r,hook,"j"']&U'.
 \end{tikzcd}
\]
The pushout of a sieve along an arbitrary functor is again a full sieve: its objects are those of $V$ and $U\setminus E$, no new arrow has target in $V$ and source outside it, and arrows between objects of $V$ remain their original arrows. This follows as well from the usual coend formula for the crossing arrows. In particular $j$ is a discrete fibration.

Factor $f$ as $E\xrightarrow{q_f}D_f\xrightarrow{\rho_f}V$ by Proposition~\ref{A:comprehensive}. Since $j\rho_f$ is a discrete fibration, this same factorization is the comprehensive factorization of $jf=p\pi$. Hence
\[
 D_f\cong U'/(p\pi)_!1_E\cong U'/p_!H
\]
as categories over $U'$. Totality of $H$ makes this category $W$-aspherical. Since $f$ was arbitrary, $T_W(E)$ follows. The empty sieve is included in this proof; whether its totality holds depends on whether $\varnothing\to e$ belongs to $W$.
\end{proof}

\subsection{A finite counterexample to the converse}

\begin{lemma}[Proved: fixed-arrow presheaves of an idempotent]\label{E:idempotent}
Let $C$ be small, let $c\in C$, and let $d:c\to c$ satisfy $d^2=d$. Define
\[
 Q_d(v)=\{\alpha:v\to c\mid d\alpha=\alpha\}.
\]
This is a split subpresheaf of $C(-,c)$, with retraction $\alpha\mapsto d\alpha$. Its category of elements $C/Q_d$ is filtered and has contractible classifying space.
\end{lemma}
\begin{proof}
The condition is preserved by precomposition, and the stated retraction is natural and restricts to the identity on $Q_d$. The category of elements is nonempty because $(c,d)$ is an object. Each $(v,\alpha)$ maps to $(c,d)$ by the arrow $\alpha:v\to c$. If $r,s:(v,\alpha)\rightrightarrows(w,\beta)$ are parallel, then $\beta r=\alpha=\beta s$, so the arrow $\beta:(w,\beta)\to(c,d)$ equalizes them. These are exactly the two filteredness conditions. There is also an explicit natural transformation from the identity functor of $C/Q_d$ to the constant functor at $(c,d)$, with component $\alpha$ at $(v,\alpha)$; its naturality is the defining equation for an arrow in the element category.
\end{proof}

\begin{theorem}[Refuted: row B035 and the first biconditional in Lead 5]\label{E:counterexample}
Let $U$ be the one-object category with endomorphism monoid $\{1,e\}$, $e^2=e$. Let $R=U(-,*)$, let $Q=\{e\}\subset R$, and form the pushout of presheaves
\[
 H=R\amalg_QR.
\]
Explicitly $H(*)=\{i,h_1,h_2\}$, and $H(e)$ is the constant function with value $i$. Then $H$ is totally $\Winfty$-aspherical, and consequently totally $W$-aspherical for every fundamental localizer $W$. The category $E=U/H$ is not totally $\Winfty$-aspherical, is not totally $W_R$-aspherical for any nonzero commutative ring $R$, and is not filtered. Thus the converse of manuscript Proposition~8 is false, even for a finite category with two arrows and a three-element presheaf.
\end{theorem}
\begin{proof}
First $Q\hookrightarrow R$ is split by the natural map sending both elements of $R(*)$ to $e$. Let $p:U\to C$ be completely arbitrary. It specifies $c=p(*)$ and $d=p(e)$ with $d^2=d$. Since left Kan extension preserves representables, coproducts, pushouts, and a specified splitting, there are canonical isomorphisms
\[
 p_!R\cong C(-,c),\qquad
 p_!Q\cong Q_d,\qquad
 p_!H\cong C(-,c)\amalg_{Q_d}C(-,c).
\]
For the second isomorphism, applying $p_!$ to the split pair gives the splitting of the idempotent induced by $d$ on $C(-,c)$; its image is exactly $Q_d$. This argument uses preservation of a split pair, not preservation of arbitrary monomorphisms by left Kan extension.

The two copies of $C(-,c)$ embed as subpresheaves of $p_!H$. Their element categories are full sieves isomorphic to $C/c$, and their intersection is $C/Q_d$. They cover $C/p_!H$. A chain in that union belongs to one of the two sieves: its last object lies in one of them and all predecessors then lie there. Therefore their nerves form an actual pushout along monomorphisms, and their realizations form a homotopy pushout. Both copies of $C/c$ have terminal object $\id_c$; the intersection is contractible by Lemma~\ref{E:idempotent}. The homotopy pushout is contractible. Since this holds for every $C$ and $p$, $H$ is totally $\Winfty$-aspherical. Cisinski's minimality theorem gives the assertion for every fundamental localizer.

The category $E$ has exactly three objects $i,h_1,h_2$. It has identities at all three objects, one additional idempotent $e_i:i\to i$, and two arrows $u_1:i\to h_1$, $u_2:i\to h_2$, with $u_je_i=u_j$. There are no other arrows. In particular $h_1$ and $h_2$ have no common successor. To exhibit a specified nonaspherical final image, let $\Psi$ be the outgoing fork with vertices $0,1,2$ and arrows $0\to1$, $0\to2$. The functor $E\to\Psi$ sends $i$ to $0$, $h_j$ to $j$, and $e_i$ to $\id_0$; it is full and the identity on these labelled objects, hence final by Proposition~\ref{B:other-final}. The four-object category $C$ of the outgoing-fork counterexample has arrows
\[
 x\overset{r,s}{\rightrightarrows}0\xrightarrow{a}1,
 \qquad 0\xrightarrow{b}2,
 \qquad ar=as,\quad br=bs,
\]
and no other nonidentity arrows beyond the two indicated composites. The inclusion $\Psi\to C$ is final: $x\backslash\Psi$ is the four-cycle poset, while the other comma categories have initial objects. Proposition~\ref{B:collage} identifies $BC$ with $S^2$. Thus the composite $E\to C$ is final and its target has $H_2(BC;R)=R$ for every nonzero $R$. This proves every asserted failure.
\end{proof}

The finite program \texttt{code/root\_presheaf\_checks.py} verifies the action laws, the pushout of sets, all arrows and composites of $E$, finality of the two displayed functors, and the nerve-chain homology of $C$. The proof of totality quantifies over all target categories and is the split-presheaf calculation above; no finite search is offered as a replacement for that universal proof.

\begin{proposition}[Proved: minimality within the two-element idempotent base]\label{E:finite-idempotent-classification}
For the fixed base $U=\{1,e\}$ of Theorem~\ref{E:counterexample}, a finite presheaf $H$ is totally $\Winfty$-aspherical if and only if $H(*)$ is nonempty and $H(e)$ has exactly one fixed point. Such a presheaf has a totally $\Winfty$-aspherical category of elements if and only if $|H(*)|\leq2$. Thus three elements are the smallest counterexample on this base.
\end{proposition}
\begin{proof}
The idempotent function $H(e)$ retracts its set onto its fixed points. The connected components of $U/H$ are its fibers, one for each fixed point. The identity test for totality therefore requires exactly one fixed point $i$, and $H(e)$ is then constant with value $i$. If $H(*)=\{i\}$, the presheaf is $Q$ and is total by Lemma~\ref{E:idempotent}. Otherwise it is the union of $|H(*)|-1$ copies of $R$, all glued along their common split subpresheaf $Q$. Every left Kan extension is the corresponding finite union of copies of $C(-,c)$ along $Q_d$. Induction using the homotopy-pushout proof in Theorem~\ref{E:counterexample} makes its category of elements contractible. This proves totality.

With one element, the element category is the filtered monoid $U$. With two elements, the nonfixed element is a terminal object of the element category. With at least three elements, two distinct nonfixed elements are maximal and have no common successor, so the element category is not total by Lemma~\ref{C:directed-affine}.
\end{proof}

\begin{corollary}[Proved: the exact localizer range of the presheaf converse]\label{E:localizers}
For a fundamental localizer $W$, the implication ``every totally $W$-aspherical presheaf $H$ has totally $W$-aspherical category of elements holds for all small bases and presheaves if and only if $S^2$ is $W$-aspherical. For every fundamental localizer in which the displayed sphere category $C$ is not aspherical, Theorem~\ref{E:counterexample} gives a totally $W$-aspherical presheaf whose category of elements is not totally $W$-aspherical. In particular this holds for $\Winfty$, every nonzero homology localizer $W_R$, and every truncation localizer $W_n$ with $n\geq2$. For $W_0$ the two notions are equivalent to $U/H$ being nonempty and connected. For $W_1$ they are equivalent to $B(U/H)$ being nonempty, connected, and simply connected.
\end{corollary}
\begin{proof}
The positive direction of the universal assertion is Corollary~\ref{D:sphere-presheaves}. The negative direction and the stated counterexamples follow from the specified final map $E\to C$. For $W_0$, every canonical comparison $U/H\to U'/p_!H$ is final and therefore induces a bijection on components. For $W_1$, a final functor between nonempty connected categories induces an epimorphism on fundamental groups. To see this directly, for each object $b$ of the target choose an arrow $b\to f(a_b)$. Connectedness of $b\backslash A$ supplies a zigzag between any two choices, and for an arrow $b\to b'$ it supplies a zigzag between the choice at $b$ and the composite through the choice at $b'$. The commuting triangles identify each target edge-path, relative to these chosen paths, with the image of a source edge-path. With the basepoint choice an identity, every target loop is therefore in the image on $\pi_1$. Thus a simply connected $U/H$ makes every canonical final target simply connected. The converse is the identity test. The same argument using all final images characterizes total $W_1$-asphericity of the category itself.
\end{proof}

\subsection{Flat presheaves and ind-right adjoints}

\begin{proposition}[Proved: the genuine filtered-presentation criterion]\label{E:flat}
For a small $U$ and a presheaf $H$, the category $U/H$ is filtered if and only if $H$ is a small filtered colimit of representable presheaves. Hence this condition is sufficient for total $W$-asphericity of $H$ for every fundamental localizer. It is not necessary, by Theorem~\ref{E:counterexample}.
\end{proposition}
\begin{proof}
There is a canonical isomorphism of presheaves
\[
 H\cong\colim_{(u,h)\in U/H}U(-,u).
\]
At an object $v$, the map sends $(\alpha:v\to u,h)$ to $H(\alpha)h$. Every element $k\in H(v)$ is represented by $(\id_v,k)$. Every other representative of $k$ receives a morphism from that representative in the indexing category of elements, proving the required identification. If $U/H$ is filtered, this is the desired presentation.

Conversely let $H\cong\colim_{i\in I}U(-,f(i))$ for a small filtered $I$. The functor $f:I\to U$ is supplied by the full faithfulness of Yoneda, after choosing the indicated representing objects. The displayed presheaf is $f_!1_I$. Proposition~\ref{A:comprehensive} gives a final functor $I\to U/H$, and Lemma~\ref{A:filtered-final} makes $U/H$ filtered. The sufficiency for totality is then Proposition~\ref{A:elements-Kan} and Theorem~\ref{A:filtered}.
\end{proof}

\begin{proposition}[Refuted: the ind-right-adjoint biconditional in Lead 5]\label{E:ind-adjoint}
Let $i:U\hookrightarrow X$ be a sieve, and let $i_0:U\to\overline U$ have the codomain of Definition~\ref{A:definitions}. There is an ind-right adjoint precisely when every $U/x$, $x\in\overline U$, is filtered. Here an ind-right adjoint means a functor $r:\overline U\to\Ind(U)$ and bijections
\[
 \Hom_{\Ind(U)}(y_U(u),r(x))\cong\Hom_X(u,x)
\]
natural in $u\in U$ and $x\in\overline U$, where $y_U$ is Yoneda and $\Ind(U)$ is the full subcategory of presheaves given by small filtered colimits of representables. Such a sieve is totally $W$-dwy\'erienne for every fundamental localizer. The converse for $\Winfty$-dwy\'erienne sieves is false, even for a finite sieve, and the converse for totally $\Winfty$-dwy\'erienne sieves is also false in general.
\end{proposition}
\begin{proof}
Yoneda identifies $r(x)$, if it exists, with the restricted representable $u\mapsto\Hom_X(u,x)$. Proposition~\ref{E:flat} gives exactly the stated criterion. Naturality in $x$ is automatic from the restricted Yoneda functor and fullness of $\Ind(U)$ in the presheaf category. If these presheaves are representable, the same argument gives an ordinary right adjoint; this is the classical Dwyer condition.

For the first failed converse, adjoin to the monoid category $U$ of Theorem~\ref{E:counterexample} an object $t$ with crossing presheaf $H$, no arrows back to $U$, and only its identity as endomorphism. The inclusion $U\hookrightarrow X$ is a sieve with $\overline U=X$. Its weights are the representable at $*$ and $H$ at $t$, both total. It is therefore $\Winfty$-dwy\'erienne. But $U/t=U/H$ is not filtered, so no ind-right adjoint exists; indeed this sieve is not totally $\Winfty$-dwy\'erienne.

For the second failed converse, let $A=\Delta_{\mathrm{inj}}$ from Theorem~\ref{C:counterexample} and include it as a sieve in $A^{\triangleright}$, the category obtained by adjoining a terminal object. The crossing category at that terminal object is $A$, which is total but not filtered. The remaining crossing categories are representable slices and hence filtered. This sieve is totally $\Winfty$-dwy\'erienne and has no ind-right adjoint. Every functor and category used here is small in the fixed universe.
\end{proof}

\subsection{A precise restricted-action criterion for total presheaves}

\begin{proposition}[Proved: what the presheaf quantifier actually tests]\label{E:restricted-actions}
Let $E=U/H$ and $\pi:E\to U$. For a nonzero commutative ring $R$, $H$ is totally $W_R$-aspherical if and only if $BE$ is $R$-acyclic and every connected component of
\[
 \El_E(L\pi),\qquad L:U\to\Set,
\]
is $R$-acyclic. For $\Winfty$, the corresponding criterion is that $BE$ is weakly contractible and these same components are integrally acyclic. The restriction that the action factors through $\pi$ is essential; Theorem~\ref{E:counterexample} disproves removing it.
\end{proposition}
\begin{proof}
For necessity, fix $L$ and form the covariant collage $p:U\hookrightarrow U_L$ of Proposition~\ref{B:collage}. The presheaf $p_!H$ agrees with $H$ on $U$; at the new source object $x$, it is $\pi_0\El_E(L\pi)$. Its category of elements is obtained from $E$ by adjoining one source object for each of those components, with that component as its crossing element category. The same nerve computation as in Proposition~\ref{B:collage} realizes this as
\[
 BE\ \amalg^{h}_{B\El_E(L\pi)}\ \pi_0\El_E(L\pi).
\]
Both the resulting space and $BE$ are $R$-acyclic under totality, so the homology exact sequence makes the positive homology of each component vanish. Empty element categories impose no condition. For $\Winfty$, apply this argument with $R=\mathbb Z$, and use the identity test for $BE$.

For sufficiency, take arbitrary $p:U\to U'$ and the final comparison $q:E\to D=U'/p_!H$. For $d=(u',c)\in D$, its comma category $d\backslash E$ is exactly component $c$ of $u'\backslash E$. The latter is $\El_E(L_{u'}\pi)$ for the covariant functor $L_{u'}(u)=\Hom_{U'}(u',p(u))$. The hypotheses therefore make all these comma categories $R$-acyclic. Lemma~\ref{C:homology-A} gives an $R$-homology isomorphism $BE\to BD$, proving total $W_R$-asphericity. For $\Winfty$, $BD$ is integrally acyclic, and finality of $q$ makes its fundamental group a quotient of the trivial group $\pi_1(BE)$, by the edge-path argument in Corollary~\ref{E:localizers}. Thus $BD$ is simply connected and acyclic, hence weakly contractible. This proves the required universal condition on $p$.
\end{proof}

\section{Dwyer-type immersions and cofibrations}
\subsection{The three attachment conditions}

\begin{definition}[Known: the manuscript's variance]\label{F:definitions}
Let $i:U\hookrightarrow X$ be a full sieve of small categories, let
$T=X\setminus U$, and let $\overline U$ be the full cosieve generated by
$U$. For $t\in T$ put
\[
 H_t:U^{\op}\longrightarrow\Set,\qquad H_t(u)=\Hom_X(u,t).
\]
The immersion is $W$-dwy\'erienne if every nonempty $H_t$ is a totally
$W$-aspherical presheaf. It is totally $W$-dwy\'erienne if every nonempty
$U/t$ is a totally $W$-aspherical category. These are precisely the
conditions that $U\to\overline U$ be strictly, respectively totally,
$W$-aspherical in XII ms.~109. A classical Dwyer map additionally requires
that $U\to\overline U$ have a right adjoint. A $W$-cofibration is a functor
$i:U\to X$ such that
\[
 i_!:U\backslash\Cat\longrightarrow X\backslash\Cat,
 \qquad (U\to C)\longmapsto(X\to X\amalg_U C),
\]
preserves morphisms whose underlying functors belong to $W$.
\end{definition}

The last definition is exactly that of an $h$-cofibration, with weak
equivalences $W$, in \citet[Definition~1.1]{F_BataninBerger2017}. We use
``classical'' to retain Thomason's right-adjoint requirement
\citep[Definition~4.1]{F_Thomason1980}; Cisinski's pseudo-Dwyer maps have
only a retraction with a relative natural transformation
\citep[Definition~1]{F_Cisinski1999}. The two requirements must not be
identified.

\begin{proposition}[Proved: cobase change and the forward part of
Th\'eor\`eme~3]\label{F:pushout}
Let $W$ be a fundamental localizer. Both classes of dwy\'erienne sieves
in Definition~\ref{F:definitions} are stable under arbitrary cobase change
in $\Cat$. Every $W$-dwy\'erienne sieve is a $W$-cofibration. In a pushout
square
\[
\begin{tikzcd}
 U\arrow[r,"i",hook]\arrow[d,"p"']&X\arrow[d,"q"]\\
 C\arrow[r,"i'",hook]&X\amalg_U C,
\end{tikzcd}
\]
the implications $p\in W\Rightarrow q\in W$ and
$i\in W\Rightarrow i'\in W$ hold. The assertions apply to arbitrary
small $C$ and arbitrary $p$, with no injectivity or fullness assumption
on $p$.
\end{proposition}

\begin{proof}
The pushout has old part $C$, complementary full subcategory $T$, and
crossing presheaves $p_!H_t$. For another functor $r:C\to D$, the
canonical isomorphism
$(rp)_!H_t\cong r_!p_!H_t$ proves preservation of totality of the
presheaf. If $U/H_t$ is totally aspherical, the canonical functor
\[
 U/H_t\longrightarrow C/(p_!H_t)
\]
is ordinary final, so its target is again totally aspherical by the
final-image criterion. An empty weight stays empty. A nonempty weight
stays nonempty, since any element gives a representative in the value
of its Kan extension at its image object. Thus the new boundary is the
same full subcategory $\dot T=\overline U\cap T$ of $T$.

For completeness, we give the gluing argument needed for the
cofibration assertion. The two cosieves $\overline U,T\subset X$ cover
$X$, and their intersection is $\dot T$. Their nerves form a literal
union. The same is true after any cobase change. The map
$U\to\overline U$ is $W$-aspherical: its slices over objects of $U$
have terminal objects, and the other slices are the categories of
elements of the nonempty $H_t$. The relative Theorem~A axiom implies
$U\to\overline U\in W$.

We recall explicitly why the forward gluing assertion used here is
valid for every fundamental localizer. Replace a union of two cosieves
$V_1\cup_{V_0}V_2$ by the Grothendieck construction of the covariant
span $V_1\leftarrow V_0\to V_2$. The functor from that construction to
the union is a weak homotopy equivalence. Over an object belonging to
just one of the cosieves its comma category has the corresponding
terminal object. Over an object of the intersection its comma
category is the Grothendieck construction of the three slice
categories, each with its compatible terminal object; contraction to
the span of these three terminal objects contracts its nerve.
An objectwise $W$-equivalence of such spans induces a $W$-equivalence
on their Grothendieck constructions, by the relative Theorem~A axiom
over the indexing category. All comparison functors to the unions
belong to $W_\infty\subseteq W$. Consequently a morphism of the three
pieces which belongs to $W$ piece by piece induces a functor in $W$
on the unions. Only this forward gluing statement is used.

Now take a morphism $f:C\to D$ under $U$ in $W$. Write $\overline C$
and $\overline D$ for the generated cosieves in the two pushouts.
The inclusions $C\to\overline C$ and $D\to\overline D$ belong to
$W$ by the first paragraph and Theorem~A. Thus
$\overline C\to\overline D$ belongs to $W$ by two-out-of-three.
Glue this map with the identities on $\dot T$ and $T$. The preceding
argument proves
$X\amalg_U C\to X\amalg_U D\in W$, which is the full
$W$-cofibration assertion, including every preliminary change of base.
The final implication $i\in W\Rightarrow i'\in W$ follows from
Proposition~\ref{F:couniversal} below.
\end{proof}

\begin{proposition}[Proved: couniversality in the fundamental-localizer
setting]\label{F:couniversal}
For every fundamental localizer $W$, a functor is a couniversal
$W$-equivalence if and only if it is both a $W$-cofibration and a member
of $W$. The class of $W$-cofibrations is closed under composition,
cobase change, and retracts in the arrow category.
\end{proposition}

\begin{proof}
The closure assertions follow by composing the pushout functors, by
the associativity isomorphism for pushouts, and by applying a functor
to a retract diagram, respectively. If every cobase change of $i$
belongs to $W$, the natural square comparing two such changes gives
the preservation property by two-out-of-three; the identity change
also gives $i\in W$.

Conversely let $i:U\to X$ be a $W$-cofibration in $W$, and let
$p:U\to C$ be arbitrary. Factor $p=w j$, with $j:U\to M$ a
Thomason cofibration and $w:M\to C$ in $W_\infty$, using
\citet[Theorem~4.9]{F_Thomason1980}. Every Thomason cofibration is a
retract of a classical Dwyer map by
\citet[Proposition~2]{F_Cisinski1999}. Proposition~\ref{F:pushout}
and retract closure therefore make $j$ a $W$-cofibration. Applying
its pushout functor to $i\in W$ gives
$M\to X\amalg_U M$ in $W$. Applying the pushout functor of $i$
to $w\in W_\infty\subseteq W$ gives
$X\amalg_U M\to X\amalg_U C$ in $W$. Since $w$ is also in
$W$, the commuting square and two-out-of-three imply
$C\to X\amalg_U C\in W$.

This proof uses the fundamental-localizer hypothesis. It does not
turn the same conclusion into a theorem for an arbitrary class of
functors satisfying only two-out-of-three. For $W_\infty$ the result
also follows directly from the left proper case of
\citet[Lemma~1.10]{F_BataninBerger2017}.
\end{proof}

\begin{proposition}[Proved: the conditional equivalences and the
boundary corollary]\label{F:conditional-theorem3}
In the square of Proposition~\ref{F:pushout}, suppose that $i$ is
$W$-dwy\'erienne. Under the full axiom W(7 bis) of XII
ms.~76, the following conditions are equivalent:
\[
 i\in W,\qquad i'\in W,\qquad
 (\dot T\hookrightarrow T)\in W,\qquad
 i\text{ is a couniversal }W\text{-equivalence}.
\]
Under the full axiom W(7) of XII ms.~77 one also has
$p\in W\Longleftrightarrow q\in W$.
The equivalence between the first and last conditions requires
neither additional axiom.
\end{proposition}

\begin{proof}
Let $\overline C$ be the cosieve generated by $C$ in
$X\amalg_U C$. The boundary $\dot T$ is unchanged, and
$U\to\overline U$ and $C\to\overline C$ belong to $W$, by
Proposition~\ref{F:pushout}. Apply the cosieve-cover form of
W(7 bis) to each of
\[
 X=\overline U\cup_{\dot T}T,
 \qquad
 X\amalg_U C=\overline C\cup_{\dot T}T.
\]
It gives
$\overline U\to X\in W\Longleftrightarrow
\dot T\to T\in W\Longleftrightarrow
\overline C\to X\amalg_U C\in W$.
Two-out-of-three gives the first three stated equivalences.
The first and last are equivalent by
Propositions~\ref{F:pushout} and~\ref{F:couniversal}.

For the second assertion, the map between these two cover
squares is the identity on $\dot T$ and $T$. W(7), in its
cosieve-cover form, therefore identifies
$q\in W$ with $\overline U\to\overline C\in W$.
The commuting square whose other two sides are
$U\to\overline U$ and $C\to\overline C$ identifies the
latter condition with $p\in W$. This proves exactly
Th\'eor\`eme~3(a,b) and Corollaire~2 under their printed
extra hypotheses; Corollaire~1 and cobase-change
Corollaire~3 are strengthened by Proposition~\ref{F:pushout}.
\end{proof}

The reverse detection statements are not consequences of the
fundamental-localizer axioms. The acyclic noncontractible example
in Part~A explains why the full W(7 bis) cannot be assumed for
$W_\infty$. There are also typographical errors at
ms.~110 and 113: the pushout is $U'\amalg_U X$, with $X$, not $X'$,
as its original second entry.

\subsection{Composition of the total classes}

\begin{lemma}[Proved: operations on totally aspherical categories]
\label{F:total-operations}
For $W=W_\infty$, or $W=W_R$ for a nonzero commutative unital ring
$R$, connected components of discrete opfibrations over
a totally $W$-aspherical category are totally $W$-aspherical. Every
nonempty cosieve in such a category is totally $W$-aspherical. If
$B$ and every value of a small covariant functor $D:B\to\Cat$ are
totally $W$-aspherical, then $\int_BD$ is totally $W$-aspherical.
\end{lemma}

\begin{proof}
We use the action criterion proved in Part~C: a nonempty connected
category $A$ is totally $W_\infty$-aspherical precisely when each
connected component of every covariant category of elements over
$A$ is weakly contractible. To recall the essential point in its
necessity, the collage test first gives acyclicity. Every connected
covering of a category of elements is itself a discrete opfibration
over $A$. A covering corresponding to a nontrivial cyclic subgroup
of the fundamental group would have nonzero $H_1$, a contradiction.
Thus those components are simply connected and acyclic. For its
sufficiency, each comma category of a final functor out of $A$ is
such a connected category of elements, and Theorem~A applies.

Discrete opfibrations compose. A further covariant category of
elements over a component of one is therefore again a discrete
opfibration over $A$, after extending by the empty set on the other
components. The criterion proves the first assertion. A nonempty
cosieve $S\subset A$ is the category of elements of the covariant
indicator functor which is a singleton on $S$ and empty elsewhere.
Totality forces common successors, so $S$ is connected: two objects
of $S$ have a common successor in $A$, and that successor is still
in $S$. The second assertion follows.

Put $E=\int_BD$. Its nerve is nonempty and connected because the
base and all fibers are nonempty and connected. Let $H:E\to\Set$
be covariant. Its category of elements is canonically the
Grothendieck construction of the categories
$E_b=\operatorname{El}(H|_{D(b)})$, with the transition functors
induced by $D$ and $H$. Every connected component of $E_b$ is
contractible by the action criterion for $D(b)$. The natural maps
$E_b\to\pi_0 E_b$ are therefore weak homotopy equivalences.
Their Grothendieck constructions give a weak homotopy equivalence
\[
 \operatorname{El}(H)\longrightarrow
 \operatorname{El}\bigl(b\mapsto\pi_0 E_b\bigr).
\]
This is the usual objectwise invariance of Grothendieck
constructions, or the relative Theorem~A axiom over $B$.
Every component of the target is contractible, by the criterion for
$B$. Therefore every component of $\operatorname{El}(H)$ is
contractible. The criterion applied to $E$ proves the last assertion.
All the constructions use the actual transition functors. No
passage to an ind-object in a homotopy category is involved.
For $W_R$, use the homological action criterion of
Theorem~\ref{D:R-action}. The same argument replaces every
occurrence of ``contractible'' by ``$R$-acyclic'' and every
weak equivalence of Grothendieck constructions by an
$R$-homology equivalence. The common-successor condition
still holds by the same nonzero-$R$ collage obstruction.
This proves all three assertions in that case as well.
\end{proof}

\begin{theorem}[Proved: composition of total functors and total
immersions for weak and homological equivalences]\label{F:total-composition}
Totally $W$-aspherical functors and totally
$W$-dwy\'erienne sieves are closed under composition when
$W=W_\infty$ or $W=W_R$ for a nonzero commutative unital ring $R$.
\end{theorem}

\begin{proof}
Let $f:X\to Y$ and $g:Y\to Z$ be totally aspherical functors.
For $z\in Z$, consider the covariant diagram on $Y/z$
\[
 (y,\beta:g(y)\to z)\longmapsto X/y.
\]
An arrow $h:y\to y'$ in $Y/z$ acts by postcomposing
$\alpha:f(x)\to y$ with $h$. Its Grothendieck construction $E_z$
is totally aspherical by Lemma~\ref{F:total-operations}. There is
a functor in $\Cat$
\[
 r:E_z\longrightarrow X/z,
 \qquad (y,\beta,x,\alpha)\longmapsto
 (x,\beta\,g(\alpha)).
\]
Its section $s:X/z\to E_z$ sends $(x,\gamma)$ to
$(f(x),\gamma,x,\id_{f(x)})$. For $c=(x,\gamma)\in X/z$
and $d=(y,\beta,x',\alpha)\in E_z$, a morphism $s(c)\to d$
is uniquely determined by a morphism $k:x\to x'$ satisfying
$\beta g(\alpha f(k))=\gamma$; its $Y$-component is
$\alpha f(k)$. This is exactly a morphism $c\to r(d)$.
Consequently $s\dashv r$, with $rs=\id$, and $r$ is ordinary
final. Totality passes to its target $X/z$.

For sieves $U\subset X\subset Y$, fix $y\in Y$ with $U/y$
nonempty. Then $X/y$ is totally aspherical. In $X/y$ let $S_y$
be the full subcategory of $(x,\beta:x\to y)$ for which $U/x$
is nonempty. It is a nonempty cosieve, hence is totally aspherical
by Lemma~\ref{F:total-operations}. The diagram
$(x,\beta)\mapsto U/x$ on $S_y$ has totally aspherical values.
Its Grothendieck construction maps to $U/y$ by composition,
with the same adjoint section as above. Therefore $U/y$ is
totally aspherical. This proves exactly the condition required
for the composite sieve. Empty boundary slices impose no condition.
\end{proof}

\begin{theorem}[Proved: strict composition for every fundamental
localizer]\label{F:strict-composition}
If $f:X\to Y$ is strictly $W$-aspherical, restriction
$f^*:Y^{\wedge}\to X^{\wedge}$ takes totally
$W$-aspherical presheaves to totally $W$-aspherical
presheaves. Consequently strictly $W$-aspherical functors
are closed under composition for every fundamental
localizer $W$.
\end{theorem}

\begin{proof}
Let $H:Y^{\op}\to\Set$ be totally aspherical and fix
an arbitrary functor $p:X\to X'$. For $y\in Y$ set
\[
 Q_y=p_!\Hom_Y(f(-),y)\in (X')^{\wedge}.
\]
These presheaves form a covariant functor
$Q:Y\to(X')^{\wedge}$. Every $Q_y$ is totally
$W$-aspherical, by the hypothesis on $f$ and
transitivity of Kan extension. Form the small collage
$Z$ with full subcategories $X',Y$, no arrows from
$Y$ to $X'$, and
$\Hom_Z(x',y)=Q_y(x')$. Composition is given by the
contravariance in $x'$ and the covariance in $y$.
Write $j:Y\hookrightarrow Z$ for the complementary
cosieve and $F=j_!H$. Since $H$ is total, $Z/F$ is
$W$-aspherical. Its restriction to $Y$ is $H$, and
on $X'$ it is the presheaf
\[
 K=\int^{y\in Y} Q_y\times H(y)
   \cong p_! f^*H.
\]
The isomorphism uses co-Yoneda in $Y^{\wedge}$,
pointwise preservation of colimits by $f^*$, and
preservation of colimits by $p_!$. All are ordinary
presheaf colimits.

The inclusion $l:X'/K\hookrightarrow Z/F$ is a
full sieve. Its slice over an object $(x',k)$ of
$X'/K$ has terminal object $((x',k),\id)$.
Its slice over $(y,h)$, with $h\in H(y)$, is
canonically isomorphic to $X'/Q_y$: an element
$q\in Q_y(x')$ gives the unique source element
$[q,h]\in K(x')$ and the corresponding crossing
arrow $(x',[q,h])\to(y,h)$. Conversely each
crossing arrow has precisely these data. Morphisms
are the arrows of $X'$ carrying $q'$ to $q$;
the condition on $[q,h]$ then follows automatically.
Thus every slice of $l$ is $W$-aspherical. By the
relative Theorem~A axiom $l\in W$, and since
$Z/F\to e$ is in $W$, so is $X'/K\to e$.
The functor $p$ was arbitrary, proving totality of
$f^*H$.

If $g:Y\to Z_0$ is strictly aspherical, apply this
result to each totally aspherical presheaf
$\Hom_{Z_0}(g(-),z)$. Its restriction along $f$
is exactly $\Hom_{Z_0}(gf(-),z)$, proving the
composition assertion.
\end{proof}

\begin{theorem}[Proved: total-functor composition for every
fundamental localizer]\label{F:total-functor-general}
For every fundamental localizer $W$, totally $W$-aspherical
functors are closed under composition. For a small covariant
diagram $D:B\to\Cat$, if $B$ and every $D(b)$ are totally
$W$-aspherical, then $\int_BD$ is totally $W$-aspherical.
\end{theorem}

\begin{proof}
Let $f:X\to Y$ and $g:Y\to Z$ be totally
$W$-aspherical, and fix $z\in Z$. Define
\[
 f_z:X/z\longrightarrow Y/z,
 \qquad (x,\gamma:gf(x)\to z)\longmapsto(f(x),\gamma),
\]
where $X/z$ is formed using $gf$ and $Y/z$ using $g$.
For $(y,\beta:g(y)\to z)$, the comma category
$(X/z)/(y,\beta)$ is canonically isomorphic to $X/y$.
Indeed an object is $x$ and $h:f(x)\to y$, since the
remaining arrow is forced to be $\gamma=\beta g(h)$.
A morphism $k:x\to x'$ is subject exactly to
$h'f(k)=h$; its other required equation follows from
this one. Thus $f_z$ is totally $W$-aspherical.
Proposition~\ref{E:forward-subterminal} makes it strictly
$W$-aspherical. The terminal presheaf of $Y/z$ is total
because $g$ is total. By Theorem~\ref{F:strict-composition}
its restriction along $f_z$, the terminal presheaf of
$X/z$, is total. This proves totality of $gf$ for all
$z$, without an extra hypothesis on $W$.

For the second assertion put $E=\int_BD$ and let
$p:E\to B$ be its projection. For $b\in B$ define
\[
 r_b:E/b\to D(b),\qquad
 (b',x,\beta:b'\to b)\longmapsto D(\beta)(x).
\]
Its section $s_b:D(b)\to E/b$ sends $y$ to
$(b,y,\id_b)$. A morphism from $(b',x,\beta)$ to
$s_b(y)$ is exactly a morphism $D(\beta)(x)\to y$
in $D(b)$, so $r_b\dashv s_b$. Hence $s_b$ is
ordinary final. Since $D(b)$ is total, its final
image $E/b$ is total. Thus $p$ is a totally
$W$-aspherical functor. So is $B\to e$; their
composite is total by the first assertion. This
says precisely that $E$ is totally $W$-aspherical.
\end{proof}

\begin{theorem}[Proved: the exact fixed-localizer criterion for
total sieve composition]\label{F:total-sieve-cosieves}
For a fixed fundamental localizer $W$, totally
$W$-dwy\'erienne sieves are closed under composition
if and only if every nonempty full cosieve of every
totally $W$-aspherical category is totally
$W$-aspherical. For $W_n$ with $n\geq2$, composition
of totally $W_n$-dwy\'erienne sieves holds whenever
both inclusions are between finite categories
(finitely many arrows).
\end{theorem}

\begin{proof}
Assume the cosieve condition, and take totally
$W$-dwy\'erienne sieves $U\subset X\subset Y$.
For $y$ with $U/y$ nonempty, let $S_y$ be the
full subcategory of $X/y$ consisting of the
objects $(x,\beta)$ for which $U/x$ is nonempty.
It is a nonempty cosieve in the total category
$X/y$, so is total by hypothesis. Every value
of the covariant diagram $(x,\beta)\mapsto U/x$
on $S_y$ is total. Theorem~\ref{F:total-functor-general}
makes its Grothendieck construction total. The
composition functor from this construction to
$U/y$ has the adjoint section explicitly given
in the proof of Theorem~\ref{F:total-composition};
it is therefore final. Thus $U/y$ is total,
proving composition of the sieves.

Conversely let $A$ be totally $W$-aspherical and
let $V\subset A$ be a nonempty full cosieve.
Form $X$ with a full old copy of $A$, a full
lower copy $V_0$ of $V$, and crossing hom sets
\[
 X(v_0,a)=A(v,a),\qquad X(a,v_0)=\varnothing.
\]
All compositions are induced by composition in
$A$ and $V$. The sieve $V_0\hookrightarrow X$
is classical Dwyer: its generated cosieve is
$V\times[1]$, and projection to the lower
copy is right adjoint to its inclusion.
The functor $r:X\to A$ which forgets the
copy is left adjoint to the old inclusion
$j:A\to X$, because
$X(x,j(a))=A(r(x),a)$ naturally. Thus $j$
is final, and $X$ is total as a final image
of $A$. Adjoining a terminal object $t$ gives
a totally $W$-dwy\'erienne sieve
$X\hookrightarrow X^{\triangleright}$: its
only new boundary category is $X$. Under
the assumed composition property, the
composite from $V_0$ is totally dwy\'erienne.
Its nonempty boundary category over $t$
is exactly $V$. Therefore $V$ is total,
proving necessity.

For the finite assertion, every nonempty boundary
slice of either inclusion is a finite category
which is totally $W_n$-aspherical. By
Corollary~\ref{D:truncated-finite} it is filtered,
hence totally $W_\infty$-aspherical. Both
inclusions are therefore totally
$W_\infty$-dwy\'erienne. Their composite has
that property by Theorem~\ref{F:total-composition},
and consequently is totally $W_n$-dwy\'erienne.
\end{proof}

\begin{theorem}[Refuted: low-degree sieve composition]
\label{F:low-sieve-composition}
For each of $W_0$ and $W_1$, neither the
$W$-dwy\'erienne sieves nor the totally
$W$-dwy\'erienne sieves are closed under
composition. In each case there is a finite
counterexample whose first factor is classical
Dwyer, whose second factor is totally
$W$-dwy\'erienne, and whose final target has
six objects. Thus the assertion of total
sieve composition for every fundamental
localizer is false, whereas total functor
composition holds by
Theorem~\ref{F:total-functor-general}.
\end{theorem}

\begin{proof}
For $W_0$ take $V$ to be the discrete category
on $a,b$. For $W_1$ take the two-object
category with exactly two nonidentity arrows
$\alpha,\beta:a\to b$. Its nerve realization
is a circle: the two edges have the same
endpoints, and there are no nondegenerate
simplices in higher dimensions. In particular,
$V$ is not $W$-aspherical in the indicated
case; the second choice is connected but
has fundamental group $\mathbb Z$.

Put $B=V^{\triangleleft}$, with new initial
object $o$, and apply the lower-copy
construction in the preceding proof to
$V\subset B$. Explicitly, $X$ has the old
objects $o,a,b$, the lower objects $a_0,b_0$,
the arrows of $B$ and of $V_0$, and
\[
 X(v_0,b')=B(v,b'),\qquad
 X(b',v_0)=\varnothing
 \quad(v\in V,\ b'\in B).
\]
The only compositions not already internal
to the two copies are precomposition and
postcomposition of these crossing arrows
in $B$. Let $Y=X^{\triangleright}$, with
new terminal object $t$. These specifications
give every object, arrow, and composite.

The inclusion $a:V_0\hookrightarrow X$
is classical Dwyer by the explicit right
adjoint above. The retraction $r:X\to B$
has a natural transformation
$\id_X\Rightarrow jr$, identity on the
old copy of $B$ and equal to the crossing
identity $v_0\to v$ on the lower copy.
Thus $BX\simeq BB$, which is contractible
since $o$ is initial. Theorem~\ref{D:zero-one}
makes $X$ totally $W_0$- and totally
$W_1$-aspherical. Consequently
$b:X\hookrightarrow Y$ is totally
$W$-dwy\'erienne, and both factors are
also strictly $W$-dwy\'erienne.

The composite's boundary category at $t$
is $V_0/t\cong V$, since there is exactly
one arrow from each old object to $t$.
For total dwy\'erienness this category
already fails the identity test. For strict
dwy\'erienness the crossing presheaf on
$V_0$ is its terminal presheaf, whose
category of elements is again $V$; its
identity test fails as well. This proves
both failures of composition.

The program \texttt{code/f\_low\_truncation.py}
constructs these categories, checks every
composition and associativity equation,
checks the right adjoint and its triangle
identities, checks the stated natural
transformation, and constructs the boundary
slice with its isomorphism to $V$. The two
targets have respectively $6$ objects and
$15$ arrows, and $6$ objects and $21$ arrows.
Their normalized simplex counts are
$(6,9,4)$ and $(6,15,16,6)$, respectively.
Both targets have trivial reduced integral
homology; the composite boundary has
$H_0=\mathbb Z^2$ in the first example
and $(H_0,H_1)=(\mathbb Z,\mathbb Z)$
in the second. Complete target arrow and
composition tables are included in the
program's output.
\end{proof}

\begin{lemma}[Proved: nonempty cosieve restrictions for homology]
\label{F:cosieve-homology}
Let $R$ be a nonzero commutative unital ring, let
$H:U^{\op}\to\Set$ be totally $W_R$-aspherical,
and let $V\subset U$ be a full cosieve. If
$H|_V$ is nonempty, it is totally $W_R$-aspherical.
\end{lemma}

\begin{proof}
Fix $p:V\to C$ and set $K=p_!(H|_V)$.
Adjoin an initial object $0$ to $C$, obtaining
$B=C^{\triangleleft}$. There is a functor
$q:U\to B$ agreeing with $p$ on $V$ and
sending every object outside $V$ to $0$.
There are no arrows from $V$ to its complement,
and all remaining mixed arrows go to the unique
arrows from $0$, so this is well defined.
For $F=q_!H$ we have $F|_C=K$ and
$F(0)=\pi_0(U/H)=\{*\}$. The latter equality
follows by testing totality of $H$ on $U\to e$,
since a discrete set is $R$-acyclic only if it
is a singleton. The category $B/F$ is exactly
the initial cone on $E=C/K$.

Define a covariant functor $L:B\to\Set$ by
$L(0)=\{a,b\}$ and $L(c)=\{*\}$ for
$c\in C$. Every arrow from $0$ collapses
the two elements, and every arrow of $C$
acts identically. Adjoin an object $w$ with
outgoing hom functor $L$, obtaining the
collage $B_L$, and let $r:B\to B_L$ be
the inclusion. The category of elements of
$L$ pulled back to $B/F$ consists of two
initial cone points on $E$. Its classifying
space is $\Sigma BE$. Because $K$ is nonempty,
this category is connected, and the ordinary
Kan extension has $(r_!F)(w)=\{*\}$.
The collage formula for the new category of
elements therefore gives
\[
 B(B_L/r_!F)\simeq
 \operatorname{hpo}\bigl(*\leftarrow
     \Sigma BE\to B(B/F)\bigr)
 \simeq\Sigma^2 BE.
\]
Here $B(B/F)$ is contractible, and the last
display uses the ordinary unreduced
suspension. Since $F$ is totally $W_R$-
aspherical, the left side is $R$-acyclic.
The reduced suspension homology isomorphism
and nonemptiness of $E$ imply that $E$ is
nonempty, connected, and $R$-acyclic. The
functor $p$ was arbitrary. Thus $H|_V$
is totally $W_R$-aspherical.
\end{proof}

\begin{theorem}[Proved: composition of homological dwy\'erienne
immersions]\label{F:strict-sieve-homology}
For every nonzero commutative unital ring $R$, the
$W_R$-dwy\'erienne sieves are closed under composition.
\end{theorem}

\begin{proof}
Let $U\subset X\subset Y$ be two such sieves,
and let $y\in Y$ receive an arrow from $U$.
The weight $H=X(-,y)$ is total on $X$.
Let $V\subset X$ be the cosieve generated by $U$.
Its restriction $H|_V$ is nonempty and is total
by Lemma~\ref{F:cosieve-homology}. The functor
$U\to V$ is strictly $W_R$-aspherical by
definition. Theorem~\ref{F:strict-composition}
therefore shows that $H|_U=Y(-,y)|_U$ is
totally $W_R$-aspherical. These are exactly
the nonempty crossing weights of the composite.
\end{proof}

\subsection{Discretely flat inclusions and strict separations}

\begin{definition}[Known: discretely flat]\label{F:DFdefinition}
For a small category $A$, let $\operatorname{disc}A$ denote the
simplicial category with its hom sets regarded as discrete
simplicial sets. A functor $i:U\to X$ is discretely flat if
\[
 \operatorname{disc}X\amalg_{\operatorname{disc}U}(-)
\]
preserves Dwyer--Kan equivalences of simplicial categories under
$\operatorname{disc}U$. It is a discretely flat cofibration when
it is also injective on objects, and a discretely flat inclusion
when it is also faithful. These are
\citet[Definition~3.4]{F_HORR2024}. A sieve satisfies both extra
injectivity conditions.
\end{definition}

\begin{proposition}[Proved: an exact weight test for discrete flatness]
\label{F:weight-test}
A sieve $i:U\hookrightarrow X$ is discretely flat if and only if,
for each $t\in X\setminus U$, the functor
\[
 H_t\otimes_U-:\operatorname{Fun}(U,\mathbf{sSet})
       \longrightarrow\mathbf{sSet}
\]
preserves objectwise weak homotopy equivalences. The coend is formed
degree by degree in sets. Every totally $W_\infty$-dwy\'erienne
sieve is discretely flat. Every sieve with discrete source is
discretely flat.
\end{proposition}

\begin{proof}
In a simplicial pushout of the sieve along
$p:\operatorname{disc}U\to C$, the object set is
$\operatorname{Ob}C\amalg\operatorname{Ob}T$. The hom simplicial
sets internal to $C$ are unchanged, those internal to $T$ are the
original discrete hom sets, and there are no arrows from $T$ to
$C$. The remaining hom simplicial sets are
\[
 (X\amalg_U C)(c,t)=H_t\otimes_U C(c,p(-)).
\]
This follows degree by degree from the ordinary collage pushout
formula; composition is induced by precomposition in $C$ and
postcomposition in $T$. If all the displayed weight functors
preserve weak equivalences, a Dwyer--Kan equivalence $C\to D$
under $U$ induces weak equivalences on all these hom objects.
Essential surjectivity is inherited from $C\to D$, together with
the identity on the unchanged objects of $T$. This proves
sufficiency.

For necessity, given a diagram $K:U\to\mathbf{sSet}$ form the
simplicial collage $C_K$ obtained from $\operatorname{disc}U$ by
adjoining an object $x$, setting $C_K(x,u)=K(u)$, and having no
arrows into $x$ except its identity. An objectwise weak equivalence
$K\to L$ induces a Dwyer--Kan equivalence $C_K\to C_L$,
identity on objects. After pushout, the map on hom objects from
$x$ to $t$ is exactly $H_t\otimes_U K\to H_t\otimes_U L$.
Discrete flatness forces this to be a weak equivalence.

We spell out the sufficient totality assertion. If a category $A$
is totally aspherical, ordinary colimits
$\operatorname{Fun}(A,\mathbf{sSet})\to\mathbf{sSet}$ preserve
objectwise weak equivalences. Indeed the augmentation from the
simplicial replacement to the ordinary colimit is a weak
equivalence: in each simplicial degree it is the map from the
nerve of the category of elements of a set-valued diagram to its
set of components, and all these components are contractible by
the action criterion. Taking the diagonal of the resulting
bisimplicial sets preserves these weak equivalences. The ordinary
colimit is thus naturally weakly equivalent to the homotopy
colimit, which preserves objectwise weak equivalences.
For $A=U/H_t$ we have the canonical degreewise isomorphism
\[
 H_t\otimes_U K\cong\colim_{(u,h)\in U/H_t}K(u).
\]
This proves the assertion for totally dwy\'erienne sieves.
If $U$ is discrete, the weight functor is simply
$\coprod_{u\in U}H_t(u)\times K(u)$, which plainly preserves
weak equivalences, including empty weights and empty $U$.
\end{proof}

\begin{proposition}[Proved: the restricted-action criterion for
discrete flatness]\label{F:DF-actions}
Let $H:U^{\op}\to\Set$, put $E=U/H$, and let
$\pi:E\to U$ be the projection. The weighted colimit
$H\otimes_U-:\mathbf{sSet}^{U}\to\mathbf{sSet}$
preserves objectwise weak equivalences if and only if
every connected component of
\[
 \operatorname{El}_{E}(L\pi),\qquad L:U\to\Set,
\]
is weakly contractible. Among weights satisfying this
condition, $H$ is totally $W_\infty$-aspherical if and
only if $E$ is nonempty and connected. Consequently a
discretely flat sieve is $W_\infty$-dwy\'erienne if and
only if each of its nonempty boundary categories $U/t$
is connected.
\end{proposition}

\begin{proof}
For a covariant set-valued functor $L$, define a
covariant simplicial diagram $PL$ by
\[
 PL(u)_n=
 \coprod_{u_0\to\cdots\to u_n\to u}L(u_0).
\]
The face and degeneracy maps are those of the usual
simplicial replacement. Its augmentation to the
discrete simplicial set $L(u)$ sends an element along
the indicated composite. For fixed $u$, $PL(u)$ is
the nerve of the category of triples
$(v,l,\alpha:v\to u)$, with the evident commuting
triangles as arrows. Each component has the terminal
triple $(u,l_u,\id_u)$ for its common image
$l_u\in L(u)$. Thus $PL\to L$ is an objectwise
weak equivalence.

The co-Yoneda isomorphism identifies
\[
 (H\otimes_U PL)_n
   \cong
 \coprod_{u_0\to\cdots\to u_n}
       L(u_0)\times H(u_n)
   \cong
 N\operatorname{El}_E(L\pi)_n.
\]
For the last identification, the element of $H(u_n)$
determines all its predecessors by contravariance,
and the element of $L(u_0)$ determines all its
successors by covariance. These identifications
respect every face and degeneracy. Also
$H\otimes_U L\cong\pi_0\operatorname{El}_E(L\pi)$.
Hence homotopical preservation by the weight forces
the map from this nerve to its set of components to
be a weak equivalence, proving necessity.

Conversely, for $K:U\to\mathbf{sSet}$ consider the
bisimplicial two-sided bar construction
\[
 B_{q,n}(H,U,K)=
 \coprod_{u_0\to\cdots\to u_q}
       H(u_q)\times K(u_0)_n.
\]
For fixed $n$, its augmentation to the set
$(H\otimes_U K)_n$ is precisely the map just
described for $L=K_n$. The assumed condition makes
that augmentation a weak equivalence of simplicial
sets in the $q$ direction. The diagonal therefore
gives a natural weak equivalence
$\operatorname{diag}B(H,U,K)\to H\otimes_U K$.
For an objectwise weak equivalence $K\to K'$, each
fixed-$q$ map of bars is a coproduct of copies of
objectwise weak equivalences and is therefore a weak
equivalence. Its diagonal is a weak equivalence as
well. Two-out-of-three proves the required
preservation by $H\otimes_U-$.

Taking $L$ constant at a singleton shows that every
component of $E$ itself is contractible. Under
nonempty connectedness, the restricted-action
criterion of Proposition~\ref{E:restricted-actions}
therefore proves total $W_\infty$-asphericity of $H$.
The converse follows from the identity test for
totality. Apply this argument weight by weight and
use Proposition~\ref{F:weight-test} for the final
assertion.
\end{proof}

\begin{example}[Proved: strict does not imply total, even among
discretely flat sieves]\label{F:strict-not-total}
Let $U$ have one object and endomorphism monoid
$M=\{1,e\}$, with $e^2=e$. Let the presheaf
$H$ have elements $i,h_1,h_2$, with $H(e)$
constant at $i$, and form $X$ by adjoining one
object $t$ with $\Hom_X(-,t)|_U=H$ and no
arrows from $t$ to $U$. Then $U\hookrightarrow X$
is $W_\infty$-dwy\'erienne and discretely flat,
but is not totally $W_\infty$-dwy\'erienne.
\end{example}

\begin{proof}
The presheaf counterexample in Part~E proves
totality of $H$ and failure of totality of $U/H$.
We recall its construction to make the new claim
about discrete flatness independently checkable.
Let $R=U(-,*)$ and $Q=\{e\}\subset R$.
Multiplication by $e$ is an idempotent retraction
of $R$ onto $Q$, and
$H=R\amalg_QR$. After any Kan extension the
two copies of $R$ remain representable and
the two copies of $Q\to R$ remain split
monomorphisms. The category of elements of the
image of $Q$ is filtered, while those of the
two images of $R$ have terminal objects; their
union is weakly contractible. Thus $H$ is total.
The objects $(i,h_1,h_2)$ of $U/H$ have
arrows from $i$ to each $h_j$, and the two
distinct $h_j$ have no common successor.
Consequently $U/H$ is not total.

For a covariant simplicial diagram $L:U\to
\mathbf{sSet}$, put $V=L(*)$ and let $d=L(e)$.
Then
\[
 H\otimes_U L\cong V\amalg_{\operatorname{im}d}V.
\]
The image of a simplicial idempotent is a
retract, so its inclusion in $V$ is a
simplicial monomorphism. If $L\to L'$ is
an objectwise weak equivalence compatible
with the idempotents, the induced map on
their images is a retract of the underlying
weak equivalence in the arrow category.
The two pushout diagrams therefore have
weakly equivalent entries and cofibration
legs. Their pushouts are weakly equivalent.
Proposition~\ref{F:weight-test} gives discrete
flatness. This also proves that the first
inclusion between the total and strict
dwy\'erienne classes is strict.
\end{proof}

\begin{proposition}[Known, with explicit consequence: discrete
flatness implies homotopical flatness]\label{F:DFh}
Every discretely flat sieve is a $W$-cofibration for every
fundamental localizer $W$. In particular it is an $h$-cofibration
for $W_\infty$.
\end{proposition}

\begin{proof}
For such an inclusion, \citet[Theorem~1.5, in its generalized
form]{F_HORR2024} show that the canonical comparison
\[
 NX\amalg_{NU}NC\longrightarrow N(X\amalg_U C)
\]
is a categorical equivalence, hence a weak homotopy equivalence.
Since $NU\to NX$ is a monomorphism, the source computes the
homotopy pushout of the three nerves. Compare this description
for a morphism $C\to D$ under $U$ in $W$. The homotopy pushout
comparison is a $W$-equivalence by objectwise invariance of the
Grothendieck construction of a span, as used in
Proposition~\ref{F:pushout}. The two comparison maps belong to
$W_\infty\subseteq W$. Two-out-of-three gives the claim.
\end{proof}

\begin{theorem}[Refuted: C035, even for a couniversal equivalence]
\label{F:C035}
A sieve which is a couniversal $W_\infty$-equivalence need not be
$W_\infty$-dwy\'erienne. The same counterexample works for every
nonzero homology localizer and every $W_n$ with $n\geq0$.
\end{theorem}

\begin{proof}
Let $X$ have objects $a,t,s$, identity arrows, two arrows
$\alpha,\beta:a\to t$, an arrow $\gamma:t\to s$, and an
arrow $\delta:a\to s$. The only composites of two nonidentity
arrows are
\[
 \gamma\alpha=\gamma\beta=\delta.
\]
Put $U=\{a\}$. This is a full sieve and generates all of $X$
as a cosieve. Its crossing presheaf at $t$ has two elements, so
$U/t$ is a discrete two-object category. Already the identity
test shows that this presheaf is not totally aspherical for any
of the localizers in the statement.

Let $p:U\to C$ be any functor, with selected object $c=p(a)$,
and put $P=X\amalg_U C$. The objects are those of $C$ together
with $t,s$. The new hom sets, for $d\in C$, are
\[
 P(d,t)=C(d,c)\times\{\alpha,\beta\},\qquad
 P(d,s)=C(d,c),\qquad P(t,s)=\{\gamma\}.
\]
Composition with $\gamma$ forgets the second coordinate. There
are no arrows from $t$ or $s$ into $C$, and the other hom sets
are as stated in $X$ or $C$. Define $r:P\to C$ by the identity
on $C$, by $r(t)=r(s)=c$, by projecting these new hom sets to
$C(d,c)$, and by $r(\gamma)=\id_c$. Let $j:C\to P$ be the
inclusion. Define $F:P\to P$ to be the identity on $C$ and to
send both $t,s$ to $s$, sending each crossing arrow to its
composite to $s$. There are natural transformations
\[
 \id_P\Longrightarrow F\Longleftarrow jr
\]
whose components outside $C$ are, respectively,
$(\gamma,\id_s)$ and the arrows $c\to s$ corresponding to
$\id_c$. Both transformations are identities on $C$. Naturality
at the two arrows to $t$ is precisely the displayed relation
$\gamma\alpha=\gamma\beta$. Thus $Bj$ is a homotopy
equivalence with inverse $Br$, for every $p$. This proves
couniversality for $W_\infty$ and hence for every fundamental
localizer. The failure of dwy\'erienne totality remains visible
in degree zero for the stated localizers. By
Proposition~\ref{F:weight-test} the sieve is also discretely flat.

The program \texttt{code/f\_check.py} verifies the seven-arrow
category, all $21$ composable triples, and the two natural
transformations. The universal proof above, rather than a finite
sample of targets, supplies the required quantifier over $p$.
\end{proof}

\begin{corollary}[Refuted: one-point and complementary-change
cofibration recognition]\label{F:onepoint-cofibration}
The ordinary-cofibration assertion at XII ms.~111
(C037) is false: a $W$-cofibration sieve whose
complementary full category is isomorphic to $e$
(one object and only its identity) need not be
$W$-dwy\'erienne. The
complementary-change recognition assertion (C038)
is also false: even requiring every change of the
complement to give a $W$-cofibration does not force
dwy\'erienness. These failures hold for $W_\infty$,
for $W_R$ with any nonzero commutative unital ring
$R$, and for every $W_n$, $n\geq0$.
\end{corollary}

\begin{proof}
Let $U=\{a\}$, and let $X$ have the two objects
$a,t$ and exactly two nonidentity arrows
$\alpha,\beta:a\to t$. Its source is discrete,
so Proposition~\ref{F:weight-test} makes $U\to X$
discretely flat, and Proposition~\ref{F:DFh}
makes it a $W$-cofibration for every fundamental
localizer. Its only boundary presheaf has two
elements, so fails the identity test for total
asphericity for all localizers listed above.

For every small category $T_1$ and the unique
functor $T_1\to\{t\}$, changing the complement
produces the collage with source $U$ and two
crossing arrows to each object of $T_1$, with
each transition preserving the two labels.
Its source is still discrete. The same two
propositions make every such changed sieve a
$W$-cofibration, including the empty-complement
case. This proves the failure of C038. In
particular, retaining full W(7) does not repair
these assertions, since nonzero homology
localizers satisfy that axiom.

The two-object category is the parallel-arrow
category explicitly constructed and checked by
\texttt{code/f\_low\_truncation.py}. Its nerve
is a circle, so this particular inclusion is
not couniversal for $W_\infty$ or for nonzero
homology localizers. The following proposition
treats that narrower variant precisely.
\end{proof}

\begin{proposition}[Proved: the exact one-point couniversality
criterion]\label{F:onepoint-couniversal}
Let $H:U^{\op}\to\Set$ be any presheaf on a small
category, and let $X$ be obtained by adjoining one
object $t$, with $X(u,t)=H(u)$, no arrows from
$t$ to $U$, and only the identity endomorphism
of $t$. Then $U\hookrightarrow X$ is a
couniversal $W_\infty$-equivalence if and only
if $H$ is totally $W_{\mathbb Z}$-aspherical.
For a nonzero commutative unital ring $R$, it
is a couniversal $W_R$-equivalence if and only
if $H$ is totally $W_R$-aspherical.
\end{proposition}

\begin{proof}
After arbitrary cobase change $p:U\to C$, the
new weight is $F=p_!H$. Write $E=C/F$ and
$\pi:E\to C$. The one-point collage formula
of Lemma~\ref{A:audit-cylinder} identifies
the specified inclusion of classifying spaces
with the inclusion into the homotopy cofiber
\[
 BC\longrightarrow
 BC\amalg^h_{BE}*.
\]
If $E$ is empty, this adds a separate point
and is not a homology equivalence over a
nonzero ring. If $E$ is nonempty, its relative
homology is $\widetilde H_{k-1}(BE;R)$.
Consequently this inclusion is an $R$-homology
equivalence exactly when $E$ is $R$-acyclic.
Quantifying over all $p$ proves the assertion
for $W_R$ and the necessity of total
$W_{\mathbb Z}$-asphericity in the first
assertion, including nonemptiness.

We prove the remaining sufficiency. First let
$F$ be any totally $W_{\mathbb Z}$-aspherical
presheaf on $C$. Its element category $E$
is nonempty, connected, and acyclic. Moreover
the homomorphism
\[
 \varphi:\pi_1(BE,e)\longrightarrow
 G=\pi_1(BC,\pi(e))
\]
induced by its projection is trivial. To
verify this, take $g\in\operatorname{im}\varphi$
and $K=\langle g\rangle\subset G$. On the
component of $C$ containing $\pi(e)$ take
the covariant set-valued local system
associated to the transitive $G$-set with
stabilizer $K$, and take the empty set on
the other components. Its pullback to $E$
is a covering action. The component through
the distinguished coset has fundamental
group $\varphi^{-1}(K)$. It is acyclic by
Proposition~\ref{E:restricted-actions}.
But $\varphi^{-1}(K)\to K$ is surjective,
since $K\subset\operatorname{im}\varphi$.
As $K$ is abelian, this surjection factors
through the zero first homology group of
that covering component. Hence $K=1$.
This applies to every $g$, proving the claim.

Consider the connected component of $BC$
to which $BE$ maps; all other components
are unchanged by the attachment. Van Kampen
and $\varphi=0$ show that adjoining the
cone on $BE$ preserves its fundamental
group. On universal covers the same
attachment adjoins a disjoint family of
cones on copies of $BE$, indexed by $G$:
the pullback of the universal cover to
$BE$ is trivial because $\varphi=0$.
The relative integral homology of the
map of universal covers is thus
\[
 \bigoplus_{G}\widetilde H_{k-1}(BE;\mathbb Z)=0.
\]
The map is a homology equivalence between
simply connected CW complexes, hence a
weak homotopy equivalence by the homological
Whitehead theorem. Together with the
fundamental-group isomorphism this proves
that the original cofiber inclusion is a
weak homotopy equivalence.

If $H$ is total $W_{\mathbb Z}$, each
$F=p_!H$ is total $W_{\mathbb Z}$, since
$q_!p_!H\cong(qp)_!H$ for every subsequent
functor $q$ with source $C$. The
preceding argument applies after every
cobase change, proving couniversality
for $W_\infty$.
\end{proof}

\begin{theorem}[Proved: a homotopical cofibration need not be
discretely flat]\label{F:h-not-DF}
There is a four-object full sieve $U\hookrightarrow X$ which
is a couniversal $W_\infty$-equivalence but is not discretely
flat. Consequently the inclusion of discretely flat sieves in
the class of $W_\infty$-cofibration sieves is strict.
\end{theorem}

\begin{proof}
Let $U=[1]=(0\xrightarrow{a}1)$. Adjoin objects $t,s$ to form
$X$. The new arrows to $t$ are
$\beta,\gamma:1\to t$ and $c:0\to t$, with
$\beta a=\gamma a=c$. Adjoin a terminal object $s$ to this
three-object category. These rules specify all arrows and all
compositions. The inclusion of $U$ is a sieve. The retraction
$X\to U$ fixes $U$ and sends $t,s$ to $1$. The endofunctor
fixing $U$ and sending $t,s$ to $s$ again gives
$\id_X\Rightarrow F\Leftarrow ir$, relative to $U$.
Products with $[1]$ preserve pushouts in $\Cat$, so a natural
transformation relative to $U$ passes to every pushout. Therefore
every pushout inclusion is a categorical homotopy equivalence.

To disprove discrete flatness, take the simplicial circle
$K=\Delta[1]/\partial\Delta[1]$ and its simplicial cone $CK$.
Define diagrams $L,M:U\to\mathbf{sSet}$ by
\[
 L(0)=K,\quad L(1)=CK,\quad L(a):K\hookrightarrow CK;
 \qquad
 M(0)=K,\quad M(1)=*,\quad M(a):K\to *.
\]
The natural transformation $L\to M$ is the identity at $0$
and the contraction $CK\to *$ at $1$, so is an objectwise
weak equivalence. The weight $H_t$ has one element at $0$,
two at $1$, and its two elements restrict to the same element.
Its coends are
\[
 H_t\otimes_U L=CK\amalg_K CK,\qquad
 H_t\otimes_U M=*\amalg_K *=*.
\]
The first realizes as the suspension of a circle, hence has
$H_2\cong\mathbb Z$; the second is a point. Thus the map is
not a weak equivalence. Proposition~\ref{F:weight-test}, or its
explicit simplicial-collage construction, disproves discrete
flatness. The finite category and its relative transformations
are checked in the accompanying program.
\end{proof}

\begin{proposition}[Proved: a strict extension of classical Dwyer
maps]\label{F:classical-strict}
The inclusion $\mathbb N\hookrightarrow\mathbb N^{\triangleright}$,
where $\mathbb N$ has its usual order and the target adds a new
terminal object, is totally $W$-dwy\'erienne for every
fundamental localizer $W$, but is not a classical Dwyer map.
\end{proposition}

\begin{proof}
The only new crossing category is $\mathbb N$, which is
filtered and therefore totally $W$-aspherical. A right adjoint
would assign to the new terminal object a terminal object of
$\mathbb N$, which does not exist. The generated cosieve is
the entire target, so no smaller neighborhood can avoid this
obstruction.
\end{proof}

Thus the proved inclusions, among sieves, are
\[
 \{\text{classical Dwyer}\}
 \varsubsetneqq
 \{\text{totally }W_\infty\text{-dwy\'erienne}\}
 \varsubsetneqq
 \{\text{discretely flat}\}
 \varsubsetneqq
 \{W_\infty\text{-cofibrations}\},
\]
and also
\[
 \{\text{totally }W_\infty\text{-dwy\'erienne}\}
 \varsubsetneqq
 \{W_\infty\text{-dwy\'erienne}\}
 \varsubsetneqq
 \{W_\infty\text{-cofibrations}\}.
\]
The first inclusion into the discretely flat class is strict as
well: Theorem~\ref{F:C035} belongs to that class and fails even
the strictly dwy\'erienne condition. In particular discretely
flat sieves are not contained in the strictly dwy\'erienne
class.

\subsection{The endpoint product is not a homotopical cofibration}

\begin{theorem}[Refuted: the proposed endpoint stability, in its
smallest standard test]\label{F:interval-failure}
Let $n\geq2$, let
\[
 I=\Phi_n=(0\to1\leftarrow2\to\cdots\leftarrow2n),
 \qquad j:\{0,2n\}\hookrightarrow I,
\]
and let $i:\{0\}\hookrightarrow[1]$ select the initial object.
The sieve
\[
 \alpha=i\mathbin{\square}j:
 A=(\{0\}\times I)\cup([1]\times\{0,2n\})
       \hookrightarrow P=[1]\times I
\]
is not a $W_\infty$-cofibration. It is not a $W$-cofibration
for any fundamental localizer $W$ in which a category of
homotopy type $S^2$ is not aspherical.
\end{theorem}

\begin{proof}
Write $l_k=(0,k)$ and $u_k=(1,k)$. The objects of $A$ are
$l_0,\ldots,l_{2n},u_0,u_{2n}$. Its nonidentity arrows form
an unoriented tree, so $BA$ is contractible. Also $BP$ is
contractible, being the product of two contractible nerves.

Define $C$ to have objects $x,z,a,b$, two arrows
$\lambda,\mu:x\to z$, arrows $r:z\to a$ and $v:z\to b$,
and their composites, with relations
\[
 r\lambda=r\mu,\qquad v\lambda=v\mu.
\]
There are no other nonidentity arrows. The full category on
$x,z$ has classifying space $S^1$; the full subcategories
obtained by adjoining $a$ and $b$ are its two terminal cones.
Their nerves cover $NC$ and intersect in that circle. Hence
$BC\simeq S^2$. Let $D=C^{\triangleright}$, with new terminal
object $s$, so $BD$ is contractible.

There is a functor $p:A\to D$ defined by
\[
 p(l_0)=z,\quad p(u_0)=b,\quad
 p(l_k)=a\ (1\leq k\leq2n),\quad p(u_{2n})=a.
\]
It sends $l_0\to l_1$ to $r$, $l_0\to u_0$ to $v$, and
every remaining nonidentity arrow to $\id_a$. This types
correctly because every remaining arrow has both endpoints
among the vertices sent to $a$. It is a weak equivalence,
since its source and target are contractible.

Compute the ordinary pushout $Q=P\amalg_A D$. Its complement
to $D$ consists of $u_1,\ldots,u_{2n-1}$, with the upper
alternating-interval arrows. The weight for $u_1$ before
pushout is subterminal on the full four-object subcategory
\[
 S=\{l_0,l_1,l_2,u_0\},\qquad
 u_0\leftarrow l_0\to l_1\leftarrow l_2.
\]
The functor $S\to\Psi=(z\to a,z\to b)$ induced by $p$
is ordinary final: its outgoing comma categories at $z,a,b$
are respectively the connected tree $S$, the arrow
$l_2\to l_1$, and the singleton $u_0$. The inclusion
$\Psi\to C$ is final. At $x$ its outgoing comma category
is the four-edge square with vertices
$(z,\lambda),(z,\mu),(a,r\lambda),(b,v\lambda)$;
at the other objects it has an initial object. Consequently
the pushed-forward crossing weight at $u_1$ is the
subterminal presheaf on $D$ supported on $C$.

For every $k\geq2$, all vertices of $A/u_k$ are sent to
$a$, all their arrows to $\id_a$, and $A/u_k$ is a nonempty
connected tree (or singleton). The pushed-forward weight is
therefore the representable $D(-,a)$. All transitions between
these latter weights are identities; the transition from
$u_2$ to $u_1$ is the canonical map $D(-,a)\to1_C$.
These statements can also be read directly from the coend
\[
 \Hom_Q(d,u_k)=\int^{v\in A}
       D(d,pv)\times\Hom_P(v,u_k).
\]
They describe every hom set and every composition of $Q$.

Let $R\subset Q$ be the full subcategory on $D$ and $u_1$.
It is exactly $C$ with two terminal cone points $s,u_1$,
with no arrow between the two points. Define $r:Q\to R$
to fix $R$ and to send each $u_k$, $k\geq2$, to $a$.
On the representable crossing hom sets it is the Yoneda
identification with $D(-,a)$, and on arrows among those
$u_k$ it is $\id_a$. The arrow $u_2\to u_1$ goes to the
unique arrow $a\to u_1$. For the full inclusion $e:R\to Q$
we have $re=\id_R$ and a natural transformation
$er\Rightarrow\id_Q$, identity on $R$, whose component
at $u_k$ is the arrow $a\to u_k$ represented by $\id_a$.
The above description of the transitions checks naturality
on every crossing and upper-interval arrow. Hence $Be$ is
a homotopy equivalence. The literal nerve union of the two
cones gives
\[
 BQ\simeq BR\simeq\Sigma BC\simeq S^3.
\]

If $\alpha$ were a $W$-cofibration, applying its pushout
functor to the morphism $p:A\to D$ under $A$, with source
the identity of $A$, would force $P\to Q$ to belong to
$W$. Since $P$ is contractible and $Q$ has type $S^3$,
this disproves the cofibration assertion whenever $S^3$
is not $W$-aspherical.

To obtain the stated stronger obstruction with $S^2$,
use instead the same assignment $p_0:A\to C$, without
adjoining $s$. The identical coend calculation shows
that $Q_0=P\amalg_A C$ retracts by a relative natural
transformation to $C^{\triangleright}_{u_1}$. Hence
$BQ_0$ is contractible. The map $C\to Q_0$ is not
in $W$ if $S^2$ is not $W$-aspherical. But $\alpha$
itself is in $W_\infty\subseteq W$. If it were a
$W$-cofibration, Proposition~\ref{F:couniversal}
would make every cobase change a $W$-equivalence,
contradicting this one. This argument also handles
$W_2$, where the preceding $S^3$ comparison alone
would detect no failure.

For $n=2$ the program computes the actual pushout from the
coend relations. It has eight objects and thirty arrows,
with normalized simplex counts $(8,22,30,20,4)$ and integral
homology ranks $(1,0,0,1,0)$. Every nonzero Smith invariant
of each boundary matrix is $1$. The program checks the
relative deformation functor and natural transformation,
and repeats the calculation for $n=3$. The proof above
establishes all $n\geq2$ without extrapolating those tests.
\end{proof}

\begin{theorem}[Proved: the maximal class allowed by C048]
\label{F:interval-impossibility}
Fix $n\geq2$ and the endpoint inclusion $j$ above. Let $W$
be a fundamental localizer for which $S^2$ is not
$W$-aspherical. If a class $\mathcal K$ of full sieve
inclusions consists of $W$-cofibrations and is closed under
$i\mapsto i\mathbin{\square}j$, then every member of
$\mathcal K$ is an isomorphism. In particular there is no
nontrivial class satisfying C048 for $W_\infty$, for
$W_R$ with a nonzero coefficient ring $R$, or for $W_m$
with $m\geq2$. The class of isomorphisms satisfies the
conditions literally and is the maximal possible class.
\end{theorem}

\begin{proof}
Suppose first that $i:U\hookrightarrow X$ has a crossing
arrow $u\to x$, with $u\in U$ and $x\notin U$. The
classifying functor $\chi:X\to[1]$, equal to $0$ on $U$
and to $1$ on its complement, together with $U\to e$,
is a morphism of arrows from $i$ to $i_0:e\to[1]$.
Choosing $u$ and the indicated crossing arrow gives a
section from $i_0$ to $i$. Their composite is the identity
of $i_0$ in the arrow category. Thus $i_0$ is a retract
of $i$. The pushout-product functor $(-)\mathbin{\square}j$
preserves retract diagrams. If $i\in\mathcal K$, the
retract $i_0\mathbin{\square}j$ of
$i\mathbin{\square}j$ would therefore be a
$W$-cofibration, by Proposition~\ref{F:couniversal}.
This contradicts Theorem~\ref{F:interval-failure}.

If $i$ has no crossing arrows, its sieve property implies
that it is a coproduct-summand inclusion
$U\hookrightarrow U\amalg V$. If it is not an
isomorphism, choose $v\in V$. Its pushout-product with
$j$ is
\[
 U\times I\amalg V\times\{0,2n\}
       \hookrightarrow U\times I\amalg V\times I.
\]
This sieve has the crossing arrow $(v,0)\to(v,1)$.
By closure it belongs to $\mathcal K$, contradicting
the first paragraph. Every member is therefore an
isomorphism. Pushouts and pushout-products of
isomorphisms are isomorphisms, proving the final assertion.
Notice that cobase-change closure was not needed for the
obstruction; adding that requirement cannot repair it.
\end{proof}

This strengthens the calculation at XII ms.~114. The
slice with two incomparable maxima does show failure
of the right-adjoint condition, but the obstruction
already concerns preservation of weak equivalences by
the entire inclusion. There is a related obstruction
to cartesian model structures on $\Cat$ in
\citet[Propositions~4.1--4.3]{F_Pavlov2024}; the theorem
above concerns the precise alternating endpoint
operation in this manuscript. A replacement by a
homotopy pushout, or by a cofibrant replacement of the
ordinary pushout-product map, retains the desired
homotopy type, but changes the specified operation.

\begin{proposition}[Proved: the degree-zero and degree-one exceptions]
\label{F:W0}
For $W_0$, every functor is a $W_0$-cofibration.
For $W_1$, every object-injective functor is a
$W_1$-cofibration. Consequently all full sieve
inclusions form a class satisfying C048 for each
of $W_0$ and $W_1$.
\end{proposition}

\begin{proof}
The functor $\pi_0:\Cat\to\Set$ is left adjoint to
the discrete-category functor and preserves pushouts.
Pushout in sets preserves isomorphisms. Thus pushout
along any functor preserves functors inducing
bijections on components. Full sieves are stable under
cobase change by the collage formula and under the
displayed endpoint pushout-product by inspection of
sources of arrows.

For $W_1$, let $\Pi:\Cat\to\mathbf{Gpd}$ be
groupoid completion, retaining the object set.
It is left adjoint to inclusion and preserves
pushouts. A functor belongs to $W_1$ precisely
when its groupoid completion is an equivalence:
the groupoid $\Pi A$ is the fundamental groupoid
of $BA$ on its vertices. The folk model structure
on categories has equivalences as weak equivalences,
object-injective functors as cofibrations, and is
left proper; these precise properties are recalled
and used in the proof of
\citet[Theorem~1.6]{F_HORR2024}. A pushout of
groupoids in categories is a groupoid, since
its arrows are generated by invertible arrows.
Thus equivalences of groupoids remain equivalences
after cobase change along object-injective
functors. Apply this to the groupoid completions
of the two successive pushouts in the definition
of a $W_1$-cofibration. Object injectivity survives
the first pushout, because object sets of
categorical pushouts are the pushouts of object
sets. It follows that the second comparison
is an equivalence after $\Pi$, as required.
\end{proof}

\begin{theorem}[Proved: equivalent forms of the remaining presheaf
question]\label{F:remaining-equivalence}
The following five assertions, each quantified over all small
categories and all indicated presheaves, are equivalent.
\begin{enumerate}
\item Every $W_\infty$-dwy\'erienne sieve is discretely flat.
\item $W_\infty$-dwy\'erienne sieves are closed under composition.
\item Every totally $W_{\mathbb Z}$-aspherical presheaf is totally
$W_\infty$-aspherical.
\item If $H$ is a totally $W_\infty$-aspherical presheaf on $U$
and $V\subset U$ is a full cosieve with $H|_V$ nonempty, then
$H|_V$ is totally $W_\infty$-aspherical.
\item Every couniversal $W_\infty$-equivalence which is a
full sieve whose complementary full category is isomorphic
to $e$ (one object and only its identity) is
$W_\infty$-dwy\'erienne.
\end{enumerate}
If assertion~2 fails, a counterexample can be chosen with a
classical Dwyer sieve as its first factor and a one-object
extension as its second factor. The equivalence theorem does
not decide whether these five assertions hold.
\end{theorem}

\begin{proof}
We first establish a construction which turns assertion~3
into a precise test. Let $H:V^{\op}\to\Set$ be totally
$W_{\mathbb Z}$-aspherical. Write $E=V/H$; this category is
nonempty, connected, and integrally acyclic. Adjoin an initial
object $0$ to $V$, obtaining $B=V^{\triangleleft}$, and define
$F:B^{\op}\to\Set$ by
\[
 F(0)=\{*\},\qquad F|_V=H.
\]
The restriction maps to $F(0)$ are the unique maps. Its category
of elements is the initial cone $E^{\triangleleft}$, hence is
contractible. We claim that $F$ is totally
$W_\infty$-aspherical.

For a covariant $L:B\to\Set$, put
$D=\operatorname{El}_E((L|_V)\pi)$. The components of $D$ are
integrally acyclic by Proposition~\ref{E:restricted-actions}.
The category
$\operatorname{El}_{B/F}(L\pi_F)$ is obtained from $D$ by
adjoining one source object for each $a\in L(0)$. The arrows
out of that source have crossing element category exactly $E$:
its functor into $D$ sends
\[
 (v,h)\longmapsto (v,h,L(0\to v)(a)).
\]
The uniqueness of $0\to v$ makes this a functor and identifies
its entire crossing category with $E$. Thus the nerve is formed
by attaching one initial cone on $NE$ for each $a$, along these
maps into $ND$. This can either be read degree by degree from
chains starting at a new source or expressed as the corresponding
homotopy pushout. The relative homology of this attachment is
\[
 H_k\bigl(B\operatorname{El}_{B/F}(L\pi_F),BD;\mathbb Z\bigr)
   \cong
 \bigoplus_{a\in L(0)}\widetilde H_{k-1}(BE;\mathbb Z)=0.
\]
Each cone attaches along a nonempty connected category, so it
neither joins distinct components of $D$ nor creates a new one.
If $D$ is empty, nonemptiness of $E$ forces $L(0)$ to be empty
as well. It follows in every case that the resulting components
are acyclic. Proposition~\ref{E:restricted-actions}, together
with contractibility of $B/F$, proves the claim.

Assume assertion~1. Apply it to the sieve obtained by adjoining
one object to $B$ with crossing presheaf $F$. By
Propositions~\ref{F:weight-test} and~\ref{F:DF-actions}, every
restricted action of $B/F$ has contractible components. Take
the covariant indicator of the cosieve $V\subset B$, empty
at $0$ and a singleton on $V$. Its action category is $E$.
Thus $E$ is contractible. The original total
$W_{\mathbb Z}$-asphericity of $H$ supplies all the acyclicity
conditions in Proposition~\ref{E:restricted-actions}; now
$H$ is totally $W_\infty$-aspherical. This proves $1\Rightarrow3$.

For $3\Rightarrow1$, let $H$ be totally $W_\infty$-aspherical
on $U$, and fix $L:U\to\Set$. Let $V=\operatorname{El}_U(L)$
with projection $\rho:V\to U$, and put $H'=\rho^*H$.
There is a canonical isomorphism
\[
 V/H'\cong\operatorname{El}_{U/H}(L\pi)=D.
\]
For any component $D_c$ of $D$, define a subpresheaf
$H'_c\subset H'$ by keeping just those elements whose objects
of $V/H'$ belong to $D_c$. This is a subpresheaf: a restriction
map gives an arrow in the element category, so cannot cross
between components. Its category of elements is exactly $D_c$,
which is nonempty, connected, and acyclic.

For any covariant $M:V\to\Set$, form the covariant $U$-action
\[
 L'(u)=\coprod_{l\in L(u)}M(u,l),
\]
with transitions induced by $L$ and $M$. Then
$\operatorname{El}_{D_c}(M\pi_c)$ is a union of connected
components of $\operatorname{El}_{U/H}(L'\pi)$: both have
objects $(u,h,l,m)$, and restricting to $D_c$ cannot remove
part of a component. The latter components are acyclic by
totality of $H$. Proposition~\ref{E:restricted-actions}
therefore makes $H'_c$ totally $W_{\mathbb Z}$-aspherical.
Assertion~3 makes it totally $W_\infty$-aspherical, so the
identity test makes $D_c$ contractible. This applies to every
$L$ and every component. Proposition~\ref{F:DF-actions}
proves that $H\otimes_U-$ is homotopical; applying this to
every boundary weight proves assertion~1. Empty weights
automatically give the constant empty coend and cause no
exception.

Assertion~3 implies assertion~4 by
Lemma~\ref{F:cosieve-homology}, with $R=\mathbb Z$, followed
by assertion~3 on the restricted presheaf. Assertion~4 implies
assertion~2 by the same restriction and strict-functor
composition argument as in
Theorem~\ref{F:strict-sieve-homology}.

Finally we prove $2\Rightarrow4$ and the stated special form
of a possible counterexample. Given $V\subset U$ and $H$ as
in assertion~4, form $X$ from the old category $U$ and a
disjoint lower copy $V_0$ of $V$ by setting
\[
 X(v_0,u)=U(v,u),\qquad X(u,v_0)=\varnothing
       \quad(v\in V,\ u\in U),
\]
with composition induced by $U$. Since $V$ is a cosieve,
the cosieve generated by $V_0$ in $X$ is the two copies of
$V$ with these crossing arrows, namely $V\times[1]$.
The inclusion $a:V_0\hookrightarrow X$ is classical Dwyer:
the inclusion into this generated cosieve has right adjoint
which forgets the copy. Let $j:U\hookrightarrow X$ be the
old inclusion. The presheaf $G=j_!H$ is total by stability
of totality under Kan extension, and the co-Yoneda formula
gives
\[
 G|_U=H,\qquad G|_{V_0}=H|_V.
\]
Adjoin one object $t$ to $X$, with crossing presheaf $G$,
to obtain a $W_\infty$-dwy\'erienne sieve $b:X\hookrightarrow Y$.
The composite $ba$ has weight $H|_V$ at $t$. Its other
nonempty boundary weights are representables on $V$;
the weights at old objects outside $V$ are empty by the
cosieve condition. Thus $ba$ is $W_\infty$-dwy\'erienne
exactly when $H|_V$ is total. This proves $2\Rightarrow4$.
Applying assertion~4 to $F$ on $V^{\triangleleft}$ constructed
above proves $4\Rightarrow3$.
Proposition~\ref{F:onepoint-couniversal} identifies assertion~5
with assertion~3: a sieve with complement isomorphic to $e$
is exactly the one-point collage of its boundary presheaf. Couniversality
forces that presheaf to be nonempty, and under this condition
the sieve is $W_\infty$-dwy\'erienne precisely when its
presheaf is total $W_\infty$. Thus all five assertions are equivalent, with
the promised restricted form of a failure of composition.
\end{proof}

\begin{proposition}[Open: the remaining class comparisons and
localizers]\label{F:open}
The arguments in this part do not decide
whether every $W_\infty$-dwy\'erienne sieve is
discretely flat, nor whether
$W_\infty$-dwy\'erienne sieves are always closed
under composition. These are equivalent to the
one-point couniversal-recognition question and the
presheaf implication in
Theorem~\ref{F:remaining-equivalence}. Strict and
total functors compose for every fundamental
localizer. For sieves, the universal assertion
of composition is false by
Theorem~\ref{F:low-sieve-composition}; total
composition for a fixed localizer is characterized
exactly by Theorem~\ref{F:total-sieve-cosieves}.
For finite $n\geq2$, strict $W_n$-dwy\'erienne
sieve composition, and total $W_n$-dwy\'erienne
sieve composition without the finite-category
hypothesis, are not settled here. The remaining
$S^2$-aspherical case of C048 is settled in
Corollary~\ref{F:C048-dichotomy} below.
\end{proposition}

The exact missing implication is identified by
Propositions~\ref{E:restricted-actions} and~\ref{F:DF-actions}:
for a total presheaf, must every acyclic component of a
restricted action be simply connected? A covering of such
a component need not itself be an action factoring through
$U/H\to U$, so the covering argument for total categories
does not supply that implication. Equivalently, by
Theorem~\ref{F:remaining-equivalence}, the unresolved question
is whether total integral-homology asphericity of a presheaf
forces its category of elements to be simply connected.
Any obstruction must have a nontrivial perfect fundamental
group. If $U/H$ is totally aspherical,
Proposition~\ref{F:weight-test} already supplies discrete
flatness. The two different interval tests
must be distinguished: the explicit weak-map
comparison has a degree-three obstruction, while
the couniversality test has a degree-two
obstruction and settles $W_2$ by
Proposition~\ref{F:couniversal}. Degree one
behaves differently because groupoid completion
preserves the relevant ordinary pushouts.

\subsection{The sphere criterion and the general-localizer answer to C048}

\begin{proposition}[Known: the localizer generated by a sphere]\label{F:sphere-localizer}
Let $W$ be any fundamental localizer, with no accessibility assumption,
and let $n\geq0$. The following conditions are equivalent:
\begin{enumerate}
\item the category $\Delta/\partial\Delta^{n+2}$ is $W$-aspherical;
\item $W_n\subseteq W$.
\end{enumerate}
In particular, $S^2$ is $W$-aspherical if and only if $W_1\subseteq W$.
Here an assertion about the asphericity of a space means the
corresponding assertion for a category having its weak homotopy type.
\end{proposition}
\begin{proof}
This is \citet[Proposition~9.2.17, equivalence (a)--(b), pp.~369--370]{Cisinski2006Prefaisceaux}.
The same source's Corollary~9.2.15 identifies its $n$-equivalences with
isomorphisms on $\pi_0$ and on all based $\pi_i$ for $1\leq i\leq n$,
so its convention agrees with Definition~\ref{D:n-convention}.
The proposition is stated for an arbitrary fundamental localizer.
Its hypotheses do not require that $W$ itself be accessible.
The simplex category $\Delta/\partial\Delta^{n+2}$ has the weak
homotopy type of the sphere $S^{n+1}$. Since
$W_\infty\subseteq W$, replacing this category by another representative
of the same weak homotopy type does not affect its $W$-asphericity.
\end{proof}

\begin{lemma}[Proved: the groupoid comparison for an ordinary pushout]\label{F:span-groupoid}
Let $i:A\to B$ be injective on objects, and let $j:A\to C$ be any functor.
Let $\mathcal G$ be the Grothendieck construction of the covariant
diagram $B\leftarrow A\to C$ indexed by the three-object outgoing
span. The canonical comparison
\[
 q:\mathcal G\longrightarrow P:=B\amalg_A C
\]
belongs to $W_1$.
\end{lemma}
\begin{proof}
Write $\Pi E$ for the groupoid completion of a category $E$.
Its components and automorphism groups are respectively the components
and fundamental groups of $BE$, by the edge-path presentation of the
nerve. It therefore suffices to prove that
$\Pi q:Q:=\Pi\mathcal G\to\Pi P$ is an equivalence of groupoids.

The category $\mathcal G$ has tagged copies $a_A,b_B,c_C$ of the objects
of $A,B,C$, and canonical arrows
\[
 u_a:a_A\longrightarrow i(a)_B,\qquad
 v_a:a_A\longrightarrow j(a)_C.
\]
Together with the arrows within the three copies, these generate all
its arrows, with the naturality relations
\[
 u_{a'}\alpha_A=i(\alpha)_B u_a,\qquad
 v_{a'}\alpha_A=j(\alpha)_C v_a
 \quad(\alpha:a\to a').
\]
In $Q$ all these arrows are invertible.

Define a functor $B':B\to Q$ on objects by
\[
 B'(b)=
 \begin{cases}
 j(a)_C,&b=i(a),\\
 b_B,&b\notin i(\operatorname{Ob}A).
 \end{cases}
\]
This is unambiguous because $i$ is injective on objects. Put
$\theta_{i(a)}=v_a u_a^{-1}:i(a)_B\to j(a)_C$ and
$\theta_b=\id_{b_B}$ off the image of $i$, and define
$B'(\beta)=\theta_{b'}\beta_B\theta_b^{-1}$ for $\beta:b\to b'$.
The displayed naturality relations give the strict equality of
functors $B'i=C_{\!C}j$, where $C_{\!C}:C\to Q$ is the tagged-copy
functor. The universal property of the ordinary categorical pushout,
followed by groupoid completion, supplies
$r:\Pi P\to Q$.

The equality $(\Pi q)r=\id_{\Pi P}$ follows on the generating copies
of $B$ and $C$, since every $\theta_b$ maps to an identity.
There is a natural isomorphism
$\id_Q\Rightarrow r(\Pi q)$ whose components are $\theta_b$ on
$b_B$, identities on $c_C$, and $v_a$ on $a_A$.
Naturality on $u_a$ is the equality $\theta_{i(a)}u_a=v_a$;
naturality on $v_a$ is immediate; the remaining naturality equations
are the conjugation formula and the two displayed relations.
This proves the required equivalence and hence $q\in W_1$.
\end{proof}

\begin{theorem}[Proved: all object-injective maps above the one-type localizer]\label{F:sphere-cofibrations}
If a fundamental localizer $W$ makes $S^2$ aspherical, then every
functor injective on objects is a $W$-cofibration in the manuscript's
sense.
\end{theorem}
\begin{proof}
Proposition~\ref{F:sphere-localizer} gives $W_1\subseteq W$.
Fix an object-injective $i:A\to B$, a functor $j:A\to C$, and a
$W$-equivalence $f:C\to C'$ under $A$. The comparison maps from
the two span Grothendieck constructions to
$B\amalg_A C$ and $B\amalg_A C'$ belong to $W_1$ by
Lemma~\ref{F:span-groupoid}, hence belong to $W$.

The functor between the two Grothendieck constructions also belongs
to $W$. For completeness, this uses only the relative Theorem~A
axiom for a fundamental localizer. For a covariant diagram
$D:J\to\Cat$ and $s\in J$, the slice of its Grothendieck construction
over $s$ retracts to $D(s)$ by
\[
 (t,x,\alpha:t\to s)\longmapsto D(\alpha)(x).
\]
Inclusion of the fibre is a section, and the arrows
$(\alpha,\id)$ give a natural transformation from the identity to
the inclusion followed by the retraction. Thus this is a
$W_\infty$-equivalence, naturally in the diagram.
An objectwise $W$-equivalence of diagrams consequently induces a
$W$-equivalence on each such slice, and the relative axiom gives
a $W$-equivalence on Grothendieck constructions.
Apply this to the objectwise map $(\id_B,\id_A,f)$ of spans.

The square of canonical comparisons commutes. Two-out-of-three
therefore gives
$B\amalg_A C\to B\amalg_A C'\in W$, exactly the required
preservation of $W$-equivalences by the pushout functor.
\end{proof}

\begin{corollary}[Proved: the complete C048 dichotomy]\label{F:C048-dichotomy}
Fix an alternating interval $\Phi_n$ with $n\geq2$ and its endpoint
sieve $j:\{0,2n\}\hookrightarrow\Phi_n$. For every fundamental
localizer $W$, the maximal class of full sieve inclusions satisfying
C048 is as follows:
\begin{enumerate}
\item If $S^2$ is $W$-aspherical, all full sieve inclusions form such a
class.
\item If $S^2$ is not $W$-aspherical, the only possible members are
isomorphisms, and the class of isomorphisms satisfies the requirements.
\end{enumerate}
No accessibility assumption on $W$ is needed.
\end{corollary}
\begin{proof}
In the first case Theorem~\ref{F:sphere-cofibrations} makes every
full sieve inclusion a $W$-cofibration. Full sieve inclusions are
preserved by arbitrary categorical cobase change. They are also
preserved by the displayed pushout-product operation: if
$U\hookrightarrow X$ is a full sieve, its pushout-product with $j$
is the inclusion of the full sieve
\[
 (U\times\Phi_n)\cup(X\times\{0,2n\})
 \ \hookrightarrow\ X\times\Phi_n.
\]
Thus the class of all full sieve inclusions has both requested
stability properties and is maximal.
The second case is Theorem~\ref{F:interval-impossibility}.
\end{proof}

\section{The fixed-origin path category}
\subsection{The exact path category and its endpoint functor}

\begin{definition}[Known: the path category used in Chapter VII]\label{G:definition}
Let $Y$ be a small category and $y_0$ an object of $Y$. A type $\tau$ is a finite word in the signs $+$ and $-$. Its interval $I_\tau$ has objects $0,\ldots,n$, where $n$ is the length of $\tau$, and is the free category on the oriented line whose $i$th edge has the indicated sign. A morphism of types is an endpoint-preserving functor $\varphi:I_\tau\to I_\rho$ whose object function is nondecreasing in the numerical ordering of the vertices. This numerical ordering is distinct from the ordering defined by the arrows of the interval.

An object of $P(Y,y_0)=\operatorname{Ch}(Y;y_0,-)$ is a pair $(\tau,c)$ with $c:I_\tau\to Y$ and $c(0)=y_0$. A morphism
\[
 (\tau,c)\longrightarrow(\rho,d)
\]
is a pair $(\varphi,\eta)$ consisting of such a type morphism and a natural transformation $\eta:c\Rightarrow d\varphi$ whose component at $0$ is $\mathrm{id}_{y_0}$. The composite of $(\varphi,\eta)$ and $(\psi,\theta)$ is $(\psi\varphi,(\theta\varphi)\eta)$. The endpoint functor
\[
 q:P(Y,y_0)\longrightarrow Y
\]
sends $(\tau,c)$ to $c(n)$ and $(\varphi,\eta)$ to $\eta_n$. The category $\operatorname{Ch}(Y;a,b)$ has fixed source $a$, fixed endpoint $b$, and requires the natural-transformation components at both endpoints to be identities.
\end{definition}

These are the corrected definitions of Chapter VII, ms.~25--31, with the increasing interval morphisms used in \S\S5--7 \citep{GrothendieckDerivateursVII}. The initially printed use of all interval functors is explicitly corrected there. In particular, all the interval collapses and endpoint inclusions below are permitted morphisms in the actual category, rather than replacements in a homotopy category.

\begin{proposition}[Proved: components of fixed-end path categories]\label{G:components}
For any small category $Y$ and objects $a,b$, evaluation of an edge word in the groupoid completion $\Pi(Y)=Y[\operatorname{Mor}(Y)^{-1}]$ induces a bijection
\[
 \pi_0\operatorname{Ch}(Y;a,b)
 \ \xrightarrow{\ \cong\ }\ \operatorname{Hom}_{\Pi(Y)}(a,b).
\]
The bijection is compatible with concatenation. It is the component-level assertion proved in Chapter VII, \S7, ms.~59--78; we give the verification needed here.
\end{proposition}

\begin{proof}
For a path $c$, multiply its direct edges and the formal inverses of its retrograde edges, in their traversal order. If $\eta:c\Rightarrow d\varphi$ is a path morphism with fixed endpoints, the commutative squares expressing naturality telescope after all arrows have been inverted. The endpoint components are identities, so $c$ and $d\varphi$ have the same value. The functor $\varphi$ preserves the endpoints and the numerical ordering; each edge in its source is sent either to an identity or to a string of consecutive equally directed edges. Thus $d\varphi$ and $d$ have the same value as well. This proves that evaluation is constant on components.

Every arrow in the groupoid completion is represented by a finite word of arrows and inverse arrows, so evaluation is surjective. For injectivity, recall the explicit presentation of the groupoid completion: words are identified by insertion or removal of identity arrows, composition or subdivision of adjacent equally directed arrows, and cancellation of an arrow next to its inverse. Each of these relations is realized by a zigzag in the fixed-end path category. Identity insertion is the pullback along an interval collapse. Composition of two direct arrows is the pullback along the increasing map from the one-edge direct interval to the two-edge direct interval that selects the endpoints; the analogous map works for two retrograde arrows. For $u:a\to b$, the two-edge path $a\xrightarrow{u}b\xleftarrow{u}a$ receives a natural transformation from the constant path at $a$ of the same type, with components $(\mathrm{id}_a,u,\mathrm{id}_a)$. The latter constant path maps to the length-zero path by the collapse to the point. For $b\xleftarrow{u}a\xrightarrow{u}b$, the natural transformation has the opposite direction, to the constant path at $b$, with components $(\mathrm{id}_b,u,\mathrm{id}_b)$. Concatenating any of these diagrams with fixed paths on either side preserves its endpoints and defines a diagram in $\operatorname{Ch}(Y;a,b)$. Thus every defining relation of the groupoid completion preserves components, proving injectivity and compatibility with concatenation.
\end{proof}

\begin{proposition}[Proved: ordinary finality of the endpoint functor]\label{G:finality}
If $BY$ is connected and simply connected, then $q:P(Y,y_0)\to Y$ is ordinary final. More precisely, every object of $y\backslash P(Y,y_0)$ is joined by a single arrow to an object of the fixed-end fibre $\operatorname{Ch}(Y;y_0,y)$, and that fibre is nonempty and connected.
\end{proposition}

\begin{proof}
The groupoid completion $\Pi(Y)$ is the fundamental groupoid of $BY$ on the vertices of the nerve. This can be read directly from the CW structure of $NY$: its edges are arrows of $Y$, its two-simplices impose the composition relations, and higher simplices do not affect the fundamental groupoid. Hence the hypothesis and Proposition~\ref{G:components} make $\operatorname{Ch}(Y;y_0,y)$ nonempty and connected for every $y$.

An arbitrary object of $y\backslash P(Y,y_0)$ consists of a path $c:I_\tau\to Y$ ending at $z$, together with an arrow $\alpha:y\to z$. Append this arrow in the retrograde direction, obtaining a path
\[
 c_\alpha:I_{\tau-}\longrightarrow Y,
 \qquad y_0\leadsto z\xleftarrow{\alpha}y.
\]
Let $\varphi:I_{\tau-}\to I_\tau$ be the identity on the original interval and collapse the last edge onto its terminal vertex. It is nondecreasing and endpoint preserving. There is a natural transformation $c_\alpha\Rightarrow c\varphi$ whose original components are identities and whose new final component is $\alpha$. Naturality on the appended retrograde edge is precisely $\alpha=\alpha$. This defines a path morphism $c_\alpha\to c$ with endpoint $\alpha$, and therefore an arrow
\[
 (c_\alpha,\mathrm{id}_y)\longrightarrow(c,\alpha)
 \quad\text{in }y\backslash P(Y,y_0).
\]
The fixed-end fibre maps into this comma category by $d\mapsto(d,\mathrm{id}_y)$. Its nonemptiness and connectedness, and the displayed arrow for every other object, prove the required nonemptiness and connectedness of the comma category. We have not used a comparison of strict fibres with homotopy fibres, or an assertion of perfectness.
\end{proof}

\begin{proposition}[Proved: a contraction for an initial origin]\label{G:contraction}
If $y_0$ is initial in $Y$, then the identity endofunctor of $P(Y,y_0)$ is joined to a constant functor by a two-step zigzag of natural transformations. In particular $BP(Y,y_0)$ is contractible.
\end{proposition}

\begin{proof}
Define $R:P(Y,y_0)\to P(Y,y_0)$ by replacing each path by the constant path at $y_0$ of the same type, and retaining its interval map on morphisms. The unique arrows $y_0\to c(i)$ give a natural transformation $R\Rightarrow\mathrm{Id}_{P(Y,y_0)}$. Their component at the origin is the identity. Initiality makes every square required for naturality commute, including squares involving a change of type. The interval collapse $I_\tau\to I_{\varnothing}$ gives another natural transformation $R\Rightarrow\mathrm{const}_{(\varnothing,y_0)}$. Realizing the nerve of these natural transformations contracts the identity map to a constant map.
\end{proof}

\begin{theorem}[Refuted: row XIIba-C005]\label{G:counterexample}
There is a finite weakly contractible category $Y$ with a specified object $y_0$ for which $P(Y,y_0)$ is not totally $W_\infty$-aspherical. The same example is not totally $W_R$-aspherical for any nonzero commutative coefficient ring $R$, and is not totally $W_n$-aspherical for any integer $n\geq2$.
\end{theorem}

\begin{proof}
Let $Y=\Psi$ have objects $z,a,b$, identities, and the two nonidentity arrows $r:z\to a$ and $s:z\to b$. Take $y_0=z$. Since $z$ is initial, $BY$ is contractible, and Proposition~\ref{G:contraction} also proves that $BP(Y,z)$ is contractible. Construct a finite category $C$ by adjoining an object $x$ and arrows $\alpha,\beta:x\rightrightarrows z$, with the relations
\[
 r\alpha=r\beta=:\gamma:x\longrightarrow a,
 \qquad s\alpha=s\beta=: \delta:x\longrightarrow b.
\]
There are no additional nonidentity arrows. Composition consists of these four displayed composites and composition with identities; hence the specification gives the whole category.

The full inclusion $j:Y\hookrightarrow C$ is ordinary final. For $c\in Y$, $c\backslash Y$ has the initial object $(c,\mathrm{id}_c)$. For $c=x$, the comma category has four objects $(z,\alpha),(z,\beta),(a,\gamma),(b,\delta)$ and precisely four nonidentity arrows, one from each of the first two objects to each of the last two. It is nonempty and connected. Its nerve is a four-edge circle, so it is deliberately not homotopy final.

Let $K$ be the full subcategory of $C$ on $x,z$. Its nerve is the graph with two parallel edges, so $BK\cong S^1$. The full subcategories $C_a$ on $x,z,a$ and $C_b$ on $x,z,b$ are respectively obtained from $K$ by adjoining a terminal object. Their nerves are the two cones on $NK$. Every nondegenerate simplex of $NC$ belongs to one of these cones, and their intersection is $NK$. Consequently
\[
 BC\cong\Sigma BK\cong S^2.
\]
By Proposition~\ref{G:finality}, $q:P(Y,z)\to Y$ is ordinary final. Ordinary final functors compose, so $jq:P(Y,z)\to C$ is ordinary final. In the presheaf variance fixed in this paper, this says
\[
 (jq)_!1_{P(Y,z)}=1_C,
 \qquad C/(jq)_!1_{P(Y,z)}\cong C.
\]
The specified outgoing functor $jq$ therefore violates the defining test of total $W_\infty$-asphericity, since $BC\cong S^2$. It also violates total $W_R$-asphericity because $H_2(BC;R)=R\ne0$, and total $W_n$-asphericity for $n\geq2$ because $\pi_2(BC)=\mathbb Z$.

The finite input categories, all their compositions, the comma-category connectedness, and the integral homology of $NC$ are checked in \texttt{code/ag\_finite\_examples.py}; its complete output is included. The path category itself is countably infinite, and the proof of finality above applies to paths of arbitrary finite length, not a bounded truncation.
\end{proof}

\begin{proposition}[Proved: a direct failure of filteredness]\label{G:notfiltered}
For the same $Y=\Psi$ and $y_0=z$, the category $P(Y,z)$ is not filtered.
\end{proposition}

\begin{proof}
The one-edge paths $r:z\to a$ and $s:z\to b$ are objects of $P(Y,z)$. If they admitted a common successor $d$ in that category, applying the endpoint functor would give arrows $a\to q(d)$ and $b\to q(d)$ in $Y$. No object of $Y$ receives arrows from both $a$ and $b$. This is already a failure of the common-successor axiom, independently of any conjectured characterization of total asphericity.
\end{proof}

The conclusion answers the manuscript's sentence ``J'en doute fort!'' at XII ms.~103 in the negative. It uses the exact path category of VII, not an equivalent contractible replacement. The earlier companion left C005 unresolved because Proposition~10 itself detects only noncontractibility of $Y$. The additional outgoing final functor $Y\to C$ detects the failure when $BY$ is contractible.

\section{Disposition of the manuscript questions}
The table records the exact scope of the statements proved above. A row marked Open is unresolved in this paper, with no claim that the question is open in the literature. The final column identifies the numbered statements and their proofs.

\small
\begin{longtable}{@{}>{\raggedright\arraybackslash}p{0.19\textwidth}>{\raggedright\arraybackslash}p{0.13\textwidth}>{\raggedright\arraybackslash}p{0.62\textwidth}@{}}
\toprule
Manuscript row or lead & Disposition & Exact result \\
\midrule
\endfirsthead
\toprule Manuscript row or lead & Disposition & Exact result \\
\midrule
\endhead
II-C-003, II ms.~36 & Refuted & $\Delta_{\mathrm{inj}}$ is totally $\Winfty$-aspherical and not filtered (Theorem~\ref{C:counterexample}). The converse holds for every weakly terminal category, hence all monoids and finite-object categories (Theorem~\ref{C:weak-terminal}). The action-category characterization is Theorem~\ref{C:classification}. \\
\addlinespace
XIIb-a-B035, ms.~94--97 & Complete localizer dichotomy & The converse from total presheaf to total element category holds for all presheaves exactly when $S^2$ is $W$-aspherical (Corollary~\ref{E:localizers}). Otherwise the same finite split-idempotent presheaf refutes it. The forward implication and subterminal converse hold for every $W$ (Proposition~\ref{E:forward-subterminal}). \\
\addlinespace
XIIba-C005, ms.~103 & Refuted & The outgoing fork with initial origin has a contractible fixed-origin path category with a final image of type $S^2$ (Theorem~\ref{G:counterexample}). The genuine interval maps and endpoint finality are proved in Part G. \\
\addlinespace
XIIba-C035, ms.~111 & Refuted & A three-object couniversal $\Winfty$-equivalence sieve is discretely flat and an $h$-cofibration but not dwy\'erienne (Theorem~\ref{F:C035}). A separate four-object example separates $h$-cofibrations from discrete flatness (Theorem~\ref{F:h-not-DF}). \\
\addlinespace
XIIba-C037, C038, ms.~111 & Refuted; narrower question reduced & The one-point-complement cofibration-recognition clauses, even with all complementary changes, are false (Corollary~\ref{F:onepoint-cofibration}). One-point couniversality is characterized in Proposition~\ref{F:onepoint-couniversal}; its weak-homotopy recognition converse is the same unresolved assertion as Theorem~\ref{F:remaining-equivalence}. \\
\addlinespace
XIIba-C045, total sieves & Exact criterion; classified cases & Composition is equivalent to nonempty cosieve stability of total categories (Theorem~\ref{F:total-sieve-cosieves}). It holds for $\Winfty$, nonzero $W_R$, and finite-category $W_n$ cases with $n\ge2$; it fails for $W_0,W_1$. The entire $S^2$-aspherical branch is classified in Theorem~\ref{D:sphere-sieve-composition}. The unrestricted higher-truncation cases remain unresolved. \\
\addlinespace
XIIba-C045, strict sieves & Partly proved; Open & Nonzero homological sieve classes compose (Theorem~\ref{F:strict-sieve-homology}); $W_0,W_1$ fail (Theorem~\ref{F:low-sieve-composition}). The $S^2$-aspherical branch is classified by Theorem~\ref{D:sphere-sieve-composition}. The $\Winfty$ case is the exact presheaf question in Theorem~\ref{F:remaining-equivalence}. \\
\addlinespace
XIIba-C047, strict functors & Proved for every $W$ & Restriction by a strictly $W$-aspherical functor preserves total presheaves; consequently strict functors compose (Theorem~\ref{F:strict-composition}). \\
\addlinespace
XIIba-C047, total functors & Proved for every $W$ & Totally $W$-aspherical functors compose for every fundamental localizer, with no extra gluing or accessibility axiom (Theorem~\ref{F:total-functor-general}). \\
\addlinespace
Discrete-flatness comparison & Exact criteria; one common Open assertion & The restricted-action test is Proposition~\ref{F:DF-actions}. Discrete flatness implies the $h$-cofibration property strictly (Proposition~\ref{F:DFh} and Theorem~\ref{F:h-not-DF}). Whether all strict $\Winfty$ sieves are discretely flat is equivalent to strict composition and total $W_{\mathbb Z}$-presheaves being total $\Winfty$ (Theorem~\ref{F:remaining-equivalence}). \\
\addlinespace
XIIba-C048 = XIIbb-A001, ms.~114 & Complete dichotomy for every $W$ & If $S^2$ is $W$-aspherical, all full sieve inclusions work; otherwise the maximal possible class is the isomorphisms (Corollary~\ref{F:C048-dichotomy}). The proof includes the actual ordinary endpoint pushout-product and requires no accessibility assumption. \\
\addlinespace
Lead 1 & Proved and strengthened & The collage is the homotopy cofiber of the specified element projection (Proposition~\ref{B:collage}). Totality forces connected covariant element categories to be themselves total and weakly contractible (Theorem~\ref{C:classification}). \\
\addlinespace
Lead 2 & Proved & The integral $H_1=\mathbb Z/g\mathbb Z$ formula holds for any small base, with the subgroup generated over all targets (Proposition~\ref{B:lead2}); directedness consolidates finite witnesses. \\
\addlinespace
Lead 3 & Proved & The right-zero monoid is cofiltered and contractible with a final image of type $S^2$ (Proposition~\ref{B:right-zero}). Proposition~\ref{B:bands} classifies the complete rectangular-band family. \\
\addlinespace
Lead 4, integral exactness & Refuted literally; corrected necessity Open & Exact colimits alone do not impose nonemptiness or connectedness. With those assumptions they imply totality. The converse is exactly Remark~\ref{D:exact-converse-open}; the affine criterion is Theorem~\ref{D:affine-exactness}. Fixed points do not supply the missing implication (Remark~\ref{D:fpp-not-exact}). \\
\addlinespace
Lead 5 & Refuted & Total presheaves need not have filtered or total element categories. An ind-right adjoint means filtered crossing categories, a stronger requirement (Propositions~\ref{E:flat} and \ref{E:ind-adjoint}). \\
\bottomrule
\end{longtable}
\normalsize


\bibliographystyle{apalike}
\bibliography{TOTASPH_references}
\end{document}
